% Generated human-review packet template.  Source records, Lean statements,
% and raw-source-to-Spec screening fields are substituted by
% scripts/review_dashboard_packet.py.
% Do not hand-edit generated packets; annotate the PDF or reconcile notes into
% the paper-local human-review log.
\documentclass[10pt]{article}
\usepackage[margin=0.43in,footskip=0.2in]{geometry}
\usepackage{fontspec}
% Use a Unicode-complete main face as well as a Unicode-complete code face.
% Reviewer reasons and source labels can contain exact Greek symbols outside
% verbatim blocks; silently dropping one would corrupt the review packet.
\setmainfont{DejaVu Serif}
% The broad DejaVu Sans Unicode coverage preserves Lean's mathematical glyphs
% (including U+2216 set difference and superscript identifiers) in verbatim
% review blocks.  Its proportional metrics are preferable to silently missing
% exact semantic text from a narrower monospace font.
\setmonofont{DejaVu Sans}
\usepackage{hyperref}
\usepackage{amssymb}
\usepackage{fvextra}
\usepackage{enumitem}
\setlength{\parindent}{0pt}
\setlength{\parskip}{0.15em}
\linespread{0.92}
\hypersetup{colorlinks=true,linkcolor=blue,urlcolor=blue,breaklinks=true}
\DefineVerbatimEnvironment{ReviewVerbatim}{Verbatim}{breaklines=true,breakanywhere=true,fontsize=\scriptsize}
\newcommand{\reviewerbox}{%
  \par\noindent\fbox{\parbox{0.96\linewidth}{\rule{0pt}{0.38in}}}\par
}
\newcommand{\reviewmatch}[1]{%
  \par\noindent\textbf{Reviewer check:} $\square$\; #1\par
  \noindent\emph{If left unchecked, explain the discrepancy or uncertainty below.}\par
}

\begin{document}
\fontsize{9}{10}\selectfont

\begin{center}
{\LARGE Human review packet: GeEtAl2024AlignmentAxioms}\\[0.35em]
\small Generated from byte-pinned source excerpts, the expanded PaperInterface
specifications, and hash-bound raw-source-to-Spec screening records.
\end{center}

\noindent\textbf{Paper:} GeEtAl2024AlignmentAxioms\\
\textbf{Paper id:} \texttt{GeEtAl2024AlignmentAxioms}\\
\textbf{Source version:} \parbox[t]{0.76\linewidth}{\raggedright NeurIPS 2024 published conference version}\\
\textbf{Source record:} \texttt{audit/paper\_statement\_map.json}\\
\textbf{Rows:} 15\\
\textbf{Generated:} 2026-09-06\\
\textbf{Source URL:} \href{https://arxiv.org/abs/2405.14758}{official source}





\section*{How to use this packet}
For each source claim, compare the exact source excerpts with the expanded
PaperInterface specification. Mark up this PDF or fill the blank reviewer box in the TeX source;
return the annotated file for reconciliation into the paper-local human-review
log.

\section*{Open the interactive dashboard}
The interactive dashboard is an optional alternative to using this PDF; it
presents the same review surface rather than an additional review requirement.
After forking and cloning (or pulling) AppliedModelingLib, run the following from the
repository root. The cache command is needed only the first time in a new clone.
\begin{ReviewVerbatim}
cd <your-repository-checkout>
lake exe cache get
python3 scripts/review_dashboard.py --paper GeEtAl2024AlignmentAxioms --serve
\end{ReviewVerbatim}
Open \url{http://127.0.0.1:8765/} in a browser and keep that terminal running
while reviewing. The dashboard records saved annotations in this paper's local
reviewer-owned review record.

\section*{Contents}
\begin{itemize}[leftmargin=1.2em]
\item \hyperlink{library-section-8af29e4f263f1bc2}{Semantic library prerequisites (21; not paper claims)}
\begin{itemize}[leftmargin=1.2em]
\item \hyperlink{library-prerequisite-e72857a369826fa4}{Applied\allowbreak{}Modeling\allowbreak{}Lib.\allowbreak{}Alignment.\allowbreak{}Axioms.\allowbreak{}C1\allowbreak{}Linear\allowbreak{}Rank\allowbreak{}Aggregation\allowbreak{}Rule}
\item \hyperlink{library-prerequisite-5fa7bdb49b050adf}{Applied\allowbreak{}Modeling\allowbreak{}Lib.\allowbreak{}Alignment.\allowbreak{}Axioms.\allowbreak{}Feasible\allowbreak{}Profile}
\item \hyperlink{library-prerequisite-a0b67c6a143b0b13}{Applied\allowbreak{}Modeling\allowbreak{}Lib.\allowbreak{}Alignment.\allowbreak{}Axioms.\allowbreak{}Feature\allowbreak{}Vector}
\item \hyperlink{library-prerequisite-b76a2acb4dafbb4d}{Applied\allowbreak{}Modeling\allowbreak{}Lib.\allowbreak{}Alignment.\allowbreak{}Axioms.\allowbreak{}Induces\allowbreak{}Ranking}
\item \hyperlink{library-prerequisite-8c9b06e1ee6dc104}{Applied\allowbreak{}Modeling\allowbreak{}Lib.\allowbreak{}Alignment.\allowbreak{}Axioms.\allowbreak{}Is\allowbreak{}Kemeny\allowbreak{}Selector}
\item \hyperlink{library-prerequisite-3c30aae25bef2671}{Applied\allowbreak{}Modeling\allowbreak{}Lib.\allowbreak{}Alignment.\allowbreak{}Axioms.\allowbreak{}Linear\allowbreak{}Feasible\allowbreak{}Ranking}
\item \hyperlink{library-prerequisite-43447e121a1884f9}{Applied\allowbreak{}Modeling\allowbreak{}Lib.\allowbreak{}Alignment.\allowbreak{}Axioms.\allowbreak{}Linear\allowbreak{}Rank\allowbreak{}Aggregation\allowbreak{}Rule}
\item \hyperlink{library-prerequisite-d58203ca3f4a8c05}{Applied\allowbreak{}Modeling\allowbreak{}Lib.\allowbreak{}Alignment.\allowbreak{}Axioms.\allowbreak{}Linear\allowbreak{}Reward\allowbreak{}Parameter}
\item \hyperlink{library-prerequisite-5d6a6008c423d214}{Applied\allowbreak{}Modeling\allowbreak{}Lib.\allowbreak{}Alignment.\allowbreak{}Axioms.\allowbreak{}Majority\allowbreak{}Consistent}
\item \hyperlink{library-prerequisite-a68d66923a3c5526}{Applied\allowbreak{}Modeling\allowbreak{}Lib.\allowbreak{}Alignment.\allowbreak{}Axioms.\allowbreak{}Nondegenerate\allowbreak{}Parameter}
\item \hyperlink{library-prerequisite-80fdf9d21d096dc8}{Applied\allowbreak{}Modeling\allowbreak{}Lib.\allowbreak{}Alignment.\allowbreak{}Axioms.\allowbreak{}Pareto\allowbreak{}Optimal}
\item \hyperlink{library-prerequisite-9bded048a60dfe0a}{Applied\allowbreak{}Modeling\allowbreak{}Lib.\allowbreak{}Alignment.\allowbreak{}Axioms.\allowbreak{}Ranking\allowbreak{}Profile}
\item \hyperlink{library-prerequisite-ad3d3904c31ec7b5}{Applied\allowbreak{}Modeling\allowbreak{}Lib.\allowbreak{}Alignment.\allowbreak{}Axioms.\allowbreak{}Ranking\allowbreak{}Separability}
\item \hyperlink{library-prerequisite-8feb80e31f32c9f1}{Applied\allowbreak{}Modeling\allowbreak{}Lib.\allowbreak{}Alignment.\allowbreak{}Axioms.\allowbreak{}Same\allowbreak{}Pairwise\allowbreak{}Majority\allowbreak{}Relation}
\item \hyperlink{library-prerequisite-191720c2353b55e9}{Applied\allowbreak{}Modeling\allowbreak{}Lib.\allowbreak{}Alignment.\allowbreak{}Axioms.\allowbreak{}Strictly\allowbreak{}Prefers}
\item \hyperlink{library-prerequisite-2a43c1465e4fa011}{Applied\allowbreak{}Modeling\allowbreak{}Lib.\allowbreak{}Alignment.\allowbreak{}Axioms.\allowbreak{}Winner\allowbreak{}Monotonic}
\item \hyperlink{library-prerequisite-82fd59c97a4067c7}{Applied\allowbreak{}Modeling\allowbreak{}Lib.\allowbreak{}Alignment.\allowbreak{}Axioms.\allowbreak{}canonical\allowbreak{}Kemeny\allowbreak{}Rule}
\item \hyperlink{library-prerequisite-8c57c8838b32b597}{Applied\allowbreak{}Modeling\allowbreak{}Lib.\allowbreak{}Alignment.\allowbreak{}Axioms.\allowbreak{}kemeny\allowbreak{}Disagreement}
\item \hyperlink{library-prerequisite-4a32ddda17499d70}{Applied\allowbreak{}Modeling\allowbreak{}Lib.\allowbreak{}Alignment.\allowbreak{}Axioms.\allowbreak{}linear\allowbreak{}Reward}
\item \hyperlink{library-prerequisite-8c0e6c25845ae78a}{Applied\allowbreak{}Modeling\allowbreak{}Lib.\allowbreak{}Social\allowbreak{}Choice.\allowbreak{}Ranking.\allowbreak{}Candidate}
\item \hyperlink{library-prerequisite-5cde54446d1ae021}{Applied\allowbreak{}Modeling\allowbreak{}Lib.\allowbreak{}Social\allowbreak{}Choice.\allowbreak{}Ranking.\allowbreak{}Ranking}
\end{itemize}
\item \hyperlink{paper-prerequisite-section-5ef5ef0364b6939c}{Paper-specific semantic prerequisites (8; not paper claims)}
\begin{itemize}[leftmargin=1.2em]
\item \hyperlink{paper-prerequisite-89cd467f0989921f}{Ge\allowbreak{}Et\allowbreak{}Al2024\allowbreak{}Alignment\allowbreak{}Axioms.\allowbreak{}Is\allowbreak{}Fixed\allowbreak{}Tie\allowbreak{}Leximax\allowbreak{}Plurality\allowbreak{}Selector}
\item \hyperlink{paper-prerequisite-c8ae4ba0fedf8db8}{Ge\allowbreak{}Et\allowbreak{}Al2024\allowbreak{}Alignment\allowbreak{}Axioms.\allowbreak{}Is\allowbreak{}Majority\allowbreak{}Loss\allowbreak{}Minimizing}
\item \hyperlink{paper-prerequisite-fa6ae496ed9c0c39}{Ge\allowbreak{}Et\allowbreak{}Al2024\allowbreak{}Alignment\allowbreak{}Axioms.\allowbreak{}Is\allowbreak{}Pareto\allowbreak{}Kemeny\allowbreak{}Selector}
\item \hyperlink{paper-prerequisite-fc38b98f65c4e72f}{Ge\allowbreak{}Et\allowbreak{}Al2024\allowbreak{}Alignment\allowbreak{}Axioms.\allowbreak{}fixed\allowbreak{}Tie\allowbreak{}LCPO}
\item \hyperlink{paper-prerequisite-e532c12e5479137f}{Ge\allowbreak{}Et\allowbreak{}Al2024\allowbreak{}Alignment\allowbreak{}Axioms.\allowbreak{}global\allowbreak{}Majority\allowbreak{}Loss\allowbreak{}Infimum}
\item \hyperlink{paper-prerequisite-8892d6fc4aee77f3}{Ge\allowbreak{}Et\allowbreak{}Al2024\allowbreak{}Alignment\allowbreak{}Axioms.\allowbreak{}majority\allowbreak{}Loss}
\item \hyperlink{paper-prerequisite-465cdc3ee9c9c6a0}{Ge\allowbreak{}Et\allowbreak{}Al2024\allowbreak{}Alignment\allowbreak{}Axioms.\allowbreak{}restricted\allowbreak{}Majority\allowbreak{}Loss\allowbreak{}Infimum}
\item \hyperlink{paper-prerequisite-efbd7cf1dcd57ea6}{Ge\allowbreak{}Et\allowbreak{}Al2024\allowbreak{}Alignment\allowbreak{}Axioms.\allowbreak{}standard\allowbreak{}Loss\allowbreak{}Formulation\allowbreak{}Spec}
\end{itemize}
\item \hyperlink{source-section-a2cb018f24443c81}{Main-text results (9)}
\begin{itemize}[leftmargin=1.2em]
\item \hyperlink{source-claim-74b7c9c683220bb0}{definition2\_\allowbreak{}2\_\allowbreak{}pairwise\allowbreak{}Majority\allowbreak{}Consistency\allowbreak{}And\allowbreak{}Consequences\allowbreak{}Spec}
\item \hyperlink{source-claim-eb42e4a359411adb}{theorem3\_\allowbreak{}1\_\allowbreak{}loss\allowbreak{}Based\allowbreak{}Impossibility\allowbreak{}Spec}
\item \hyperlink{source-claim-31b67361459405c7}{lemma3\_\allowbreak{}2\_\allowbreak{}unconstrained\allowbreak{}Core\allowbreak{}Separation\allowbreak{}Spec}
\item \hyperlink{source-claim-df13a69840ea5ff9}{lemma3\_\allowbreak{}3\_\allowbreak{}core\allowbreak{}Minimizers\allowbreak{}Spec}
\item \hyperlink{source-claim-1682b709e6333170}{lemma3\_\allowbreak{}4\_\allowbreak{}corrected\allowbreak{}Copy\allowbreak{}Cone\allowbreak{}Spec}
\item \hyperlink{source-claim-03ca4893d06d5c73}{lemma3\_\allowbreak{}5\_\allowbreak{}corrected\allowbreak{}Perturbation\allowbreak{}Gap\allowbreak{}Spec}
\item \hyperlink{source-claim-bbe7cd3338235aa4}{theorem3\_\allowbreak{}6\_\allowbreak{}majority\allowbreak{}Loss\allowbreak{}Spec}
\item \hyperlink{source-claim-d5b15047bfa92842}{theorem3\_\allowbreak{}7\_\allowbreak{}C1\allowbreak{}Fails\allowbreak{}Pareto\allowbreak{}And\allowbreak{}Consequences\allowbreak{}Spec}
\item \hyperlink{source-claim-11bcc2903523a4df}{theorem4\_\allowbreak{}3\_\allowbreak{}constructed\allowbreak{}Fixed\allowbreak{}Tie\allowbreak{}LCPO\allowbreak{}Spec}
\end{itemize}
\item \hyperlink{source-section-1cf4eccafb9ce99a}{Appendix results (6)}
\begin{itemize}[leftmargin=1.2em]
\item \hyperlink{source-claim-42097890307741b0}{appendix\allowbreak{}B\allowbreak{}Explicit\allowbreak{}Infeasible\allowbreak{}PMC\allowbreak{}Spec}
\item \hyperlink{source-claim-19bc4515dc72f018}{theorem\allowbreak{}C\_\allowbreak{}2\_\allowbreak{}copeland\allowbreak{}And\allowbreak{}LCPO\allowbreak{}Fail\allowbreak{}Separability\allowbreak{}Spec}
\item \hyperlink{source-claim-db88b9868ac40a8b}{theorem\allowbreak{}C\_\allowbreak{}3\_\allowbreak{}linear\allowbreak{}Kemeny\allowbreak{}Spec}
\item \hyperlink{source-claim-aba4949a527a8006}{theorem\allowbreak{}C\_\allowbreak{}4\_\allowbreak{}linear\allowbreak{}Kemeny\allowbreak{}Fails\allowbreak{}Pareto\allowbreak{}And\allowbreak{}Majority\allowbreak{}Consistency\allowbreak{}Spec}
\item \hyperlink{source-claim-e6c3a6a88028de60}{theorem\allowbreak{}C\_\allowbreak{}5\_\allowbreak{}pareto\allowbreak{}Kemeny\allowbreak{}Fails\allowbreak{}Separability\allowbreak{}And\allowbreak{}Majority\allowbreak{}Consistency\allowbreak{}Spec}
\item \hyperlink{source-claim-c7e3d91ee78bfd29}{theorem\allowbreak{}C\_\allowbreak{}6\_\allowbreak{}fixed\allowbreak{}Tie\allowbreak{}Leximax\allowbreak{}Plurality\allowbreak{}Spec}
\end{itemize}
\end{itemize}

\clearpage
\hypertarget{library-section-8af29e4f263f1bc2}{}
\section*{Material library prerequisites}
\hypertarget{library-prerequisite-e72857a369826fa4}{}
\subsection*{\small\ttfamily\raggedright Applied\allowbreak{}Modeling\allowbreak{}Lib.\allowbreak{}Alignment.\allowbreak{}Axioms.\allowbreak{}C1\allowbreak{}Linear\allowbreak{}Rank\allowbreak{}Aggregation\allowbreak{}Rule}
\textbf{Source locator:} {\footnotesize\raggedright cited publication, lines 196-\allowbreak{}198\par}
\paragraph{Verbatim paper-source connection}
\begin{ReviewVerbatim}
We pay special attention to a prominent family of rules from social choice theory referred to as C1
rules [14], whose outputs depend only on majority relationships, i.e., they only need to know for
each pair of candidates (a, b) whether the majority prefers a or b.

[Next verbatim source excerpt]

However, despite this good news for PMC, we show that this does not help in achieving PO. In fact,
our negative result extends to every C1 linear rank aggregation rule. Note that if f minimizing ell in
the majority formulation breaks ties consistently (i.e., if multiple feasible rankings are optimal, then
it consistently chooses the same one), then it is C1. We then have the following result, whose proof
is relegated to Appendix A.6.
\end{ReviewVerbatim}



\paragraph{Lean-expanded library semantic target}
\begin{ReviewVerbatim}
type:
{Voter : Type u_1} →
  [Fintype Voter] →
    {n : ℕ} →
      (feasible : AppliedModelingLib.SocialChoice.Ranking.Ranking n → Prop) →
        AppliedModelingLib.Alignment.Axioms.LinearRankAggregationRule Voter n feasible → Prop

value:
fun {Voter : Type u_1} [Fintype Voter] {n : ℕ} (feasible : AppliedModelingLib.SocialChoice.Ranking.Ranking n → Prop)
    (rule : AppliedModelingLib.Alignment.Axioms.LinearRankAggregationRule Voter n feasible) =>
  ∀ (firstProfile secondProfile : AppliedModelingLib.Alignment.Axioms.RankingProfile Voter n),
    AppliedModelingLib.Alignment.Axioms.SamePairwiseMajorityRelation firstProfile secondProfile →
      rule.1 firstProfile = rule.1 secondProfile
\end{ReviewVerbatim}

\paragraph{Exact Lean library declaration}
\begin{ReviewVerbatim}
/--
The paper's `C1` condition: a rule's output depends only on the directed
strict-pairwise-majority relation of its input profile.
-/
def C1LinearRankAggregationRule {Voter : Type*} [Fintype Voter] {n : ℕ}
    (feasible : Ranking n → Prop) (rule : LinearRankAggregationRule Voter n feasible) : Prop :=
  ∀ firstProfile secondProfile,
    SamePairwiseMajorityRelation firstProfile secondProfile →
      rule.run firstProfile = rule.run secondProfile
\end{ReviewVerbatim}

\paragraph{Recorded source-to-library screening}
\textbf{Verdict:} \texttt{matches}
\\
{\raggedright
\textbf{Reason:} Compared the source definition that a C1 rule depends only on pairwise majority relationships and all Theorem 3.\allowbreak{}7/\allowbreak{}consequence uses.\allowbreak{} The target compares every directed strict-\allowbreak{}majority edge in two profiles and requires identical rule outputs exactly when those relations coincide;\allowbreak{} tied pairs are represented by absence of either strict edge, so no majority information is omitted or added.\allowbreak{}
\par}
\reviewmatch{Matches source input}
\noindent\textbf{Reviewer annotation}\par
\reviewerbox
\clearpage
\hypertarget{library-prerequisite-5fa7bdb49b050adf}{}
\subsection*{\small\ttfamily\raggedright Applied\allowbreak{}Modeling\allowbreak{}Lib.\allowbreak{}Alignment.\allowbreak{}Axioms.\allowbreak{}Feasible\allowbreak{}Profile}
\textbf{Source locator:} {\footnotesize\raggedright cited publication, lines 158-\allowbreak{}198\par}
\paragraph{Verbatim paper-source connection}
\begin{ReviewVerbatim}
2     The Linear Social Choice Model
Let C be a set of m distinct prompt/responses, referred to as candidates, and let V = {1, . . . , n} be
a set of n human participants, known as voters. We denote by Rd the d-dimensional real space in
which both candidate feature vectors and chosen parameter vectors lie.
Each candidate c ∈ C is associated with a distinct feature vector xc ∈ Rd . A parameter vector
θ ∈ Rd induces a linear reward function rθ : C 7→ R defined by taking the dot product with feature
vectors rθ (c) = ⟨θ, xc ⟩. We will primarily be interested in how these parameterized functions rank
the candidates by reward. Let Ra≻b = {θ | rθ (a) ≥ rθ (b)} be the region where the reward of a is
at least as large as that of b. Note that Ra≻b and Rb≻a split Rd into two half spaces, separated by
the hyperplane orthogonal to xa − xb . Parameter vectors θ on the hyperplane have rθ (a) = rθ (b),
while rankings in the interior of either half-space strictly rank one over the other.
For a ranking σ over the candidates, we say that θ induces σ, denoted θ ▷ σ, if a ≻σ b implies
rθ (a) ≥ rθ (b). Let Rσ = {θ | θ ▷ σ} be the set of vectors θ thatTinduce it. Note that this can
be written as the intersection of corresponding half spaces Rσ = a,b:a≻σ b Ra≻b . Further, the


                                                    3

[form-feed in source extraction]
collection of {Rσ } essentially form a partition of Rd , covering the space and intersecting only at
their boundaries.
We call a θ non-degenerate if it is fully on one side of each of the separating hyperplanes, i.e.,
rθ (a) ̸= rθ (b) for all a, b ∈ C. Non-degenerate parameter vectors lie in the interior of some Rσ ,
and thus induce exactly one ranking. We call σ feasible if Rσ has a nonempty interior, i.e., is induced
by some nondegenerate θ.3
Each voter i ∈ V submits a ranking over the candidate σi . We assume that the feature space is
rich enough that voter preferences can be captured via non-degenerate parameter vectors. In other
words, we assume that each σi is feasible. We refer to the vector of voter rankings π = (σi )i∈V as
a profile. Further, for two candidates a, b, we write na≻b (π) := |{i ∈ V | a ≻σi b}| for the number
of voters that prefer a to b, and wa≻b (π) = na≻b (π)/n for the proportion of such voters. When the
profile π is clear from context, we may shorten these to na≻b and wa≻b , respectively.
We define a parameter aggregation rule as a function that takes as input a profile π and outputs a
parameter vector θ∗ . Our goal is to design parameter aggregation rules such that rθ∗ satisfies desir-
able properties with respect to the voter preferences. However, as the properties we care about will
only be with respect to how rθ∗ ranks the candidates, it will be more convenient to work with what
we call linear rank aggregation rules that take as input a profile π and output a feasible ranking σ.
There is a natural way to interpret a parameter aggregation rule as a linear rank aggregation rule,
namely, output any feasible ranking induced by θ∗ . The exact properties of the parameter aggrega-
tion rule could in principle be sensitive to the tie-breaking of non-degenerate outputs, however, all
of our results will be robust to such tie-breaking.4
We pay special attention to a prominent family of rules from social choice theory referred to as C1
rules [14], whose outputs depend only on majority relationships, i.e., they only need to know for
each pair of candidates (a, b) whether the majority prefers a or b.
\end{ReviewVerbatim}



\paragraph{Lean-expanded library semantic target}
\begin{ReviewVerbatim}
type:
{Voter : Type u_1} →
  {n : ℕ} →
    (AppliedModelingLib.SocialChoice.Ranking.Ranking n → Prop) →
      AppliedModelingLib.Alignment.Axioms.RankingProfile Voter n → Prop

value:
fun {Voter : Type u_1} {n : ℕ} (feasible : AppliedModelingLib.SocialChoice.Ranking.Ranking n → Prop)
    (profile : AppliedModelingLib.Alignment.Axioms.RankingProfile Voter n) =>
  ∀ (voter : Voter), feasible (profile voter)
\end{ReviewVerbatim}

\paragraph{Exact Lean library declaration}
\begin{ReviewVerbatim}
/-- A profile whose submitted rankings all belong to a supplied feasible set. -/
def FeasibleProfile {Voter : Type*} {n : ℕ}
    (feasible : Ranking n → Prop) (profile : RankingProfile Voter n) : Prop :=
  ∀ voter, feasible (profile voter)
\end{ReviewVerbatim}

\paragraph{Recorded source-to-library screening}
\textbf{Verdict:} \texttt{matches}
\\
{\raggedright
\textbf{Reason:} The model says each submitted voter ranking is feasible;\allowbreak{} this is exactly the universal assertion that every profile entry satisfies the supplied feasibility predicate.\allowbreak{}
\par}
\reviewmatch{Matches source input}
\noindent\textbf{Reviewer annotation}\par
\reviewerbox
\clearpage
\hypertarget{library-prerequisite-a0b67c6a143b0b13}{}
\subsection*{\small\ttfamily\raggedright Applied\allowbreak{}Modeling\allowbreak{}Lib.\allowbreak{}Alignment.\allowbreak{}Axioms.\allowbreak{}Feature\allowbreak{}Vector}
\textbf{Source locator:} {\footnotesize\raggedright cited publication, lines 158-\allowbreak{}198\par}
\paragraph{Verbatim paper-source connection}
\begin{ReviewVerbatim}
2     The Linear Social Choice Model
Let C be a set of m distinct prompt/responses, referred to as candidates, and let V = {1, . . . , n} be
a set of n human participants, known as voters. We denote by Rd the d-dimensional real space in
which both candidate feature vectors and chosen parameter vectors lie.
Each candidate c ∈ C is associated with a distinct feature vector xc ∈ Rd . A parameter vector
θ ∈ Rd induces a linear reward function rθ : C 7→ R defined by taking the dot product with feature
vectors rθ (c) = ⟨θ, xc ⟩. We will primarily be interested in how these parameterized functions rank
the candidates by reward. Let Ra≻b = {θ | rθ (a) ≥ rθ (b)} be the region where the reward of a is
at least as large as that of b. Note that Ra≻b and Rb≻a split Rd into two half spaces, separated by
the hyperplane orthogonal to xa − xb . Parameter vectors θ on the hyperplane have rθ (a) = rθ (b),
while rankings in the interior of either half-space strictly rank one over the other.
For a ranking σ over the candidates, we say that θ induces σ, denoted θ ▷ σ, if a ≻σ b implies
rθ (a) ≥ rθ (b). Let Rσ = {θ | θ ▷ σ} be the set of vectors θ thatTinduce it. Note that this can
be written as the intersection of corresponding half spaces Rσ = a,b:a≻σ b Ra≻b . Further, the


                                                    3

[form-feed in source extraction]
collection of {Rσ } essentially form a partition of Rd , covering the space and intersecting only at
their boundaries.
We call a θ non-degenerate if it is fully on one side of each of the separating hyperplanes, i.e.,
rθ (a) ̸= rθ (b) for all a, b ∈ C. Non-degenerate parameter vectors lie in the interior of some Rσ ,
and thus induce exactly one ranking. We call σ feasible if Rσ has a nonempty interior, i.e., is induced
by some nondegenerate θ.3
Each voter i ∈ V submits a ranking over the candidate σi . We assume that the feature space is
rich enough that voter preferences can be captured via non-degenerate parameter vectors. In other
words, we assume that each σi is feasible. We refer to the vector of voter rankings π = (σi )i∈V as
a profile. Further, for two candidates a, b, we write na≻b (π) := |{i ∈ V | a ≻σi b}| for the number
of voters that prefer a to b, and wa≻b (π) = na≻b (π)/n for the proportion of such voters. When the
profile π is clear from context, we may shorten these to na≻b and wa≻b , respectively.
We define a parameter aggregation rule as a function that takes as input a profile π and outputs a
parameter vector θ∗ . Our goal is to design parameter aggregation rules such that rθ∗ satisfies desir-
able properties with respect to the voter preferences. However, as the properties we care about will
only be with respect to how rθ∗ ranks the candidates, it will be more convenient to work with what
we call linear rank aggregation rules that take as input a profile π and output a feasible ranking σ.
There is a natural way to interpret a parameter aggregation rule as a linear rank aggregation rule,
namely, output any feasible ranking induced by θ∗ . The exact properties of the parameter aggrega-
tion rule could in principle be sensitive to the tie-breaking of non-degenerate outputs, however, all
of our results will be robust to such tie-breaking.4
We pay special attention to a prominent family of rules from social choice theory referred to as C1
rules [14], whose outputs depend only on majority relationships, i.e., they only need to know for
each pair of candidates (a, b) whether the majority prefers a or b.
\end{ReviewVerbatim}



\paragraph{Lean-expanded library semantic target}
\begin{ReviewVerbatim}
type:
ℕ → Type

value:
fun (dimension : ℕ) => Fin dimension → ℝ
\end{ReviewVerbatim}

\paragraph{Exact Lean library declaration}
\begin{ReviewVerbatim}
/-- A finite-dimensional real feature vector. -/
abbrev FeatureVector (dimension : ℕ) := Fin dimension → ℝ
\end{ReviewVerbatim}

\paragraph{Recorded source-to-library screening}
\textbf{Verdict:} \texttt{matches}
\\
{\raggedright
\textbf{Reason:} The source places candidate features in R\textasciicircum{}d;\allowbreak{} Fin dimension to R is its finite-\allowbreak{}coordinate representation.\allowbreak{}
\par}
\reviewmatch{Matches source input}
\noindent\textbf{Reviewer annotation}\par
\reviewerbox
\clearpage
\hypertarget{library-prerequisite-b76a2acb4dafbb4d}{}
\subsection*{\small\ttfamily\raggedright Applied\allowbreak{}Modeling\allowbreak{}Lib.\allowbreak{}Alignment.\allowbreak{}Axioms.\allowbreak{}Induces\allowbreak{}Ranking}
\textbf{Source locator:} {\footnotesize\raggedright cited publication, lines 158-\allowbreak{}198\par}
\paragraph{Verbatim paper-source connection}
\begin{ReviewVerbatim}
2     The Linear Social Choice Model
Let C be a set of m distinct prompt/responses, referred to as candidates, and let V = {1, . . . , n} be
a set of n human participants, known as voters. We denote by Rd the d-dimensional real space in
which both candidate feature vectors and chosen parameter vectors lie.
Each candidate c ∈ C is associated with a distinct feature vector xc ∈ Rd . A parameter vector
θ ∈ Rd induces a linear reward function rθ : C 7→ R defined by taking the dot product with feature
vectors rθ (c) = ⟨θ, xc ⟩. We will primarily be interested in how these parameterized functions rank
the candidates by reward. Let Ra≻b = {θ | rθ (a) ≥ rθ (b)} be the region where the reward of a is
at least as large as that of b. Note that Ra≻b and Rb≻a split Rd into two half spaces, separated by
the hyperplane orthogonal to xa − xb . Parameter vectors θ on the hyperplane have rθ (a) = rθ (b),
while rankings in the interior of either half-space strictly rank one over the other.
For a ranking σ over the candidates, we say that θ induces σ, denoted θ ▷ σ, if a ≻σ b implies
rθ (a) ≥ rθ (b). Let Rσ = {θ | θ ▷ σ} be the set of vectors θ thatTinduce it. Note that this can
be written as the intersection of corresponding half spaces Rσ = a,b:a≻σ b Ra≻b . Further, the


                                                    3

[form-feed in source extraction]
collection of {Rσ } essentially form a partition of Rd , covering the space and intersecting only at
their boundaries.
We call a θ non-degenerate if it is fully on one side of each of the separating hyperplanes, i.e.,
rθ (a) ̸= rθ (b) for all a, b ∈ C. Non-degenerate parameter vectors lie in the interior of some Rσ ,
and thus induce exactly one ranking. We call σ feasible if Rσ has a nonempty interior, i.e., is induced
by some nondegenerate θ.3
Each voter i ∈ V submits a ranking over the candidate σi . We assume that the feature space is
rich enough that voter preferences can be captured via non-degenerate parameter vectors. In other
words, we assume that each σi is feasible. We refer to the vector of voter rankings π = (σi )i∈V as
a profile. Further, for two candidates a, b, we write na≻b (π) := |{i ∈ V | a ≻σi b}| for the number
of voters that prefer a to b, and wa≻b (π) = na≻b (π)/n for the proportion of such voters. When the
profile π is clear from context, we may shorten these to na≻b and wa≻b , respectively.
We define a parameter aggregation rule as a function that takes as input a profile π and outputs a
parameter vector θ∗ . Our goal is to design parameter aggregation rules such that rθ∗ satisfies desir-
able properties with respect to the voter preferences. However, as the properties we care about will
only be with respect to how rθ∗ ranks the candidates, it will be more convenient to work with what
we call linear rank aggregation rules that take as input a profile π and output a feasible ranking σ.
There is a natural way to interpret a parameter aggregation rule as a linear rank aggregation rule,
namely, output any feasible ranking induced by θ∗ . The exact properties of the parameter aggrega-
tion rule could in principle be sensitive to the tie-breaking of non-degenerate outputs, however, all
of our results will be robust to such tie-breaking.4
We pay special attention to a prominent family of rules from social choice theory referred to as C1
rules [14], whose outputs depend only on majority relationships, i.e., they only need to know for
each pair of candidates (a, b) whether the majority prefers a or b.
\end{ReviewVerbatim}



\paragraph{Lean-expanded library semantic target}
\begin{ReviewVerbatim}
type:
{n dimension : ℕ} →
  (AppliedModelingLib.SocialChoice.Ranking.Candidate n → AppliedModelingLib.Alignment.Axioms.FeatureVector dimension) →
    AppliedModelingLib.Alignment.Axioms.LinearRewardParameter dimension →
      AppliedModelingLib.SocialChoice.Ranking.Ranking n → Prop

value:
fun {n dimension : ℕ}
    (features :
      AppliedModelingLib.SocialChoice.Ranking.Candidate n → AppliedModelingLib.Alignment.Axioms.FeatureVector dimension)
    (parameter : AppliedModelingLib.Alignment.Axioms.LinearRewardParameter dimension)
    (ranking : AppliedModelingLib.SocialChoice.Ranking.Ranking n) =>
  ∀ (first second : AppliedModelingLib.SocialChoice.Ranking.Candidate n),
    AppliedModelingLib.Alignment.Axioms.StrictlyPrefers ranking first second →
      AppliedModelingLib.Alignment.Axioms.linearReward parameter features second ≤
        AppliedModelingLib.Alignment.Axioms.linearReward parameter features first
\end{ReviewVerbatim}

\paragraph{Exact Lean library declaration}
\begin{ReviewVerbatim}
/-- A parameter weakly respects every strict comparison in `ranking`. -/
def InducesRanking {n dimension : ℕ}
    (features : Candidate n → FeatureVector dimension)
    (parameter : LinearRewardParameter dimension) (ranking : Ranking n) : Prop :=
  ∀ first second,
    StrictlyPrefers ranking first second →
      linearReward parameter features second ≤ linearReward parameter features first
\end{ReviewVerbatim}

\paragraph{Recorded source-to-library screening}
\textbf{Verdict:} \texttt{matches}
\\
{\raggedright
\textbf{Reason:} The source defines induction by a preferred-\allowbreak{}to-\allowbreak{}less-\allowbreak{}preferred pair having weakly larger reward;\allowbreak{} the Lean inequality reward(second) <= reward(first) states the same condition.\allowbreak{}
\par}
\reviewmatch{Matches source input}
\noindent\textbf{Reviewer annotation}\par
\reviewerbox
\clearpage
\hypertarget{library-prerequisite-8c9b06e1ee6dc104}{}
\subsection*{\small\ttfamily\raggedright Applied\allowbreak{}Modeling\allowbreak{}Lib.\allowbreak{}Alignment.\allowbreak{}Axioms.\allowbreak{}Is\allowbreak{}Kemeny\allowbreak{}Selector}
\textbf{Source locator:} {\footnotesize\raggedright cited publication, lines 1423-\allowbreak{}1449\par}
\paragraph{Verbatim paper-source connection}
\begin{ReviewVerbatim}
C.2   Linear Kemeny Rule

Next, we consider a different rule from social choice theory, the Kemeny rule, which can be trans-
formed to the linear setting while maintaining the separability can be achieved along with PMC.
Given an input profile π, the Kemeny rule returns a ranking σ ∗ that minimizes the total number of
pairwise disagreements with voters rankings, i.e.
                                            X X
                              σ ∗ ∈ arg min                1(b ≻σ a)
                                        σ∈S m
                                                 i∈V (a,b):a≻σi b
        m
where S contains all possible permutations of the m candidates. This expression can be equiva-
lently written as                        X
                          σ ∗ ∈ arg min        na≻b (π) · 1(b ≻σ a).
                                      σ∈S m
                                                a̸=b∈C
Here, we define the linear Kemeny rule, which outputs a parameter vector θ∗ that induces a ranking
that minimizes the total number of pairwise disagreements with voters’ rankings, i.e.,
                                        X
                          θ∗ ∈ arg min       na≻b (π) · 1(rθ (b) > rθ (a)).
                                 θ∈Rd    a̸=b∈C
Note that this rule conforms to the standard loss formulation, where the loss function is binary: it is
0 if two rankings agree with respect to the relative ranking of a pair of candidates and 1 otherwise.
Since binary loss is not convex, it does not fit in the impossibility result of Theorem 3.1.
Note that Kemeny is generally NP-hard to compute; however, even for linear Kemeny, there is
at least an exponential time aglorithm by brute-force computing the score of every ranking, and
determining whether or not it is feasible.
\end{ReviewVerbatim}



\paragraph{Lean-expanded library semantic target}
\begin{ReviewVerbatim}
type:
{Voter : Type u_1} →
  [Fintype Voter] →
    {n : ℕ} →
      (feasible : AppliedModelingLib.SocialChoice.Ranking.Ranking n → Prop) →
        AppliedModelingLib.Alignment.Axioms.LinearRankAggregationRule Voter n feasible → Prop

value:
fun {Voter : Type u_1} [Fintype Voter] {n : ℕ} (feasible : AppliedModelingLib.SocialChoice.Ranking.Ranking n → Prop)
    (rule : AppliedModelingLib.Alignment.Axioms.LinearRankAggregationRule Voter n feasible) =>
  ∀ (profile : AppliedModelingLib.Alignment.Axioms.RankingProfile Voter n),
    feasible (rule.1 profile) ∧
      ∀ (contender : AppliedModelingLib.SocialChoice.Ranking.Ranking n),
        feasible contender →
          AppliedModelingLib.Alignment.Axioms.kemenyDisagreement profile (rule.1 profile) ≤
            AppliedModelingLib.Alignment.Axioms.kemenyDisagreement profile contender
\end{ReviewVerbatim}

\paragraph{Exact Lean library declaration}
\begin{ReviewVerbatim}
/-- A rule selects a Kemeny minimizer on every finite profile. -/
def IsKemenySelector {Voter : Type*} [Fintype Voter] {n : ℕ}
    (feasible : Ranking n → Prop) (rule : LinearRankAggregationRule Voter n feasible) : Prop :=
  ∀ profile, IsKemenyMinimizer feasible profile (rule.run profile)
\end{ReviewVerbatim}

\paragraph{Recorded source-to-library screening}
\textbf{Verdict:} \texttt{matches}
\\
{\raggedright
\textbf{Reason:} Compared the linear-\allowbreak{}Kemeny argmin model and the Theorem C.\allowbreak{}4 use.\allowbreak{} On every finite profile the expanded predicate requires the rule output to be feasible and to have no larger exact Kemeny disagreement count than any feasible contender, which is precisely minimization over linearly realizable rankings;\allowbreak{} it does not assume uniqueness or a hidden tie rule.\allowbreak{}
\par}
\reviewmatch{Matches source input}
\noindent\textbf{Reviewer annotation}\par
\reviewerbox
\clearpage
\hypertarget{library-prerequisite-3c30aae25bef2671}{}
\subsection*{\small\ttfamily\raggedright Applied\allowbreak{}Modeling\allowbreak{}Lib.\allowbreak{}Alignment.\allowbreak{}Axioms.\allowbreak{}Linear\allowbreak{}Feasible\allowbreak{}Ranking}
\textbf{Source locator:} {\footnotesize\raggedright cited publication, lines 158-\allowbreak{}198\par}
\paragraph{Verbatim paper-source connection}
\begin{ReviewVerbatim}
2     The Linear Social Choice Model
Let C be a set of m distinct prompt/responses, referred to as candidates, and let V = {1, . . . , n} be
a set of n human participants, known as voters. We denote by Rd the d-dimensional real space in
which both candidate feature vectors and chosen parameter vectors lie.
Each candidate c ∈ C is associated with a distinct feature vector xc ∈ Rd . A parameter vector
θ ∈ Rd induces a linear reward function rθ : C 7→ R defined by taking the dot product with feature
vectors rθ (c) = ⟨θ, xc ⟩. We will primarily be interested in how these parameterized functions rank
the candidates by reward. Let Ra≻b = {θ | rθ (a) ≥ rθ (b)} be the region where the reward of a is
at least as large as that of b. Note that Ra≻b and Rb≻a split Rd into two half spaces, separated by
the hyperplane orthogonal to xa − xb . Parameter vectors θ on the hyperplane have rθ (a) = rθ (b),
while rankings in the interior of either half-space strictly rank one over the other.
For a ranking σ over the candidates, we say that θ induces σ, denoted θ ▷ σ, if a ≻σ b implies
rθ (a) ≥ rθ (b). Let Rσ = {θ | θ ▷ σ} be the set of vectors θ thatTinduce it. Note that this can
be written as the intersection of corresponding half spaces Rσ = a,b:a≻σ b Ra≻b . Further, the


                                                    3

[form-feed in source extraction]
collection of {Rσ } essentially form a partition of Rd , covering the space and intersecting only at
their boundaries.
We call a θ non-degenerate if it is fully on one side of each of the separating hyperplanes, i.e.,
rθ (a) ̸= rθ (b) for all a, b ∈ C. Non-degenerate parameter vectors lie in the interior of some Rσ ,
and thus induce exactly one ranking. We call σ feasible if Rσ has a nonempty interior, i.e., is induced
by some nondegenerate θ.3
Each voter i ∈ V submits a ranking over the candidate σi . We assume that the feature space is
rich enough that voter preferences can be captured via non-degenerate parameter vectors. In other
words, we assume that each σi is feasible. We refer to the vector of voter rankings π = (σi )i∈V as
a profile. Further, for two candidates a, b, we write na≻b (π) := |{i ∈ V | a ≻σi b}| for the number
of voters that prefer a to b, and wa≻b (π) = na≻b (π)/n for the proportion of such voters. When the
profile π is clear from context, we may shorten these to na≻b and wa≻b , respectively.
We define a parameter aggregation rule as a function that takes as input a profile π and outputs a
parameter vector θ∗ . Our goal is to design parameter aggregation rules such that rθ∗ satisfies desir-
able properties with respect to the voter preferences. However, as the properties we care about will
only be with respect to how rθ∗ ranks the candidates, it will be more convenient to work with what
we call linear rank aggregation rules that take as input a profile π and output a feasible ranking σ.
There is a natural way to interpret a parameter aggregation rule as a linear rank aggregation rule,
namely, output any feasible ranking induced by θ∗ . The exact properties of the parameter aggrega-
tion rule could in principle be sensitive to the tie-breaking of non-degenerate outputs, however, all
of our results will be robust to such tie-breaking.4
We pay special attention to a prominent family of rules from social choice theory referred to as C1
rules [14], whose outputs depend only on majority relationships, i.e., they only need to know for
each pair of candidates (a, b) whether the majority prefers a or b.
\end{ReviewVerbatim}



\paragraph{Lean-expanded library semantic target}
\begin{ReviewVerbatim}
type:
{n dimension : ℕ} →
  (AppliedModelingLib.SocialChoice.Ranking.Candidate n → AppliedModelingLib.Alignment.Axioms.FeatureVector dimension) →
    AppliedModelingLib.SocialChoice.Ranking.Ranking n → Prop

value:
fun {n dimension : ℕ}
    (features :
      AppliedModelingLib.SocialChoice.Ranking.Candidate n → AppliedModelingLib.Alignment.Axioms.FeatureVector dimension)
    (ranking : AppliedModelingLib.SocialChoice.Ranking.Ranking n) =>
  ∃ (parameter : AppliedModelingLib.Alignment.Axioms.LinearRewardParameter dimension),
    AppliedModelingLib.Alignment.Axioms.NondegenerateParameter features parameter ∧
      AppliedModelingLib.Alignment.Axioms.InducesRanking features parameter ranking
\end{ReviewVerbatim}

\paragraph{Exact Lean library declaration}
\begin{ReviewVerbatim}
/-- A ranking is feasible when some nondegenerate linear parameter induces it. -/
def LinearFeasibleRanking {n dimension : ℕ}
    (features : Candidate n → FeatureVector dimension) (ranking : Ranking n) : Prop :=
  ∃ parameter : LinearRewardParameter dimension,
    NondegenerateParameter features parameter ∧ InducesRanking features parameter ranking
\end{ReviewVerbatim}

\paragraph{Recorded source-to-library screening}
\textbf{Verdict:} \texttt{matches}
\\
{\raggedright
\textbf{Reason:} The source calls a ranking feasible exactly when some nondegenerate parameter induces it, which is the existential conjunction in the declaration.\allowbreak{}
\par}
\reviewmatch{Matches source input}
\noindent\textbf{Reviewer annotation}\par
\reviewerbox
\clearpage
\hypertarget{library-prerequisite-43447e121a1884f9}{}
\subsection*{\small\ttfamily\raggedright Applied\allowbreak{}Modeling\allowbreak{}Lib.\allowbreak{}Alignment.\allowbreak{}Axioms.\allowbreak{}Linear\allowbreak{}Rank\allowbreak{}Aggregation\allowbreak{}Rule}
\textbf{Source locator:} {\footnotesize\raggedright cited publication, lines 158-\allowbreak{}198\par}
\paragraph{Verbatim paper-source connection}
\begin{ReviewVerbatim}
2     The Linear Social Choice Model
Let C be a set of m distinct prompt/responses, referred to as candidates, and let V = {1, . . . , n} be
a set of n human participants, known as voters. We denote by Rd the d-dimensional real space in
which both candidate feature vectors and chosen parameter vectors lie.
Each candidate c ∈ C is associated with a distinct feature vector xc ∈ Rd . A parameter vector
θ ∈ Rd induces a linear reward function rθ : C 7→ R defined by taking the dot product with feature
vectors rθ (c) = ⟨θ, xc ⟩. We will primarily be interested in how these parameterized functions rank
the candidates by reward. Let Ra≻b = {θ | rθ (a) ≥ rθ (b)} be the region where the reward of a is
at least as large as that of b. Note that Ra≻b and Rb≻a split Rd into two half spaces, separated by
the hyperplane orthogonal to xa − xb . Parameter vectors θ on the hyperplane have rθ (a) = rθ (b),
while rankings in the interior of either half-space strictly rank one over the other.
For a ranking σ over the candidates, we say that θ induces σ, denoted θ ▷ σ, if a ≻σ b implies
rθ (a) ≥ rθ (b). Let Rσ = {θ | θ ▷ σ} be the set of vectors θ thatTinduce it. Note that this can
be written as the intersection of corresponding half spaces Rσ = a,b:a≻σ b Ra≻b . Further, the


                                                    3

[form-feed in source extraction]
collection of {Rσ } essentially form a partition of Rd , covering the space and intersecting only at
their boundaries.
We call a θ non-degenerate if it is fully on one side of each of the separating hyperplanes, i.e.,
rθ (a) ̸= rθ (b) for all a, b ∈ C. Non-degenerate parameter vectors lie in the interior of some Rσ ,
and thus induce exactly one ranking. We call σ feasible if Rσ has a nonempty interior, i.e., is induced
by some nondegenerate θ.3
Each voter i ∈ V submits a ranking over the candidate σi . We assume that the feature space is
rich enough that voter preferences can be captured via non-degenerate parameter vectors. In other
words, we assume that each σi is feasible. We refer to the vector of voter rankings π = (σi )i∈V as
a profile. Further, for two candidates a, b, we write na≻b (π) := |{i ∈ V | a ≻σi b}| for the number
of voters that prefer a to b, and wa≻b (π) = na≻b (π)/n for the proportion of such voters. When the
profile π is clear from context, we may shorten these to na≻b and wa≻b , respectively.
We define a parameter aggregation rule as a function that takes as input a profile π and outputs a
parameter vector θ∗ . Our goal is to design parameter aggregation rules such that rθ∗ satisfies desir-
able properties with respect to the voter preferences. However, as the properties we care about will
only be with respect to how rθ∗ ranks the candidates, it will be more convenient to work with what
we call linear rank aggregation rules that take as input a profile π and output a feasible ranking σ.
There is a natural way to interpret a parameter aggregation rule as a linear rank aggregation rule,
namely, output any feasible ranking induced by θ∗ . The exact properties of the parameter aggrega-
tion rule could in principle be sensitive to the tie-breaking of non-degenerate outputs, however, all
of our results will be robust to such tie-breaking.4
We pay special attention to a prominent family of rules from social choice theory referred to as C1
rules [14], whose outputs depend only on majority relationships, i.e., they only need to know for
each pair of candidates (a, b) whether the majority prefers a or b.
\end{ReviewVerbatim}



\paragraph{Lean-expanded library semantic target}
\begin{ReviewVerbatim}
type:
Type u_1 → (n : ℕ) → (AppliedModelingLib.SocialChoice.Ranking.Ranking n → Prop) → Type u_1

constructor AppliedModelingLib.Alignment.Axioms.LinearRankAggregationRule.mk:
{Voter : Type u_1} →
  {n : ℕ} →
    {feasible : AppliedModelingLib.SocialChoice.Ranking.Ranking n → Prop} →
      (run :
          AppliedModelingLib.Alignment.Axioms.RankingProfile Voter n →
            AppliedModelingLib.SocialChoice.Ranking.Ranking n) →
        (∀ (profile : AppliedModelingLib.Alignment.Axioms.RankingProfile Voter n), feasible (run profile)) →
          AppliedModelingLib.Alignment.Axioms.LinearRankAggregationRule Voter n feasible
\end{ReviewVerbatim}

\paragraph{Exact Lean library declaration}
\begin{ReviewVerbatim}
/--
A linear-rank-aggregation rule with an explicit output-feasibility guarantee.
The feature vectors/parameters which justify `feasible` stay external to this
generic social-choice interface.
-/
structure LinearRankAggregationRule (Voter : Type*) (n : ℕ)
    (feasible : Ranking n → Prop) where
  run : RankingProfile Voter n → Ranking n
  output_feasible : ∀ profile, feasible (run profile)
\end{ReviewVerbatim}

\paragraph{Recorded source-to-library screening}
\textbf{Verdict:} \texttt{matches}
\\
{\raggedright
\textbf{Reason:} The source's linear rank aggregation rule maps a profile to a feasible ranking;\allowbreak{} run plus output\_\allowbreak{}feasible gives precisely that total-\allowbreak{}rule requirement.\allowbreak{}
\par}
\reviewmatch{Matches source input}
\noindent\textbf{Reviewer annotation}\par
\reviewerbox
\clearpage
\hypertarget{library-prerequisite-d58203ca3f4a8c05}{}
\subsection*{\small\ttfamily\raggedright Applied\allowbreak{}Modeling\allowbreak{}Lib.\allowbreak{}Alignment.\allowbreak{}Axioms.\allowbreak{}Linear\allowbreak{}Reward\allowbreak{}Parameter}
\textbf{Source locator:} {\footnotesize\raggedright cited publication, lines 158-\allowbreak{}198\par}
\paragraph{Verbatim paper-source connection}
\begin{ReviewVerbatim}
2     The Linear Social Choice Model
Let C be a set of m distinct prompt/responses, referred to as candidates, and let V = {1, . . . , n} be
a set of n human participants, known as voters. We denote by Rd the d-dimensional real space in
which both candidate feature vectors and chosen parameter vectors lie.
Each candidate c ∈ C is associated with a distinct feature vector xc ∈ Rd . A parameter vector
θ ∈ Rd induces a linear reward function rθ : C 7→ R defined by taking the dot product with feature
vectors rθ (c) = ⟨θ, xc ⟩. We will primarily be interested in how these parameterized functions rank
the candidates by reward. Let Ra≻b = {θ | rθ (a) ≥ rθ (b)} be the region where the reward of a is
at least as large as that of b. Note that Ra≻b and Rb≻a split Rd into two half spaces, separated by
the hyperplane orthogonal to xa − xb . Parameter vectors θ on the hyperplane have rθ (a) = rθ (b),
while rankings in the interior of either half-space strictly rank one over the other.
For a ranking σ over the candidates, we say that θ induces σ, denoted θ ▷ σ, if a ≻σ b implies
rθ (a) ≥ rθ (b). Let Rσ = {θ | θ ▷ σ} be the set of vectors θ thatTinduce it. Note that this can
be written as the intersection of corresponding half spaces Rσ = a,b:a≻σ b Ra≻b . Further, the


                                                    3

[form-feed in source extraction]
collection of {Rσ } essentially form a partition of Rd , covering the space and intersecting only at
their boundaries.
We call a θ non-degenerate if it is fully on one side of each of the separating hyperplanes, i.e.,
rθ (a) ̸= rθ (b) for all a, b ∈ C. Non-degenerate parameter vectors lie in the interior of some Rσ ,
and thus induce exactly one ranking. We call σ feasible if Rσ has a nonempty interior, i.e., is induced
by some nondegenerate θ.3
Each voter i ∈ V submits a ranking over the candidate σi . We assume that the feature space is
rich enough that voter preferences can be captured via non-degenerate parameter vectors. In other
words, we assume that each σi is feasible. We refer to the vector of voter rankings π = (σi )i∈V as
a profile. Further, for two candidates a, b, we write na≻b (π) := |{i ∈ V | a ≻σi b}| for the number
of voters that prefer a to b, and wa≻b (π) = na≻b (π)/n for the proportion of such voters. When the
profile π is clear from context, we may shorten these to na≻b and wa≻b , respectively.
We define a parameter aggregation rule as a function that takes as input a profile π and outputs a
parameter vector θ∗ . Our goal is to design parameter aggregation rules such that rθ∗ satisfies desir-
able properties with respect to the voter preferences. However, as the properties we care about will
only be with respect to how rθ∗ ranks the candidates, it will be more convenient to work with what
we call linear rank aggregation rules that take as input a profile π and output a feasible ranking σ.
There is a natural way to interpret a parameter aggregation rule as a linear rank aggregation rule,
namely, output any feasible ranking induced by θ∗ . The exact properties of the parameter aggrega-
tion rule could in principle be sensitive to the tie-breaking of non-degenerate outputs, however, all
of our results will be robust to such tie-breaking.4
We pay special attention to a prominent family of rules from social choice theory referred to as C1
rules [14], whose outputs depend only on majority relationships, i.e., they only need to know for
each pair of candidates (a, b) whether the majority prefers a or b.
\end{ReviewVerbatim}



\paragraph{Lean-expanded library semantic target}
\begin{ReviewVerbatim}
type:
ℕ → Type

value:
fun (dimension : ℕ) => AppliedModelingLib.Alignment.Axioms.FeatureVector dimension
\end{ReviewVerbatim}

\paragraph{Exact Lean library declaration}
\begin{ReviewVerbatim}
/-- A parameter vector for a linear reward model. -/
abbrev LinearRewardParameter (dimension : ℕ) := FeatureVector dimension
\end{ReviewVerbatim}

\paragraph{Recorded source-to-library screening}
\textbf{Verdict:} \texttt{matches}
\\
{\raggedright
\textbf{Reason:} The source takes a parameter theta in R\textasciicircum{}d;\allowbreak{} this abbreviation is the same finite-\allowbreak{}dimensional real coordinate space.\allowbreak{}
\par}
\reviewmatch{Matches source input}
\noindent\textbf{Reviewer annotation}\par
\reviewerbox
\clearpage
\hypertarget{library-prerequisite-5d6a6008c423d214}{}
\subsection*{\small\ttfamily\raggedright Applied\allowbreak{}Modeling\allowbreak{}Lib.\allowbreak{}Alignment.\allowbreak{}Axioms.\allowbreak{}Majority\allowbreak{}Consistent}
\textbf{Source locator:} {\footnotesize\raggedright cited publication, lines 441-\allowbreak{}443\par}
\paragraph{Verbatim paper-source connection}
\begin{ReviewVerbatim}
Definition 4.1 (majority consistency). A linear rank aggregation rule satisfies majority consistency
if when a candidate a is ranked first by a majority of voters in the input profile, a is ranked first in
the output ranking.
\end{ReviewVerbatim}



\paragraph{Lean-expanded library semantic target}
\begin{ReviewVerbatim}
type:
{Voter : Type u_1} →
  [Fintype Voter] →
    {n : ℕ} →
      (feasible : AppliedModelingLib.SocialChoice.Ranking.Ranking n → Prop) →
        AppliedModelingLib.Alignment.Axioms.LinearRankAggregationRule Voter n feasible → Prop

value:
fun {Voter : Type u_1} [Fintype Voter] {n : ℕ} (feasible : AppliedModelingLib.SocialChoice.Ranking.Ranking n → Prop)
    (rule : AppliedModelingLib.Alignment.Axioms.LinearRankAggregationRule Voter n feasible) =>
  ∀ (profile : AppliedModelingLib.Alignment.Axioms.RankingProfile Voter n)
    (candidate : AppliedModelingLib.SocialChoice.Ranking.Candidate n),
    AppliedModelingLib.Alignment.Axioms.FeasibleProfile feasible profile →
      Fintype.card Voter < 2 * {voter : Voter | (profile voter) 0 = candidate}.card → (rule.1 profile) 0 = candidate
\end{ReviewVerbatim}

\paragraph{Exact Lean library declaration}
\begin{ReviewVerbatim}
/--
Definition 4.1 of Ge et al. (2024), specialized to finite ranking profiles.
-/
def MajorityConsistent {Voter : Type*} [Fintype Voter] {n : ℕ}
    (feasible : Ranking n → Prop)
    (rule : LinearRankAggregationRule Voter n feasible) : Prop :=
  ∀ profile candidate,
    FeasibleProfile feasible profile → MajorityTopChoice profile candidate →
      firstChoice (rule.run profile) = candidate
\end{ReviewVerbatim}

\paragraph{Recorded source-to-library screening}
\textbf{Verdict:} \texttt{matches}
\\
{\raggedright
\textbf{Reason:} Compared Definition 4.\allowbreak{}1 and its positive and negative theorem uses.\allowbreak{} After requiring a source-\allowbreak{}feasible profile, the target uses the exact strict-\allowbreak{}cardinality condition card(V)<2*topCount and requires the aggregate first position to be that candidate, precisely 'ranked first by a majority' rather than by a mere plurality or weak half.\allowbreak{}
\par}
\reviewmatch{Matches source input}
\noindent\textbf{Reviewer annotation}\par
\reviewerbox
\clearpage
\hypertarget{library-prerequisite-a68d66923a3c5526}{}
\subsection*{\small\ttfamily\raggedright Applied\allowbreak{}Modeling\allowbreak{}Lib.\allowbreak{}Alignment.\allowbreak{}Axioms.\allowbreak{}Nondegenerate\allowbreak{}Parameter}
\textbf{Source locator:} {\footnotesize\raggedright cited publication, lines 158-\allowbreak{}198\par}
\paragraph{Verbatim paper-source connection}
\begin{ReviewVerbatim}
2     The Linear Social Choice Model
Let C be a set of m distinct prompt/responses, referred to as candidates, and let V = {1, . . . , n} be
a set of n human participants, known as voters. We denote by Rd the d-dimensional real space in
which both candidate feature vectors and chosen parameter vectors lie.
Each candidate c ∈ C is associated with a distinct feature vector xc ∈ Rd . A parameter vector
θ ∈ Rd induces a linear reward function rθ : C 7→ R defined by taking the dot product with feature
vectors rθ (c) = ⟨θ, xc ⟩. We will primarily be interested in how these parameterized functions rank
the candidates by reward. Let Ra≻b = {θ | rθ (a) ≥ rθ (b)} be the region where the reward of a is
at least as large as that of b. Note that Ra≻b and Rb≻a split Rd into two half spaces, separated by
the hyperplane orthogonal to xa − xb . Parameter vectors θ on the hyperplane have rθ (a) = rθ (b),
while rankings in the interior of either half-space strictly rank one over the other.
For a ranking σ over the candidates, we say that θ induces σ, denoted θ ▷ σ, if a ≻σ b implies
rθ (a) ≥ rθ (b). Let Rσ = {θ | θ ▷ σ} be the set of vectors θ thatTinduce it. Note that this can
be written as the intersection of corresponding half spaces Rσ = a,b:a≻σ b Ra≻b . Further, the


                                                    3

[form-feed in source extraction]
collection of {Rσ } essentially form a partition of Rd , covering the space and intersecting only at
their boundaries.
We call a θ non-degenerate if it is fully on one side of each of the separating hyperplanes, i.e.,
rθ (a) ̸= rθ (b) for all a, b ∈ C. Non-degenerate parameter vectors lie in the interior of some Rσ ,
and thus induce exactly one ranking. We call σ feasible if Rσ has a nonempty interior, i.e., is induced
by some nondegenerate θ.3
Each voter i ∈ V submits a ranking over the candidate σi . We assume that the feature space is
rich enough that voter preferences can be captured via non-degenerate parameter vectors. In other
words, we assume that each σi is feasible. We refer to the vector of voter rankings π = (σi )i∈V as
a profile. Further, for two candidates a, b, we write na≻b (π) := |{i ∈ V | a ≻σi b}| for the number
of voters that prefer a to b, and wa≻b (π) = na≻b (π)/n for the proportion of such voters. When the
profile π is clear from context, we may shorten these to na≻b and wa≻b , respectively.
We define a parameter aggregation rule as a function that takes as input a profile π and outputs a
parameter vector θ∗ . Our goal is to design parameter aggregation rules such that rθ∗ satisfies desir-
able properties with respect to the voter preferences. However, as the properties we care about will
only be with respect to how rθ∗ ranks the candidates, it will be more convenient to work with what
we call linear rank aggregation rules that take as input a profile π and output a feasible ranking σ.
There is a natural way to interpret a parameter aggregation rule as a linear rank aggregation rule,
namely, output any feasible ranking induced by θ∗ . The exact properties of the parameter aggrega-
tion rule could in principle be sensitive to the tie-breaking of non-degenerate outputs, however, all
of our results will be robust to such tie-breaking.4
We pay special attention to a prominent family of rules from social choice theory referred to as C1
rules [14], whose outputs depend only on majority relationships, i.e., they only need to know for
each pair of candidates (a, b) whether the majority prefers a or b.
\end{ReviewVerbatim}



\paragraph{Lean-expanded library semantic target}
\begin{ReviewVerbatim}
type:
{n dimension : ℕ} →
  (AppliedModelingLib.SocialChoice.Ranking.Candidate n → AppliedModelingLib.Alignment.Axioms.FeatureVector dimension) →
    AppliedModelingLib.Alignment.Axioms.LinearRewardParameter dimension → Prop

value:
fun {n dimension : ℕ}
    (features :
      AppliedModelingLib.SocialChoice.Ranking.Candidate n → AppliedModelingLib.Alignment.Axioms.FeatureVector dimension)
    (parameter : AppliedModelingLib.Alignment.Axioms.LinearRewardParameter dimension) =>
  ∀ (first second : AppliedModelingLib.SocialChoice.Ranking.Candidate n),
    first ≠ second →
      AppliedModelingLib.Alignment.Axioms.linearReward parameter features first ≠
        AppliedModelingLib.Alignment.Axioms.linearReward parameter features second
\end{ReviewVerbatim}

\paragraph{Exact Lean library declaration}
\begin{ReviewVerbatim}
/-- A parameter gives distinct rewards to every pair of distinct candidates. -/
def NondegenerateParameter {n dimension : ℕ}
    (features : Candidate n → FeatureVector dimension)
    (parameter : LinearRewardParameter dimension) : Prop :=
  ∀ first second, first ≠ second →
    linearReward parameter features first ≠ linearReward parameter features second
\end{ReviewVerbatim}

\paragraph{Recorded source-to-library screening}
\textbf{Verdict:} \texttt{matches}
\\
{\raggedright
\textbf{Reason:} The source's nondegeneracy condition requires unequal rewards across candidate-\allowbreak{}separating comparisons;\allowbreak{} Lean states this for every distinct candidate pair.\allowbreak{}
\par}
\reviewmatch{Matches source input}
\noindent\textbf{Reviewer annotation}\par
\reviewerbox
\clearpage
\hypertarget{library-prerequisite-80fdf9d21d096dc8}{}
\subsection*{\small\ttfamily\raggedright Applied\allowbreak{}Modeling\allowbreak{}Lib.\allowbreak{}Alignment.\allowbreak{}Axioms.\allowbreak{}Pareto\allowbreak{}Optimal}
\textbf{Source locator:} {\footnotesize\raggedright cited publication, lines 201-\allowbreak{}203\par}
\paragraph{Verbatim paper-source connection}
\begin{ReviewVerbatim}
Definition 2.1 (Pareto Optimality). A linear rank aggregation rule f satisfies Pareto optimality if,
whenever every voter prefers candidate a over candidate b on π, i.e., wa≻b (π) = 1, then candidate
a is ranked higher than candidate b in the output ranking, i.e., a ≻f (π) b.
\end{ReviewVerbatim}



\paragraph{Lean-expanded library semantic target}
\begin{ReviewVerbatim}
type:
{Voter : Type u_1} →
  {n : ℕ} →
    (feasible : AppliedModelingLib.SocialChoice.Ranking.Ranking n → Prop) →
      AppliedModelingLib.Alignment.Axioms.LinearRankAggregationRule Voter n feasible → Prop

value:
fun {Voter : Type u_1} {n : ℕ} (feasible : AppliedModelingLib.SocialChoice.Ranking.Ranking n → Prop)
    (rule : AppliedModelingLib.Alignment.Axioms.LinearRankAggregationRule Voter n feasible) =>
  ∀ (profile : AppliedModelingLib.Alignment.Axioms.RankingProfile Voter n)
    (first second : AppliedModelingLib.SocialChoice.Ranking.Candidate n),
    AppliedModelingLib.Alignment.Axioms.FeasibleProfile feasible profile →
      (∀ (voter : Voter), AppliedModelingLib.Alignment.Axioms.StrictlyPrefers (profile voter) first second) →
        AppliedModelingLib.Alignment.Axioms.StrictlyPrefers (rule.1 profile) first second
\end{ReviewVerbatim}

\paragraph{Exact Lean library declaration}
\begin{ReviewVerbatim}
/--
Definition 2.1 of Ge et al. (2024), specialized to finite ranking profiles:
unanimously preferred candidates remain ordered that way by the rule.
-/
def ParetoOptimal {Voter : Type*} {n : ℕ} (feasible : Ranking n → Prop)
    (rule : LinearRankAggregationRule Voter n feasible) : Prop :=
  ∀ profile first second,
    FeasibleProfile feasible profile → UniversallyPreferred profile first second →
      StrictlyPrefers (rule.run profile) first second
\end{ReviewVerbatim}

\paragraph{Recorded source-to-library screening}
\textbf{Verdict:} \texttt{matches}
\\
{\raggedright
\textbf{Reason:} Compared Definition 2.\allowbreak{}1 and all PO success/\allowbreak{}failure uses.\allowbreak{} For every source-\allowbreak{}feasible profile and candidate pair, unanimous strict preference by all voters implies the same strict comparison in the output ranking, exactly the wa>b=1 condition and conclusion;\allowbreak{} no majority or score surrogate replaces unanimity.\allowbreak{}
\par}
\reviewmatch{Matches source input}
\noindent\textbf{Reviewer annotation}\par
\reviewerbox
\clearpage
\hypertarget{library-prerequisite-9bded048a60dfe0a}{}
\subsection*{\small\ttfamily\raggedright Applied\allowbreak{}Modeling\allowbreak{}Lib.\allowbreak{}Alignment.\allowbreak{}Axioms.\allowbreak{}Ranking\allowbreak{}Profile}
\textbf{Source locator:} {\footnotesize\raggedright cited publication, lines 158-\allowbreak{}198\par}
\paragraph{Verbatim paper-source connection}
\begin{ReviewVerbatim}
2     The Linear Social Choice Model
Let C be a set of m distinct prompt/responses, referred to as candidates, and let V = {1, . . . , n} be
a set of n human participants, known as voters. We denote by Rd the d-dimensional real space in
which both candidate feature vectors and chosen parameter vectors lie.
Each candidate c ∈ C is associated with a distinct feature vector xc ∈ Rd . A parameter vector
θ ∈ Rd induces a linear reward function rθ : C 7→ R defined by taking the dot product with feature
vectors rθ (c) = ⟨θ, xc ⟩. We will primarily be interested in how these parameterized functions rank
the candidates by reward. Let Ra≻b = {θ | rθ (a) ≥ rθ (b)} be the region where the reward of a is
at least as large as that of b. Note that Ra≻b and Rb≻a split Rd into two half spaces, separated by
the hyperplane orthogonal to xa − xb . Parameter vectors θ on the hyperplane have rθ (a) = rθ (b),
while rankings in the interior of either half-space strictly rank one over the other.
For a ranking σ over the candidates, we say that θ induces σ, denoted θ ▷ σ, if a ≻σ b implies
rθ (a) ≥ rθ (b). Let Rσ = {θ | θ ▷ σ} be the set of vectors θ thatTinduce it. Note that this can
be written as the intersection of corresponding half spaces Rσ = a,b:a≻σ b Ra≻b . Further, the


                                                    3

[form-feed in source extraction]
collection of {Rσ } essentially form a partition of Rd , covering the space and intersecting only at
their boundaries.
We call a θ non-degenerate if it is fully on one side of each of the separating hyperplanes, i.e.,
rθ (a) ̸= rθ (b) for all a, b ∈ C. Non-degenerate parameter vectors lie in the interior of some Rσ ,
and thus induce exactly one ranking. We call σ feasible if Rσ has a nonempty interior, i.e., is induced
by some nondegenerate θ.3
Each voter i ∈ V submits a ranking over the candidate σi . We assume that the feature space is
rich enough that voter preferences can be captured via non-degenerate parameter vectors. In other
words, we assume that each σi is feasible. We refer to the vector of voter rankings π = (σi )i∈V as
a profile. Further, for two candidates a, b, we write na≻b (π) := |{i ∈ V | a ≻σi b}| for the number
of voters that prefer a to b, and wa≻b (π) = na≻b (π)/n for the proportion of such voters. When the
profile π is clear from context, we may shorten these to na≻b and wa≻b , respectively.
We define a parameter aggregation rule as a function that takes as input a profile π and outputs a
parameter vector θ∗ . Our goal is to design parameter aggregation rules such that rθ∗ satisfies desir-
able properties with respect to the voter preferences. However, as the properties we care about will
only be with respect to how rθ∗ ranks the candidates, it will be more convenient to work with what
we call linear rank aggregation rules that take as input a profile π and output a feasible ranking σ.
There is a natural way to interpret a parameter aggregation rule as a linear rank aggregation rule,
namely, output any feasible ranking induced by θ∗ . The exact properties of the parameter aggrega-
tion rule could in principle be sensitive to the tie-breaking of non-degenerate outputs, however, all
of our results will be robust to such tie-breaking.4
We pay special attention to a prominent family of rules from social choice theory referred to as C1
rules [14], whose outputs depend only on majority relationships, i.e., they only need to know for
each pair of candidates (a, b) whether the majority prefers a or b.
\end{ReviewVerbatim}



\paragraph{Lean-expanded library semantic target}
\begin{ReviewVerbatim}
type:
Type u_1 → ℕ → Type u_1

value:
fun (Voter : Type u_1) (n : ℕ) => Voter → AppliedModelingLib.SocialChoice.Ranking.Ranking n
\end{ReviewVerbatim}

\paragraph{Exact Lean library declaration}
\begin{ReviewVerbatim}
/-- A profile assigns one complete ranking to each voter. -/
abbrev RankingProfile (Voter : Type*) (n : ℕ) := Voter → Ranking n
\end{ReviewVerbatim}

\paragraph{Recorded source-to-library screening}
\textbf{Verdict:} \texttt{matches}
\\
{\raggedright
\textbf{Reason:} The source defines the profile as the voter-\allowbreak{}indexed vector pi = (sigma\_\allowbreak{}i);\allowbreak{} Voter to Ranking is that representation.\allowbreak{}
\par}
\reviewmatch{Matches source input}
\noindent\textbf{Reviewer annotation}\par
\reviewerbox
\clearpage
\hypertarget{library-prerequisite-ad3d3904c31ec7b5}{}
\subsection*{\small\ttfamily\raggedright Applied\allowbreak{}Modeling\allowbreak{}Lib.\allowbreak{}Alignment.\allowbreak{}Axioms.\allowbreak{}Ranking\allowbreak{}Separability}
\textbf{Source locator:} {\footnotesize\raggedright cited publication, lines 1389-\allowbreak{}1391\par}
\paragraph{Verbatim paper-source connection}
\begin{ReviewVerbatim}
Definition C.1 (Separability). A ranking aggregation rule satisfies ranking separability (or separa-
bility for short) if, when two profiles yield identical output rankings, when combined into a single
profile, this should also produce the same output ranking.
\end{ReviewVerbatim}



\paragraph{Lean-expanded library semantic target}
\begin{ReviewVerbatim}
type:
{n : ℕ} →
  (feasible : AppliedModelingLib.SocialChoice.Ranking.Ranking n → Prop) →
    ((voterCount : ℕ) → AppliedModelingLib.Alignment.Axioms.LinearRankAggregationRule (Fin voterCount) n feasible) →
      Prop

value:
fun {n : ℕ} (feasible : AppliedModelingLib.SocialChoice.Ranking.Ranking n → Prop)
    (rule :
      (voterCount : ℕ) → AppliedModelingLib.Alignment.Axioms.LinearRankAggregationRule (Fin voterCount) n feasible) =>
  ∀ (leftCount rightCount : ℕ) (left : AppliedModelingLib.Alignment.Axioms.RankingProfile (Fin leftCount) n)
    (right : AppliedModelingLib.Alignment.Axioms.RankingProfile (Fin rightCount) n),
    (rule leftCount).1 left = (rule rightCount).1 right →
      ((rule (leftCount + rightCount)).1 fun (i : Fin (leftCount + rightCount)) => Fin.addCases left right i) =
        (rule leftCount).1 left
\end{ReviewVerbatim}

\paragraph{Exact Lean library declaration}
\begin{ReviewVerbatim}
/--
Definition C.1 of Ge et al. (2024), with a rule family indexed by the finite
number of voters.  The conclusion compares the two original equal outputs to
the output on their literal concatenation.
-/
def RankingSeparability {n : ℕ} (feasible : Ranking n → Prop)
    (rule : ∀ voterCount : ℕ, LinearRankAggregationRule (Fin voterCount) n feasible) : Prop :=
  ∀ leftCount rightCount
    (left : RankingProfile (Fin leftCount) n)
    (right : RankingProfile (Fin rightCount) n),
      (rule leftCount).run left = (rule rightCount).run right →
        (rule (leftCount + rightCount)).run (rankingProfileAppend left right) =
          (rule leftCount).run left
\end{ReviewVerbatim}

\paragraph{Recorded source-to-library screening}
\textbf{Verdict:} \texttt{matches}
\\
{\raggedright
\textbf{Reason:} Compared Definition C.\allowbreak{}1 and all four separability uses.\allowbreak{} The target quantifies two finite profiles, assumes their selected output rankings are identical, forms their literal disjoint concatenation with left and right voter counts, and requires the combined output to equal that common ranking.\allowbreak{} This is exactly the source property, including exact ranking equality rather than score equivalence.\allowbreak{}
\par}
\reviewmatch{Matches source input}
\noindent\textbf{Reviewer annotation}\par
\reviewerbox
\clearpage
\hypertarget{library-prerequisite-8feb80e31f32c9f1}{}
\subsection*{\small\ttfamily\raggedright Applied\allowbreak{}Modeling\allowbreak{}Lib.\allowbreak{}Alignment.\allowbreak{}Axioms.\allowbreak{}Same\allowbreak{}Pairwise\allowbreak{}Majority\allowbreak{}Relation}
\textbf{Source locator:} {\footnotesize\raggedright cited publication, lines 196-\allowbreak{}198\par}
\paragraph{Verbatim paper-source connection}
\begin{ReviewVerbatim}
We pay special attention to a prominent family of rules from social choice theory referred to as C1
rules [14], whose outputs depend only on majority relationships, i.e., they only need to know for
each pair of candidates (a, b) whether the majority prefers a or b.

[Next verbatim source excerpt]

However, despite this good news for PMC, we show that this does not help in achieving PO. In fact,
our negative result extends to every C1 linear rank aggregation rule. Note that if f minimizing ell in
the majority formulation breaks ties consistently (i.e., if multiple feasible rankings are optimal, then
it consistently chooses the same one), then it is C1. We then have the following result, whose proof
is relegated to Appendix A.6.
\end{ReviewVerbatim}



\paragraph{Lean-expanded library semantic target}
\begin{ReviewVerbatim}
type:
{Voter : Type u_1} →
  [Fintype Voter] →
    {n : ℕ} →
      AppliedModelingLib.Alignment.Axioms.RankingProfile Voter n →
        AppliedModelingLib.Alignment.Axioms.RankingProfile Voter n → Prop

value:
fun {Voter : Type u_1} [Fintype Voter] {n : ℕ}
    (firstProfile secondProfile : AppliedModelingLib.Alignment.Axioms.RankingProfile Voter n) =>
  ∀ (first second : AppliedModelingLib.SocialChoice.Ranking.Candidate n),
    Fintype.card Voter <
        2 *
          {voter : Voter | AppliedModelingLib.Alignment.Axioms.StrictlyPrefers (firstProfile voter) first second}.card ↔
      Fintype.card Voter <
        2 *
          {voter : Voter | AppliedModelingLib.Alignment.Axioms.StrictlyPrefers (secondProfile voter) first second}.card
\end{ReviewVerbatim}

\paragraph{Exact Lean library declaration}
\begin{ReviewVerbatim}
/-- Two profiles induce the same directed strict-pairwise-majority relation. -/
def SamePairwiseMajorityRelation {Voter : Type*} [Fintype Voter] {n : ℕ}
    (firstProfile secondProfile : RankingProfile Voter n) : Prop :=
  ∀ first second,
    StrictMajorityPrefers firstProfile first second ↔
      StrictMajorityPrefers secondProfile first second
\end{ReviewVerbatim}

\paragraph{Recorded source-to-library screening}
\textbf{Verdict:} \texttt{matches}
\\
{\raggedright
\textbf{Reason:} Compared the C1 source description and its Theorem 3.\allowbreak{}7 uses.\allowbreak{} The expanded relation checks, for every ordered candidate pair, equivalence of the strict-\allowbreak{}more-\allowbreak{}than-\allowbreak{}half voter condition in the two profiles.\allowbreak{} It therefore identifies exactly the same directed majority graph, including majority ties, and no cardinal or ranking information beyond that graph.\allowbreak{}
\par}
\reviewmatch{Matches source input}
\noindent\textbf{Reviewer annotation}\par
\reviewerbox
\clearpage
\hypertarget{library-prerequisite-191720c2353b55e9}{}
\subsection*{\small\ttfamily\raggedright Applied\allowbreak{}Modeling\allowbreak{}Lib.\allowbreak{}Alignment.\allowbreak{}Axioms.\allowbreak{}Strictly\allowbreak{}Prefers}
\textbf{Source locator:} {\footnotesize\raggedright cited publication, lines 158-\allowbreak{}198\par}
\paragraph{Verbatim paper-source connection}
\begin{ReviewVerbatim}
2     The Linear Social Choice Model
Let C be a set of m distinct prompt/responses, referred to as candidates, and let V = {1, . . . , n} be
a set of n human participants, known as voters. We denote by Rd the d-dimensional real space in
which both candidate feature vectors and chosen parameter vectors lie.
Each candidate c ∈ C is associated with a distinct feature vector xc ∈ Rd . A parameter vector
θ ∈ Rd induces a linear reward function rθ : C 7→ R defined by taking the dot product with feature
vectors rθ (c) = ⟨θ, xc ⟩. We will primarily be interested in how these parameterized functions rank
the candidates by reward. Let Ra≻b = {θ | rθ (a) ≥ rθ (b)} be the region where the reward of a is
at least as large as that of b. Note that Ra≻b and Rb≻a split Rd into two half spaces, separated by
the hyperplane orthogonal to xa − xb . Parameter vectors θ on the hyperplane have rθ (a) = rθ (b),
while rankings in the interior of either half-space strictly rank one over the other.
For a ranking σ over the candidates, we say that θ induces σ, denoted θ ▷ σ, if a ≻σ b implies
rθ (a) ≥ rθ (b). Let Rσ = {θ | θ ▷ σ} be the set of vectors θ thatTinduce it. Note that this can
be written as the intersection of corresponding half spaces Rσ = a,b:a≻σ b Ra≻b . Further, the


                                                    3

[form-feed in source extraction]
collection of {Rσ } essentially form a partition of Rd , covering the space and intersecting only at
their boundaries.
We call a θ non-degenerate if it is fully on one side of each of the separating hyperplanes, i.e.,
rθ (a) ̸= rθ (b) for all a, b ∈ C. Non-degenerate parameter vectors lie in the interior of some Rσ ,
and thus induce exactly one ranking. We call σ feasible if Rσ has a nonempty interior, i.e., is induced
by some nondegenerate θ.3
Each voter i ∈ V submits a ranking over the candidate σi . We assume that the feature space is
rich enough that voter preferences can be captured via non-degenerate parameter vectors. In other
words, we assume that each σi is feasible. We refer to the vector of voter rankings π = (σi )i∈V as
a profile. Further, for two candidates a, b, we write na≻b (π) := |{i ∈ V | a ≻σi b}| for the number
of voters that prefer a to b, and wa≻b (π) = na≻b (π)/n for the proportion of such voters. When the
profile π is clear from context, we may shorten these to na≻b and wa≻b , respectively.
We define a parameter aggregation rule as a function that takes as input a profile π and outputs a
parameter vector θ∗ . Our goal is to design parameter aggregation rules such that rθ∗ satisfies desir-
able properties with respect to the voter preferences. However, as the properties we care about will
only be with respect to how rθ∗ ranks the candidates, it will be more convenient to work with what
we call linear rank aggregation rules that take as input a profile π and output a feasible ranking σ.
There is a natural way to interpret a parameter aggregation rule as a linear rank aggregation rule,
namely, output any feasible ranking induced by θ∗ . The exact properties of the parameter aggrega-
tion rule could in principle be sensitive to the tie-breaking of non-degenerate outputs, however, all
of our results will be robust to such tie-breaking.4
We pay special attention to a prominent family of rules from social choice theory referred to as C1
rules [14], whose outputs depend only on majority relationships, i.e., they only need to know for
each pair of candidates (a, b) whether the majority prefers a or b.
\end{ReviewVerbatim}



\paragraph{Lean-expanded library semantic target}
\begin{ReviewVerbatim}
type:
{n : ℕ} →
  AppliedModelingLib.SocialChoice.Ranking.Ranking n →
    AppliedModelingLib.SocialChoice.Ranking.Candidate n → AppliedModelingLib.SocialChoice.Ranking.Candidate n → Prop

value:
fun {n : ℕ} (ranking : AppliedModelingLib.SocialChoice.Ranking.Ranking n)
    (first second : AppliedModelingLib.SocialChoice.Ranking.Candidate n) =>
  (Equiv.symm ranking) first < (Equiv.symm ranking) second
\end{ReviewVerbatim}

\paragraph{Exact Lean library declaration}
\begin{ReviewVerbatim}
/-- Candidate `first` is strictly above `second` in a ranking. -/
def StrictlyPrefers {n : ℕ} (ranking : Ranking n)
    (first second : Candidate n) : Prop :=
  rankOf ranking first < rankOf ranking second
\end{ReviewVerbatim}

\paragraph{Recorded source-to-library screening}
\textbf{Verdict:} \texttt{matches}
\\
{\raggedright
\textbf{Reason:} The source relation a succeeds\_\allowbreak{}sigma b is strict rank order;\allowbreak{} Lean compares their inverse-\allowbreak{}permutation positions, so smaller position means the same strict preference.\allowbreak{}
\par}
\reviewmatch{Matches source input}
\noindent\textbf{Reviewer annotation}\par
\reviewerbox
\clearpage
\hypertarget{library-prerequisite-2a43c1465e4fa011}{}
\subsection*{\small\ttfamily\raggedright Applied\allowbreak{}Modeling\allowbreak{}Lib.\allowbreak{}Alignment.\allowbreak{}Axioms.\allowbreak{}Winner\allowbreak{}Monotonic}
\textbf{Source locator:} {\footnotesize\raggedright cited publication, lines 452-\allowbreak{}454\par}
\paragraph{Verbatim paper-source connection}
\begin{ReviewVerbatim}
Definition 4.2 (winner monotonicity). A linear rank aggregation rule satisfies winner monotonicity
if, when a candidate a is ranked first in the output ranking, elevating a in any voter’s preference
does not cause a to lose their top position in the updated aggregate ranking.
\end{ReviewVerbatim}



\paragraph{Lean-expanded library semantic target}
\begin{ReviewVerbatim}
type:
{Voter : Type u_1} →
  {n : ℕ} →
    (feasible : AppliedModelingLib.SocialChoice.Ranking.Ranking n → Prop) →
      AppliedModelingLib.Alignment.Axioms.LinearRankAggregationRule Voter n feasible → Prop

value:
fun {Voter : Type u_1} {n : ℕ} (feasible : AppliedModelingLib.SocialChoice.Ranking.Ranking n → Prop)
    (rule : AppliedModelingLib.Alignment.Axioms.LinearRankAggregationRule Voter n feasible) =>
  ∀ (original updated : AppliedModelingLib.Alignment.Axioms.RankingProfile Voter n) (voter : Voter)
    (candidate : AppliedModelingLib.SocialChoice.Ranking.Candidate n),
    AppliedModelingLib.Alignment.Axioms.FeasibleProfile feasible original →
      AppliedModelingLib.Alignment.Axioms.FeasibleProfile feasible updated →
        (rule.1 original) 0 = candidate →
          (((Equiv.symm (updated voter)) candidate ≤ (Equiv.symm (original voter)) candidate ∧
                ∀ (first second : AppliedModelingLib.SocialChoice.Ranking.Candidate n),
                  first ≠ candidate →
                    second ≠ candidate →
                      (AppliedModelingLib.Alignment.Axioms.StrictlyPrefers (original voter) first second ↔
                        AppliedModelingLib.Alignment.Axioms.StrictlyPrefers (updated voter) first second)) ∧
              ∀ (otherVoter : Voter), otherVoter ≠ voter → updated otherVoter = original otherVoter) →
            (rule.1 updated) 0 = candidate
\end{ReviewVerbatim}

\paragraph{Exact Lean library declaration}
\begin{ReviewVerbatim}
/--
Definition 4.2 of Ge et al. (2024), specialized to feasible finite profiles.
Elevating an output winner in one submitted ranking leaves it the output
winner after the update.
-/
def WinnerMonotonic {Voter : Type*} {n : ℕ} (feasible : Ranking n → Prop)
    (rule : LinearRankAggregationRule Voter n feasible) : Prop :=
  ∀ original updated voter candidate,
    FeasibleProfile feasible original → FeasibleProfile feasible updated →
      firstChoice (rule.run original) = candidate →
      ProfileElevatesCandidate original updated voter candidate →
        firstChoice (rule.run updated) = candidate
\end{ReviewVerbatim}

\paragraph{Recorded source-to-library screening}
\textbf{Verdict:} \texttt{matches}
\\
{\raggedright
\textbf{Reason:} Compared Definition 4.\allowbreak{}2 and its Theorems 4.\allowbreak{}3/\allowbreak{}C.\allowbreak{}6 uses.\allowbreak{} The target changes exactly one voter, weakly raises the named output winner's position, preserves the relative order of every other candidate and every other voter's ranking, requires both profiles source-\allowbreak{}feasible, and concludes that the same candidate remains first.\allowbreak{} This is the source elevation operation and winner conclusion.\allowbreak{}
\par}
\reviewmatch{Matches source input}
\noindent\textbf{Reviewer annotation}\par
\reviewerbox
\clearpage
\hypertarget{library-prerequisite-82fd59c97a4067c7}{}
\subsection*{\small\ttfamily\raggedright Applied\allowbreak{}Modeling\allowbreak{}Lib.\allowbreak{}Alignment.\allowbreak{}Axioms.\allowbreak{}canonical\allowbreak{}Kemeny\allowbreak{}Rule}
\textbf{Source locator:} {\footnotesize\raggedright cited publication, lines 1423-\allowbreak{}1449\par}
\paragraph{Verbatim paper-source connection}
\begin{ReviewVerbatim}
C.2   Linear Kemeny Rule

Next, we consider a different rule from social choice theory, the Kemeny rule, which can be trans-
formed to the linear setting while maintaining the separability can be achieved along with PMC.
Given an input profile π, the Kemeny rule returns a ranking σ ∗ that minimizes the total number of
pairwise disagreements with voters rankings, i.e.
                                            X X
                              σ ∗ ∈ arg min                1(b ≻σ a)
                                        σ∈S m
                                                 i∈V (a,b):a≻σi b
        m
where S contains all possible permutations of the m candidates. This expression can be equiva-
lently written as                        X
                          σ ∗ ∈ arg min        na≻b (π) · 1(b ≻σ a).
                                      σ∈S m
                                                a̸=b∈C
Here, we define the linear Kemeny rule, which outputs a parameter vector θ∗ that induces a ranking
that minimizes the total number of pairwise disagreements with voters’ rankings, i.e.,
                                        X
                          θ∗ ∈ arg min       na≻b (π) · 1(rθ (b) > rθ (a)).
                                 θ∈Rd    a̸=b∈C
Note that this rule conforms to the standard loss formulation, where the loss function is binary: it is
0 if two rankings agree with respect to the relative ranking of a pair of candidates and 1 otherwise.
Since binary loss is not convex, it does not fit in the impossibility result of Theorem 3.1.
Note that Kemeny is generally NP-hard to compute; however, even for linear Kemeny, there is
at least an exponential time aglorithm by brute-force computing the score of every ranking, and
determining whether or not it is feasible.
\end{ReviewVerbatim}



\paragraph{Lean-expanded library semantic target}
\begin{ReviewVerbatim}
type:
{n : ℕ} →
  (feasible : AppliedModelingLib.SocialChoice.Ranking.Ranking n → Prop) →
    (fallback : AppliedModelingLib.SocialChoice.Ranking.Ranking n) →
      feasible fallback →
        (voterCount : ℕ) → AppliedModelingLib.Alignment.Axioms.LinearRankAggregationRule (Fin voterCount) n feasible

value:
fun {n : ℕ} (feasible : AppliedModelingLib.SocialChoice.Ranking.Ranking n → Prop)
    (fallback : AppliedModelingLib.SocialChoice.Ranking.Ranking n) (hfallback : feasible fallback) (voterCount : ℕ) =>
  {
    run := fun (profile : AppliedModelingLib.Alignment.Axioms.RankingProfile (Fin voterCount) n) =>
      ↑(let minimizers : Finset (AppliedModelingLib.SocialChoice.Ranking.Ranking n) :=
          {output : AppliedModelingLib.SocialChoice.Ranking.Ranking n |
            feasible output ∧
              ∀ (contender : AppliedModelingLib.SocialChoice.Ranking.Ranking n),
                feasible contender →
                  AppliedModelingLib.Alignment.Axioms.kemenyDisagreement profile output ≤
                    AppliedModelingLib.Alignment.Axioms.kemenyDisagreement profile contender};
        have hminimizers : minimizers.Nonempty := ⋯;
        let tieKeys : Finset ℕ :=
          Finset.image
            (fun (ranking : AppliedModelingLib.SocialChoice.Ranking.Ranking n) =>
              ↑((Fintype.equivFin (AppliedModelingLib.SocialChoice.Ranking.Ranking n)) ranking))
            minimizers;
        have htieKeys : tieKeys.Nonempty := ⋯;
        let key : ℕ := tieKeys.min' htieKeys;
        have hkeyMem : key ∈ tieKeys := ⋯;
        let output : AppliedModelingLib.SocialChoice.Ranking.Ranking n := Classical.choose ⋯;
        have houtput :
          output ∈ minimizers ∧
            ↑((Fintype.equivFin (AppliedModelingLib.SocialChoice.Ranking.Ranking n)) output) = key :=
          ⋯;
        ⟨output, ⋯⟩),
    output_feasible := ⋯ }
\end{ReviewVerbatim}

\paragraph{Exact Lean library declaration}
\begin{ReviewVerbatim}
/-- The canonical finite Kemeny rule. -/
noncomputable def canonicalKemenyRule {n : ℕ}
    (feasible : Ranking n → Prop) (fallback : Ranking n) (hfallback : feasible fallback)
    (voterCount : ℕ) : LinearRankAggregationRule (Fin voterCount) n feasible where
  run profile := (canonicalKemenySelection feasible fallback hfallback profile).val
  output_feasible profile :=
    (canonicalKemenySelection feasible fallback hfallback profile).property.1.1
\end{ReviewVerbatim}

\paragraph{Recorded source-to-library screening}
\textbf{Verdict:} \texttt{matches}
\\
{\raggedright
\textbf{Reason:} Compared the linear-\allowbreak{}Kemeny rule model and Theorem C.\allowbreak{}3.\allowbreak{} The expanded rule chooses a feasible ranking minimizing the exact disagreement objective and then applies one profile-\allowbreak{}independent finite ranking key only among co-\allowbreak{}minimizers;\allowbreak{} this is a deterministic choice permitted by the source argmin, and the fixed key supplies the consistent selection used by the additive separability argument.\allowbreak{}
\par}
\reviewmatch{Matches source input}
\noindent\textbf{Reviewer annotation}\par
\reviewerbox
\clearpage
\hypertarget{library-prerequisite-8c57c8838b32b597}{}
\subsection*{\small\ttfamily\raggedright Applied\allowbreak{}Modeling\allowbreak{}Lib.\allowbreak{}Alignment.\allowbreak{}Axioms.\allowbreak{}kemeny\allowbreak{}Disagreement}
\textbf{Source locator:} {\footnotesize\raggedright cited publication, lines 1423-\allowbreak{}1449\par}
\paragraph{Verbatim paper-source connection}
\begin{ReviewVerbatim}
C.2   Linear Kemeny Rule

Next, we consider a different rule from social choice theory, the Kemeny rule, which can be trans-
formed to the linear setting while maintaining the separability can be achieved along with PMC.
Given an input profile π, the Kemeny rule returns a ranking σ ∗ that minimizes the total number of
pairwise disagreements with voters rankings, i.e.
                                            X X
                              σ ∗ ∈ arg min                1(b ≻σ a)
                                        σ∈S m
                                                 i∈V (a,b):a≻σi b
        m
where S contains all possible permutations of the m candidates. This expression can be equiva-
lently written as                        X
                          σ ∗ ∈ arg min        na≻b (π) · 1(b ≻σ a).
                                      σ∈S m
                                                a̸=b∈C
Here, we define the linear Kemeny rule, which outputs a parameter vector θ∗ that induces a ranking
that minimizes the total number of pairwise disagreements with voters’ rankings, i.e.,
                                        X
                          θ∗ ∈ arg min       na≻b (π) · 1(rθ (b) > rθ (a)).
                                 θ∈Rd    a̸=b∈C
Note that this rule conforms to the standard loss formulation, where the loss function is binary: it is
0 if two rankings agree with respect to the relative ranking of a pair of candidates and 1 otherwise.
Since binary loss is not convex, it does not fit in the impossibility result of Theorem 3.1.
Note that Kemeny is generally NP-hard to compute; however, even for linear Kemeny, there is
at least an exponential time aglorithm by brute-force computing the score of every ranking, and
determining whether or not it is feasible.
\end{ReviewVerbatim}



\paragraph{Lean-expanded library semantic target}
\begin{ReviewVerbatim}
type:
{Voter : Type u_1} →
  [Fintype Voter] →
    {n : ℕ} →
      AppliedModelingLib.Alignment.Axioms.RankingProfile Voter n → AppliedModelingLib.SocialChoice.Ranking.Ranking n → ℕ

value:
fun {Voter : Type u_1} [Fintype Voter] {n : ℕ} (profile : AppliedModelingLib.Alignment.Axioms.RankingProfile Voter n)
    (output : AppliedModelingLib.SocialChoice.Ranking.Ranking n) =>
  ∑ voter : Voter,
    ∑ first : AppliedModelingLib.SocialChoice.Ranking.Candidate n,
      ∑ second : AppliedModelingLib.SocialChoice.Ranking.Candidate n,
        if
            AppliedModelingLib.Alignment.Axioms.StrictlyPrefers (profile voter) first second ∧
              AppliedModelingLib.Alignment.Axioms.StrictlyPrefers output second first then
          1
        else 0
\end{ReviewVerbatim}

\paragraph{Exact Lean library declaration}
\begin{ReviewVerbatim}
/--
The Kemeny disagreement count of an output ranking against a finite profile.
For every ordered candidate pair submitted in one voter's ranking, it charges
one exactly when the output reverses that pair.  This is the source's
`n_{a ≻ b}(π) · 1(b ≻_σ a)` objective written as a direct voter sum.
-/
def kemenyDisagreement {Voter : Type*} [Fintype Voter] {n : ℕ}
    (profile : RankingProfile Voter n) (output : Ranking n) : ℕ :=
  ∑ voter, ∑ first, ∑ second,
    if StrictlyPrefers (profile voter) first second ∧ StrictlyPrefers output second first then 1
    else 0
\end{ReviewVerbatim}

\paragraph{Recorded source-to-library screening}
\textbf{Verdict:} \texttt{matches}
\\
{\raggedright
\textbf{Reason:} Compared both displayed Kemeny formulas and all C.\allowbreak{}3-\allowbreak{}C.\allowbreak{}5 uses.\allowbreak{} The target sums over voters and ordered candidate pairs and charges one exactly when the voter ranks first above second while the output ranks second above first.\allowbreak{} Regrouping by pair yields n\_\allowbreak{}\{a>b\}(pi)*1(b>\_\allowbreak{}sigma a), with self-\allowbreak{}pairs contributing zero, exactly the source disagreement count.\allowbreak{}
\par}
\reviewmatch{Matches source input}
\noindent\textbf{Reviewer annotation}\par
\reviewerbox
\clearpage
\hypertarget{library-prerequisite-4a32ddda17499d70}{}
\subsection*{\small\ttfamily\raggedright Applied\allowbreak{}Modeling\allowbreak{}Lib.\allowbreak{}Alignment.\allowbreak{}Axioms.\allowbreak{}linear\allowbreak{}Reward}
\textbf{Source locator:} {\footnotesize\raggedright cited publication, lines 158-\allowbreak{}198\par}
\paragraph{Verbatim paper-source connection}
\begin{ReviewVerbatim}
2     The Linear Social Choice Model
Let C be a set of m distinct prompt/responses, referred to as candidates, and let V = {1, . . . , n} be
a set of n human participants, known as voters. We denote by Rd the d-dimensional real space in
which both candidate feature vectors and chosen parameter vectors lie.
Each candidate c ∈ C is associated with a distinct feature vector xc ∈ Rd . A parameter vector
θ ∈ Rd induces a linear reward function rθ : C 7→ R defined by taking the dot product with feature
vectors rθ (c) = ⟨θ, xc ⟩. We will primarily be interested in how these parameterized functions rank
the candidates by reward. Let Ra≻b = {θ | rθ (a) ≥ rθ (b)} be the region where the reward of a is
at least as large as that of b. Note that Ra≻b and Rb≻a split Rd into two half spaces, separated by
the hyperplane orthogonal to xa − xb . Parameter vectors θ on the hyperplane have rθ (a) = rθ (b),
while rankings in the interior of either half-space strictly rank one over the other.
For a ranking σ over the candidates, we say that θ induces σ, denoted θ ▷ σ, if a ≻σ b implies
rθ (a) ≥ rθ (b). Let Rσ = {θ | θ ▷ σ} be the set of vectors θ thatTinduce it. Note that this can
be written as the intersection of corresponding half spaces Rσ = a,b:a≻σ b Ra≻b . Further, the


                                                    3

[form-feed in source extraction]
collection of {Rσ } essentially form a partition of Rd , covering the space and intersecting only at
their boundaries.
We call a θ non-degenerate if it is fully on one side of each of the separating hyperplanes, i.e.,
rθ (a) ̸= rθ (b) for all a, b ∈ C. Non-degenerate parameter vectors lie in the interior of some Rσ ,
and thus induce exactly one ranking. We call σ feasible if Rσ has a nonempty interior, i.e., is induced
by some nondegenerate θ.3
Each voter i ∈ V submits a ranking over the candidate σi . We assume that the feature space is
rich enough that voter preferences can be captured via non-degenerate parameter vectors. In other
words, we assume that each σi is feasible. We refer to the vector of voter rankings π = (σi )i∈V as
a profile. Further, for two candidates a, b, we write na≻b (π) := |{i ∈ V | a ≻σi b}| for the number
of voters that prefer a to b, and wa≻b (π) = na≻b (π)/n for the proportion of such voters. When the
profile π is clear from context, we may shorten these to na≻b and wa≻b , respectively.
We define a parameter aggregation rule as a function that takes as input a profile π and outputs a
parameter vector θ∗ . Our goal is to design parameter aggregation rules such that rθ∗ satisfies desir-
able properties with respect to the voter preferences. However, as the properties we care about will
only be with respect to how rθ∗ ranks the candidates, it will be more convenient to work with what
we call linear rank aggregation rules that take as input a profile π and output a feasible ranking σ.
There is a natural way to interpret a parameter aggregation rule as a linear rank aggregation rule,
namely, output any feasible ranking induced by θ∗ . The exact properties of the parameter aggrega-
tion rule could in principle be sensitive to the tie-breaking of non-degenerate outputs, however, all
of our results will be robust to such tie-breaking.4
We pay special attention to a prominent family of rules from social choice theory referred to as C1
rules [14], whose outputs depend only on majority relationships, i.e., they only need to know for
each pair of candidates (a, b) whether the majority prefers a or b.
\end{ReviewVerbatim}



\paragraph{Lean-expanded library semantic target}
\begin{ReviewVerbatim}
type:
{n dimension : ℕ} →
  AppliedModelingLib.Alignment.Axioms.LinearRewardParameter dimension →
    (AppliedModelingLib.SocialChoice.Ranking.Candidate n →
        AppliedModelingLib.Alignment.Axioms.FeatureVector dimension) →
      AppliedModelingLib.SocialChoice.Ranking.Candidate n → ℝ

value:
fun {n dimension : ℕ} (parameter : AppliedModelingLib.Alignment.Axioms.LinearRewardParameter dimension)
    (features :
      AppliedModelingLib.SocialChoice.Ranking.Candidate n → AppliedModelingLib.Alignment.Axioms.FeatureVector dimension)
    (candidate : AppliedModelingLib.SocialChoice.Ranking.Candidate n) =>
  ∑ coordinate : Fin dimension, parameter coordinate * features candidate coordinate
\end{ReviewVerbatim}

\paragraph{Exact Lean library declaration}
\begin{ReviewVerbatim}
/-- The dot-product reward of one candidate under a feature-linear model. -/
noncomputable def linearReward {n dimension : ℕ}
    (parameter : LinearRewardParameter dimension)
    (features : Candidate n → FeatureVector dimension)
    (candidate : Candidate n) : ℝ :=
  ∑ coordinate, parameter coordinate * features candidate coordinate
\end{ReviewVerbatim}

\paragraph{Recorded source-to-library screening}
\textbf{Verdict:} \texttt{matches}
\\
{\raggedright
\textbf{Reason:} The source defines r\_\allowbreak{}theta(c) as the dot product of theta and the candidate feature vector;\allowbreak{} the finite coordinate sum is exactly that dot product.\allowbreak{}
\par}
\reviewmatch{Matches source input}
\noindent\textbf{Reviewer annotation}\par
\reviewerbox
\clearpage
\hypertarget{library-prerequisite-8c0e6c25845ae78a}{}
\subsection*{\small\ttfamily\raggedright Applied\allowbreak{}Modeling\allowbreak{}Lib.\allowbreak{}Social\allowbreak{}Choice.\allowbreak{}Ranking.\allowbreak{}Candidate}
\textbf{Source locator:} {\footnotesize\raggedright cited publication, lines 158-\allowbreak{}198\par}
\paragraph{Verbatim paper-source connection}
\begin{ReviewVerbatim}
2     The Linear Social Choice Model
Let C be a set of m distinct prompt/responses, referred to as candidates, and let V = {1, . . . , n} be
a set of n human participants, known as voters. We denote by Rd the d-dimensional real space in
which both candidate feature vectors and chosen parameter vectors lie.
Each candidate c ∈ C is associated with a distinct feature vector xc ∈ Rd . A parameter vector
θ ∈ Rd induces a linear reward function rθ : C 7→ R defined by taking the dot product with feature
vectors rθ (c) = ⟨θ, xc ⟩. We will primarily be interested in how these parameterized functions rank
the candidates by reward. Let Ra≻b = {θ | rθ (a) ≥ rθ (b)} be the region where the reward of a is
at least as large as that of b. Note that Ra≻b and Rb≻a split Rd into two half spaces, separated by
the hyperplane orthogonal to xa − xb . Parameter vectors θ on the hyperplane have rθ (a) = rθ (b),
while rankings in the interior of either half-space strictly rank one over the other.
For a ranking σ over the candidates, we say that θ induces σ, denoted θ ▷ σ, if a ≻σ b implies
rθ (a) ≥ rθ (b). Let Rσ = {θ | θ ▷ σ} be the set of vectors θ thatTinduce it. Note that this can
be written as the intersection of corresponding half spaces Rσ = a,b:a≻σ b Ra≻b . Further, the


                                                    3

[form-feed in source extraction]
collection of {Rσ } essentially form a partition of Rd , covering the space and intersecting only at
their boundaries.
We call a θ non-degenerate if it is fully on one side of each of the separating hyperplanes, i.e.,
rθ (a) ̸= rθ (b) for all a, b ∈ C. Non-degenerate parameter vectors lie in the interior of some Rσ ,
and thus induce exactly one ranking. We call σ feasible if Rσ has a nonempty interior, i.e., is induced
by some nondegenerate θ.3
Each voter i ∈ V submits a ranking over the candidate σi . We assume that the feature space is
rich enough that voter preferences can be captured via non-degenerate parameter vectors. In other
words, we assume that each σi is feasible. We refer to the vector of voter rankings π = (σi )i∈V as
a profile. Further, for two candidates a, b, we write na≻b (π) := |{i ∈ V | a ≻σi b}| for the number
of voters that prefer a to b, and wa≻b (π) = na≻b (π)/n for the proportion of such voters. When the
profile π is clear from context, we may shorten these to na≻b and wa≻b , respectively.
We define a parameter aggregation rule as a function that takes as input a profile π and outputs a
parameter vector θ∗ . Our goal is to design parameter aggregation rules such that rθ∗ satisfies desir-
able properties with respect to the voter preferences. However, as the properties we care about will
only be with respect to how rθ∗ ranks the candidates, it will be more convenient to work with what
we call linear rank aggregation rules that take as input a profile π and output a feasible ranking σ.
There is a natural way to interpret a parameter aggregation rule as a linear rank aggregation rule,
namely, output any feasible ranking induced by θ∗ . The exact properties of the parameter aggrega-
tion rule could in principle be sensitive to the tie-breaking of non-degenerate outputs, however, all
of our results will be robust to such tie-breaking.4
We pay special attention to a prominent family of rules from social choice theory referred to as C1
rules [14], whose outputs depend only on majority relationships, i.e., they only need to know for
each pair of candidates (a, b) whether the majority prefers a or b.
\end{ReviewVerbatim}



\paragraph{Lean-expanded library semantic target}
\begin{ReviewVerbatim}
type:
ℕ → Type

value:
fun (n : ℕ) => Fin (n + 2)
\end{ReviewVerbatim}

\paragraph{Exact Lean library declaration}
\begin{ReviewVerbatim}
/--
Candidate set with at least two candidates. The parameter `n` means there are
`n + 2` candidates.
-/
abbrev Candidate (n : ℕ) := Fin (n + 2)
\end{ReviewVerbatim}

\paragraph{Recorded source-to-library screening}
\textbf{Verdict:} \texttt{matches}
\\
{\raggedright
\textbf{Reason:} The source uses a finite set C of distinct candidates;\allowbreak{} Fin (n + 2) is a concrete finite indexing of such a candidate set and adds no preference or reward claim.\allowbreak{}
\par}
\reviewmatch{Matches source input}
\noindent\textbf{Reviewer annotation}\par
\reviewerbox
\clearpage
\hypertarget{library-prerequisite-5cde54446d1ae021}{}
\subsection*{\small\ttfamily\raggedright Applied\allowbreak{}Modeling\allowbreak{}Lib.\allowbreak{}Social\allowbreak{}Choice.\allowbreak{}Ranking.\allowbreak{}Ranking}
\textbf{Source locator:} {\footnotesize\raggedright cited publication, lines 158-\allowbreak{}198\par}
\paragraph{Verbatim paper-source connection}
\begin{ReviewVerbatim}
2     The Linear Social Choice Model
Let C be a set of m distinct prompt/responses, referred to as candidates, and let V = {1, . . . , n} be
a set of n human participants, known as voters. We denote by Rd the d-dimensional real space in
which both candidate feature vectors and chosen parameter vectors lie.
Each candidate c ∈ C is associated with a distinct feature vector xc ∈ Rd . A parameter vector
θ ∈ Rd induces a linear reward function rθ : C 7→ R defined by taking the dot product with feature
vectors rθ (c) = ⟨θ, xc ⟩. We will primarily be interested in how these parameterized functions rank
the candidates by reward. Let Ra≻b = {θ | rθ (a) ≥ rθ (b)} be the region where the reward of a is
at least as large as that of b. Note that Ra≻b and Rb≻a split Rd into two half spaces, separated by
the hyperplane orthogonal to xa − xb . Parameter vectors θ on the hyperplane have rθ (a) = rθ (b),
while rankings in the interior of either half-space strictly rank one over the other.
For a ranking σ over the candidates, we say that θ induces σ, denoted θ ▷ σ, if a ≻σ b implies
rθ (a) ≥ rθ (b). Let Rσ = {θ | θ ▷ σ} be the set of vectors θ thatTinduce it. Note that this can
be written as the intersection of corresponding half spaces Rσ = a,b:a≻σ b Ra≻b . Further, the


                                                    3

[form-feed in source extraction]
collection of {Rσ } essentially form a partition of Rd , covering the space and intersecting only at
their boundaries.
We call a θ non-degenerate if it is fully on one side of each of the separating hyperplanes, i.e.,
rθ (a) ̸= rθ (b) for all a, b ∈ C. Non-degenerate parameter vectors lie in the interior of some Rσ ,
and thus induce exactly one ranking. We call σ feasible if Rσ has a nonempty interior, i.e., is induced
by some nondegenerate θ.3
Each voter i ∈ V submits a ranking over the candidate σi . We assume that the feature space is
rich enough that voter preferences can be captured via non-degenerate parameter vectors. In other
words, we assume that each σi is feasible. We refer to the vector of voter rankings π = (σi )i∈V as
a profile. Further, for two candidates a, b, we write na≻b (π) := |{i ∈ V | a ≻σi b}| for the number
of voters that prefer a to b, and wa≻b (π) = na≻b (π)/n for the proportion of such voters. When the
profile π is clear from context, we may shorten these to na≻b and wa≻b , respectively.
We define a parameter aggregation rule as a function that takes as input a profile π and outputs a
parameter vector θ∗ . Our goal is to design parameter aggregation rules such that rθ∗ satisfies desir-
able properties with respect to the voter preferences. However, as the properties we care about will
only be with respect to how rθ∗ ranks the candidates, it will be more convenient to work with what
we call linear rank aggregation rules that take as input a profile π and output a feasible ranking σ.
There is a natural way to interpret a parameter aggregation rule as a linear rank aggregation rule,
namely, output any feasible ranking induced by θ∗ . The exact properties of the parameter aggrega-
tion rule could in principle be sensitive to the tie-breaking of non-degenerate outputs, however, all
of our results will be robust to such tie-breaking.4
We pay special attention to a prominent family of rules from social choice theory referred to as C1
rules [14], whose outputs depend only on majority relationships, i.e., they only need to know for
each pair of candidates (a, b) whether the majority prefers a or b.
\end{ReviewVerbatim}



\paragraph{Lean-expanded library semantic target}
\begin{ReviewVerbatim}
type:
ℕ → Type

value:
fun (n : ℕ) => Equiv.Perm (AppliedModelingLib.SocialChoice.Ranking.Candidate n)
\end{ReviewVerbatim}

\paragraph{Exact Lean library declaration}
\begin{ReviewVerbatim}
/-- A ranking is a permutation of the candidate set. -/
abbrev Ranking (n : ℕ) := Equiv.Perm (Candidate n)
\end{ReviewVerbatim}

\paragraph{Recorded source-to-library screening}
\textbf{Verdict:} \texttt{matches}
\\
{\raggedright
\textbf{Reason:} The source treats rankings as permutations of the m candidates;\allowbreak{} Equiv.\allowbreak{}Perm Candidate is exactly a permutation ranking.\allowbreak{}
\par}
\reviewmatch{Matches source input}
\noindent\textbf{Reviewer annotation}\par
\reviewerbox
\clearpage
\hypertarget{paper-prerequisite-section-5ef5ef0364b6939c}{}
\section*{Paper-specific semantic prerequisites}
\hypertarget{paper-prerequisite-89cd467f0989921f}{}
\subsection*{\small\ttfamily\raggedright Ge\allowbreak{}Et\allowbreak{}Al2024\allowbreak{}Alignment\allowbreak{}Axioms.\allowbreak{}Is\allowbreak{}Fixed\allowbreak{}Tie\allowbreak{}Leximax\allowbreak{}Plurality\allowbreak{}Selector}
\noindent\textbf{Paper-local declaration:}\par
{\footnotesize\ttfamily\raggedright Ge\allowbreak{}Et\allowbreak{}Al2024\allowbreak{}Alignment\allowbreak{}Axioms.\allowbreak{}Is\allowbreak{}Fixed\allowbreak{}Tie\allowbreak{}Leximax\allowbreak{}Plurality\allowbreak{}Selector\par}
\textbf{Source locator:} {\footnotesize\raggedright cited publication, lines 1530-\allowbreak{}1539\par}
\paragraph{Verbatim paper-source connection}
\begin{ReviewVerbatim}
C.3   Leximax Plurality

Plurality is probably the most ubiquitous voting rule in the world. Its ranking variant ranks the can-
didates in decreasing order with respect to their plurality scores. The plurality score of a candidate
is equal to the number of her appearances in the first position. This rule is known to satisfy several
axioms but in linear social choice cannot be directly applied, as not all rankings are feasibly.
Similarly to leximax Copeland, we define Leximax Plurality as follows. It ranks first the candidate
with the highest plurality score that can be ranked first under some parameter vector. Subject to
this first position, it ranks second the candidate with the highest plurality score that can feasibly be
ranked second, and so on, until all the positions are filled.
\end{ReviewVerbatim}



\paragraph{Lean-elaborated paper signature}
\begin{ReviewVerbatim}
type:
{Voter : Type u_1} →
  [Fintype Voter] →
    {n : ℕ} →
      (feasible : AppliedModelingLib.SocialChoice.Ranking.Ranking n → Prop) →
        AppliedModelingLib.Alignment.Axioms.LinearRankAggregationRule Voter n feasible → Prop

value:
fun {Voter : Type u_1} [Fintype Voter] {n : ℕ} (feasible : AppliedModelingLib.SocialChoice.Ranking.Ranking n → Prop)
    (rule : AppliedModelingLib.Alignment.Axioms.LinearRankAggregationRule Voter n feasible) =>
  ∀ (profile : AppliedModelingLib.Alignment.Axioms.RankingProfile Voter n),
    feasible (rule.1 profile) ∧
      ∀ (position candidate : AppliedModelingLib.SocialChoice.Ranking.Candidate n),
        (∃ (contender : AppliedModelingLib.SocialChoice.Ranking.Ranking n),
            feasible contender ∧
              (∀ earlier < position, contender earlier = (rule.1 profile) earlier) ∧ contender position = candidate) →
          ¬(((∑ voter : Voter, if (profile voter) 0 = (rule.1 profile) position then 1 else 0) <
                ∑ voter : Voter, if (profile voter) 0 = candidate then 1 else 0) ∨
              ((∑ voter : Voter, if (profile voter) 0 = candidate then 1 else 0) =
                  ∑ voter : Voter, if (profile voter) 0 = (rule.1 profile) position then 1 else 0) ∧
                candidate < (rule.1 profile) position)
\end{ReviewVerbatim}


\paragraph{Recorded source-to-declaration screening}
\textbf{Verdict:} \texttt{matches}
\\
{\raggedright
\textbf{Reason:} The expanded predicate requires a feasible output and, at every prefix, rejects any feasible next candidate with a higher plurality score or an equal score and lower fixed index.\allowbreak{} This is the source's sequential highest-\allowbreak{}feasible-\allowbreak{}plurality construction with one deterministic choice among tied score maximizers.\allowbreak{}
\par}
\reviewmatch{Matches source input}
\noindent\textbf{Reviewer annotation}\par
\reviewerbox
\clearpage
\hypertarget{paper-prerequisite-c8ae4ba0fedf8db8}{}
\subsection*{\small\ttfamily\raggedright Ge\allowbreak{}Et\allowbreak{}Al2024\allowbreak{}Alignment\allowbreak{}Axioms.\allowbreak{}Is\allowbreak{}Majority\allowbreak{}Loss\allowbreak{}Minimizing}
\noindent\textbf{Paper-local declaration:}\par
{\footnotesize\ttfamily\raggedright Ge\allowbreak{}Et\allowbreak{}Al2024\allowbreak{}Alignment\allowbreak{}Axioms.\allowbreak{}Is\allowbreak{}Majority\allowbreak{}Loss\allowbreak{}Minimizing\par}
\textbf{Source locator:} {\footnotesize\raggedright cited publication, lines 370-\allowbreak{}386\par}
\paragraph{Verbatim paper-source connection}
\begin{ReviewVerbatim}
3.2   Majority-Based Loss Formulation

Despite the negative results for loss-function-based rules, we may hope for a remedy using slightly
different information. Specifically, we consider a similar loss-based function that rather than getting
penalized for disagreeing with each voter only gets penalized if it disagrees with a majority of voters.
That is, we choose θ minimizing
                                       X
                    Lmaj (θ; π, ell) =         I[wa≻b (π) > 1/2] · ell(rθ(b) − rθ(a) ).
                                       a̸=b∈C

Defining a parameter aggregation function based on this loss suffers from the same caveat as before,
that in some cases no optimal θ exists. Nevertheless, we can apply an analogous fix for a ranking
variant. We say that a linear rank aggregation rule f minimizes ell in the majority formulation if for
all σ = f (π),
                               inf σ Lmaj (θ; π, ell) = inf Lmaj (θ; π, ell).
                               θ:θ∈R                      θ
\end{ReviewVerbatim}



\paragraph{Lean-elaborated paper signature}
\begin{ReviewVerbatim}
type:
{Voter : Type u_1} →
  [Fintype Voter] →
    {n dimension : ℕ} →
      (features :
          AppliedModelingLib.SocialChoice.Ranking.Candidate n →
            AppliedModelingLib.Alignment.Axioms.FeatureVector dimension) →
        (ℝ → ℝ) →
          AppliedModelingLib.Alignment.Axioms.LinearRankAggregationRule Voter n
              (AppliedModelingLib.Alignment.Axioms.LinearFeasibleRanking features) →
            Prop

value:
fun {Voter : Type u_1} [Fintype Voter] {n dimension : ℕ}
    (features :
      AppliedModelingLib.SocialChoice.Ranking.Candidate n → AppliedModelingLib.Alignment.Axioms.FeatureVector dimension)
    (loss : ℝ → ℝ)
    (rule :
      AppliedModelingLib.Alignment.Axioms.LinearRankAggregationRule Voter n
        (AppliedModelingLib.Alignment.Axioms.LinearFeasibleRanking features)) =>
  ∀ (profile : AppliedModelingLib.Alignment.Axioms.RankingProfile Voter n),
    AppliedModelingLib.Alignment.Axioms.FeasibleProfile
        (AppliedModelingLib.Alignment.Axioms.LinearFeasibleRanking features) profile →
      GeEtAl2024AlignmentAxioms.restrictedMajorityLossInfimum features profile loss (rule.1 profile) =
        GeEtAl2024AlignmentAxioms.globalMajorityLossInfimum features profile loss
\end{ReviewVerbatim}


\paragraph{Recorded source-to-declaration screening}
\textbf{Verdict:} \texttt{matches}
\\
{\raggedright
\textbf{Reason:} For every feasible profile, the expanded target equates the infimum of the majority-\allowbreak{}loss objective over parameters inducing the rule output with its unrestricted infimum.\allowbreak{} This is the source's definition of a rule minimizing the majority formulation.\allowbreak{}
\par}
\reviewmatch{Matches source input}
\noindent\textbf{Reviewer annotation}\par
\reviewerbox
\clearpage
\hypertarget{paper-prerequisite-fa6ae496ed9c0c39}{}
\subsection*{\small\ttfamily\raggedright Ge\allowbreak{}Et\allowbreak{}Al2024\allowbreak{}Alignment\allowbreak{}Axioms.\allowbreak{}Is\allowbreak{}Pareto\allowbreak{}Kemeny\allowbreak{}Selector}
\noindent\textbf{Paper-local declaration:}\par
{\footnotesize\ttfamily\raggedright Ge\allowbreak{}Et\allowbreak{}Al2024\allowbreak{}Alignment\allowbreak{}Axioms.\allowbreak{}Is\allowbreak{}Pareto\allowbreak{}Kemeny\allowbreak{}Selector\par}
\textbf{Source locator:} {\footnotesize\raggedright cited publication, lines 1506-\allowbreak{}1508\par}
\paragraph{Verbatim paper-source connection}
\begin{ReviewVerbatim}
A possibly easy fix to the problem that linear Kemeny does not satisfy PO would be to enforce the
PO criterion, as we did for the LCPO, i.e., to restrict to parameter vectors that respect PO. However,
we show that if we do that, then linear Kemeny subject to PO does not satisfy separability anymore.
\end{ReviewVerbatim}



\paragraph{Lean-elaborated paper signature}
\begin{ReviewVerbatim}
type:
((voterCount : ℕ) →
    AppliedModelingLib.Alignment.Axioms.LinearRankAggregationRule (Fin voterCount) 18
      (AppliedModelingLib.Alignment.Axioms.LinearFeasibleRanking
        ![![2000000, 0, 0, 0, 0, 0, 0], ![0, 2000000, 0, 0, 0, 0, 0], ![0, 200000, 0, 0, 0, 0, 0],
          ![0, 100000, 100000, 0, 0, 0, 0], ![0, 0, 200000, 0, 0, 0, 0], ![0, 0, 20000, 0, 0, 0, 0],
          ![0, 0, 10000, 10000, 0, 0, 0], ![0, 0, 0, 20000, 0, 0, 0], ![0, 0, 0, 2000, 0, 0, 0],
          ![0, 0, 0, 1000, 1000, 0, 0], ![0, 0, 0, 0, 2000, 0, 0], ![0, 0, 0, 0, 200, 0, 0], ![0, 0, 0, 0, 100, 100, 0],
          ![0, 0, 0, 0, 0, 200, 0], ![0, 0, 0, 0, 0, 20, 0], ![0, 0, 0, 0, 0, 10, 10], ![0, 0, 0, 0, 0, 0, 20],
          ![0, 0, 0, 0, 0, 0, 2], ![1, 0, 0, 0, 0, 0, 1], ![2, 0, 0, 0, 0, 0, 0]])) →
  Prop

value:
fun
    (rule :
      (voterCount : ℕ) →
        AppliedModelingLib.Alignment.Axioms.LinearRankAggregationRule (Fin voterCount) 18
          (AppliedModelingLib.Alignment.Axioms.LinearFeasibleRanking
            ![![2000000, 0, 0, 0, 0, 0, 0], ![0, 2000000, 0, 0, 0, 0, 0], ![0, 200000, 0, 0, 0, 0, 0],
              ![0, 100000, 100000, 0, 0, 0, 0], ![0, 0, 200000, 0, 0, 0, 0], ![0, 0, 20000, 0, 0, 0, 0],
              ![0, 0, 10000, 10000, 0, 0, 0], ![0, 0, 0, 20000, 0, 0, 0], ![0, 0, 0, 2000, 0, 0, 0],
              ![0, 0, 0, 1000, 1000, 0, 0], ![0, 0, 0, 0, 2000, 0, 0], ![0, 0, 0, 0, 200, 0, 0],
              ![0, 0, 0, 0, 100, 100, 0], ![0, 0, 0, 0, 0, 200, 0], ![0, 0, 0, 0, 0, 20, 0], ![0, 0, 0, 0, 0, 10, 10],
              ![0, 0, 0, 0, 0, 0, 20], ![0, 0, 0, 0, 0, 0, 2], ![1, 0, 0, 0, 0, 0, 1], ![2, 0, 0, 0, 0, 0, 0]])) =>
  ∀ (voterCount : ℕ) (profile : AppliedModelingLib.Alignment.Axioms.RankingProfile (Fin voterCount) 18),
    AppliedModelingLib.Alignment.Axioms.FeasibleProfile
        (AppliedModelingLib.Alignment.Axioms.LinearFeasibleRanking
          ![![2000000, 0, 0, 0, 0, 0, 0], ![0, 2000000, 0, 0, 0, 0, 0], ![0, 200000, 0, 0, 0, 0, 0],
            ![0, 100000, 100000, 0, 0, 0, 0], ![0, 0, 200000, 0, 0, 0, 0], ![0, 0, 20000, 0, 0, 0, 0],
            ![0, 0, 10000, 10000, 0, 0, 0], ![0, 0, 0, 20000, 0, 0, 0], ![0, 0, 0, 2000, 0, 0, 0],
            ![0, 0, 0, 1000, 1000, 0, 0], ![0, 0, 0, 0, 2000, 0, 0], ![0, 0, 0, 0, 200, 0, 0],
            ![0, 0, 0, 0, 100, 100, 0], ![0, 0, 0, 0, 0, 200, 0], ![0, 0, 0, 0, 0, 20, 0], ![0, 0, 0, 0, 0, 10, 10],
            ![0, 0, 0, 0, 0, 0, 20], ![0, 0, 0, 0, 0, 0, 2], ![1, 0, 0, 0, 0, 0, 1], ![2, 0, 0, 0, 0, 0, 0]])
        profile →
      (AppliedModelingLib.Alignment.Axioms.LinearFeasibleRanking
            ![![2000000, 0, 0, 0, 0, 0, 0], ![0, 2000000, 0, 0, 0, 0, 0], ![0, 200000, 0, 0, 0, 0, 0],
              ![0, 100000, 100000, 0, 0, 0, 0], ![0, 0, 200000, 0, 0, 0, 0], ![0, 0, 20000, 0, 0, 0, 0],
              ![0, 0, 10000, 10000, 0, 0, 0], ![0, 0, 0, 20000, 0, 0, 0], ![0, 0, 0, 2000, 0, 0, 0],
              ![0, 0, 0, 1000, 1000, 0, 0], ![0, 0, 0, 0, 2000, 0, 0], ![0, 0, 0, 0, 200, 0, 0],
              ![0, 0, 0, 0, 100, 100, 0], ![0, 0, 0, 0, 0, 200, 0], ![0, 0, 0, 0, 0, 20, 0], ![0, 0, 0, 0, 0, 10, 10],
              ![0, 0, 0, 0, 0, 0, 20], ![0, 0, 0, 0, 0, 0, 2], ![1, 0, 0, 0, 0, 0, 1], ![2, 0, 0, 0, 0, 0, 0]]
            ((rule voterCount).1 profile) ∧
          ∀ (first second : AppliedModelingLib.SocialChoice.Ranking.Candidate 18),
            (∀ (voter : Fin voterCount),
                AppliedModelingLib.Alignment.Axioms.StrictlyPrefers (profile voter) first second) →
              AppliedModelingLib.Alignment.Axioms.StrictlyPrefers ((rule voterCount).1 profile) first second) ∧
        ∀ (contender : AppliedModelingLib.SocialChoice.Ranking.Ranking 18),
          (AppliedModelingLib.Alignment.Axioms.LinearFeasibleRanking
                ![![2000000, 0, 0, 0, 0, 0, 0], ![0, 2000000, 0, 0, 0, 0, 0], ![0, 200000, 0, 0, 0, 0, 0],
                  ![0, 100000, 100000, 0, 0, 0, 0], ![0, 0, 200000, 0, 0, 0, 0], ![0, 0, 20000, 0, 0, 0, 0],
                  ![0, 0, 10000, 10000, 0, 0, 0], ![0, 0, 0, 20000, 0, 0, 0], ![0, 0, 0, 2000, 0, 0, 0],
                  ![0, 0, 0, 1000, 1000, 0, 0], ![0, 0, 0, 0, 2000, 0, 0], ![0, 0, 0, 0, 200, 0, 0],
                  ![0, 0, 0, 0, 100, 100, 0], ![0, 0, 0, 0, 0, 200, 0], ![0, 0, 0, 0, 0, 20, 0],
                  ![0, 0, 0, 0, 0, 10, 10], ![0, 0, 0, 0, 0, 0, 20], ![0, 0, 0, 0, 0, 0, 2], ![1, 0, 0, 0, 0, 0, 1],
                  ![2, 0, 0, 0, 0, 0, 0]]
                contender ∧
              ∀ (first second : AppliedModelingLib.SocialChoice.Ranking.Candidate 18),
                (∀ (voter : Fin voterCount),
                    AppliedModelingLib.Alignment.Axioms.StrictlyPrefers (profile voter) first second) →
                  AppliedModelingLib.Alignment.Axioms.StrictlyPrefers contender first second) →
            AppliedModelingLib.Alignment.Axioms.kemenyDisagreement profile ((rule voterCount).1 profile) ≤
              AppliedModelingLib.Alignment.Axioms.kemenyDisagreement profile contender
\end{ReviewVerbatim}


\paragraph{Recorded source-to-declaration screening}
\textbf{Verdict:} \texttt{matches}
\\
{\raggedright
\textbf{Reason:} For the displayed source feature instance, the expanded target requires the output to be feasible, to respect every unanimous comparison, and to minimize Kemeny disagreement among all such rankings on every feasible profile.\allowbreak{} This is exactly Kemeny restricted to PO-\allowbreak{}respecting parameter-\allowbreak{}induced rankings;\allowbreak{} it adds no global tie rule.\allowbreak{}
\par}
\reviewmatch{Matches source input}
\noindent\textbf{Reviewer annotation}\par
\reviewerbox
\clearpage
\hypertarget{paper-prerequisite-fc38b98f65c4e72f}{}
\subsection*{\small\ttfamily\raggedright Ge\allowbreak{}Et\allowbreak{}Al2024\allowbreak{}Alignment\allowbreak{}Axioms.\allowbreak{}fixed\allowbreak{}Tie\allowbreak{}LCPO}
\noindent\textbf{Paper-local declaration:}\par
{\footnotesize\ttfamily\raggedright Ge\allowbreak{}Et\allowbreak{}Al2024\allowbreak{}Alignment\allowbreak{}Axioms.\allowbreak{}fixed\allowbreak{}Tie\allowbreak{}LCPO\par}
\textbf{Source locator:} {\footnotesize\raggedright cited publication, lines 416-\allowbreak{}438\par}
\paragraph{Verbatim paper-source connection}
\begin{ReviewVerbatim}
The Copeland rule assigns a Copeland score to each alternative equal to the number of other alter-
natives it beats in a pairwise competition, i.e., the score for a is |{b | wa≻b > 1/2}|. It then ranks
the candidates in descending order according to their Copeland scores (breaking ties arbitrarily). It


                                                     7

[form-feed in source extraction]
is known that Copeland satisfies PO, PMC, and additional axiomatic properties. However, in linear
social choice, since not every ranking is feasible, we cannot always output the Copeland ranking.
We, therefore, define a new linear rank aggregation rule, which we call leximax Copeland. This
rule chooses a feasible ranking as follows. It ranks first the candidate with the highest Copeland
score that can be feasibly ranked first under some parameter vector θ. Subject to this first position,
it ranks second the candidate with the highest Copeland score which can be feasibly ranked second,
and continues this process for subsequent positions.
Copeland’s rule is a C1 rule because it only requires the majority relationships between the can-
didates. Analogously, leximax Copeland is also a C1 linear rank aggregation rule. Therefore, by
Theorem 3.7, it does not satisfy the PO criterion. To address this issue, we define a variant called
leximax Copeland subject to PO (LCPO), which incorporates the PO criterion. Under LCPO, for
every pair of alternatives where one dominates the other, the rule restricts rankings to place the
dominating alternative above the dominated one.
The rule remains well-defined since the set of feasible rankings when enforcing the PO criterion is
non-empty, as whenever a dominates b, all the rankings in the input profile rank a above b. Note that
if the Copeland ranking is feasible, then this rule outputs that ranking, since unrestricted Copeland
satisfies PO.
\end{ReviewVerbatim}



\paragraph{Lean-elaborated paper signature}
\begin{ReviewVerbatim}
type:
{Voter : Type u_1} →
  [Fintype Voter] →
    [Nonempty Voter] →
      {n : ℕ} →
        (feasible : AppliedModelingLib.SocialChoice.Ranking.Ranking n → Prop) →
          (fallback : AppliedModelingLib.SocialChoice.Ranking.Ranking n) →
            feasible fallback → AppliedModelingLib.Alignment.Axioms.LinearRankAggregationRule Voter n feasible

value:
fun {Voter : Type u_1} [Fintype Voter] [Nonempty Voter] {n : ℕ}
    (feasible : AppliedModelingLib.SocialChoice.Ranking.Ranking n → Prop)
    (fallback : AppliedModelingLib.SocialChoice.Ranking.Ranking n) (hfallback : feasible fallback) =>
  {
    run := fun (profile : AppliedModelingLib.Alignment.Axioms.RankingProfile Voter n) =>
      if hprofile : AppliedModelingLib.Alignment.Axioms.FeasibleProfile feasible profile then
        ↑(let searchDomain : Finset (AppliedModelingLib.SocialChoice.Ranking.Ranking n) :=
            {ranking : AppliedModelingLib.SocialChoice.Ranking.Ranking n |
              feasible ranking ∧
                ∀ (first second : AppliedModelingLib.SocialChoice.Ranking.Candidate n),
                  (∀ (voter : Voter),
                      AppliedModelingLib.Alignment.Axioms.StrictlyPrefers (profile voter) first second) →
                    AppliedModelingLib.Alignment.Axioms.StrictlyPrefers ranking first second};
          let keyVectors : Finset (Lex (AppliedModelingLib.SocialChoice.Ranking.Candidate n → ℕ)) :=
            Finset.image
              (fun (ranking : AppliedModelingLib.SocialChoice.Ranking.Ranking n) =>
                toLex fun (position : AppliedModelingLib.SocialChoice.Ranking.Candidate n) =>
                  {opponent : AppliedModelingLib.SocialChoice.Ranking.Candidate n |
                          Fintype.card Voter <
                            2 *
                              {voter : Voter |
                                  AppliedModelingLib.Alignment.Axioms.StrictlyPrefers (profile voter) (ranking position)
                                    opponent}.card}.card *
                      Fintype.card (AppliedModelingLib.SocialChoice.Ranking.Candidate n) +
                    (Fintype.card (AppliedModelingLib.SocialChoice.Ranking.Candidate n) - 1 - ↑(ranking position)))
              searchDomain;
          have hsearchDomain : searchDomain.Nonempty := ⋯;
          have hkeyVectors : keyVectors.Nonempty := ⋯;
          have hExists :
            ∃ ranking ∈ searchDomain,
              (toLex fun (position : AppliedModelingLib.SocialChoice.Ranking.Candidate n) =>
                  {opponent : AppliedModelingLib.SocialChoice.Ranking.Candidate n |
                          Fintype.card Voter <
                            2 *
                              {voter : Voter |
                                  AppliedModelingLib.Alignment.Axioms.StrictlyPrefers (profile voter) (ranking position)
                                    opponent}.card}.card *
                      Fintype.card (AppliedModelingLib.SocialChoice.Ranking.Candidate n) +
                    (Fintype.card (AppliedModelingLib.SocialChoice.Ranking.Candidate n) - 1 - ↑(ranking position))) =
                let searchDomain : Finset (AppliedModelingLib.SocialChoice.Ranking.Ranking n) :=
                  {ranking : AppliedModelingLib.SocialChoice.Ranking.Ranking n |
                    feasible ranking ∧
                      ∀ (first second : AppliedModelingLib.SocialChoice.Ranking.Candidate n),
                        (∀ (voter : Voter),
                            AppliedModelingLib.Alignment.Axioms.StrictlyPrefers (profile voter) first second) →
                          AppliedModelingLib.Alignment.Axioms.StrictlyPrefers ranking first second};
                let keyVectors : Finset (Lex (AppliedModelingLib.SocialChoice.Ranking.Candidate n → ℕ)) :=
                  Finset.image
                    (fun (ranking : AppliedModelingLib.SocialChoice.Ranking.Ranking n) =>
                      toLex fun (position : AppliedModelingLib.SocialChoice.Ranking.Candidate n) =>
                        {opponent : AppliedModelingLib.SocialChoice.Ranking.Candidate n |
                                Fintype.card Voter <
                                  2 *
                                    {voter : Voter |
                                        AppliedModelingLib.Alignment.Axioms.StrictlyPrefers (profile voter)
                                          (ranking position) opponent}.card}.card *
                            Fintype.card (AppliedModelingLib.SocialChoice.Ranking.Candidate n) +
                          (Fintype.card (AppliedModelingLib.SocialChoice.Ranking.Candidate n) - 1 -
                            ↑(ranking position)))
                    searchDomain;
                have hsearchDomain : searchDomain.Nonempty := ⋯;
                have hkeyVectors : keyVectors.Nonempty := ⋯;
                keyVectors.max' hkeyVectors :=
            ⋯;
          ⟨Classical.choose hExists, ⋯⟩)
      else fallback,
    output_feasible := ⋯ }
\end{ReviewVerbatim}


\paragraph{Recorded source-to-declaration screening}
\textbf{Verdict:} \texttt{matches}
\\
{\raggedright
\textbf{Reason:} On source-\allowbreak{}feasible profiles the expanded run searches rankings that are both feasible and Pareto-\allowbreak{}respecting, then lexicographically maximizes the positionwise Copeland keys.\allowbreak{} This implements the source's LCPO restriction and sequential highest-\allowbreak{}Copeland selection;\allowbreak{} its candidate-\allowbreak{}order tie break instantiates the source's arbitrary ties.\allowbreak{}
\par}
\reviewmatch{Matches source input}
\noindent\textbf{Reviewer annotation}\par
\reviewerbox
\clearpage
\hypertarget{paper-prerequisite-e532c12e5479137f}{}
\subsection*{\small\ttfamily\raggedright Ge\allowbreak{}Et\allowbreak{}Al2024\allowbreak{}Alignment\allowbreak{}Axioms.\allowbreak{}global\allowbreak{}Majority\allowbreak{}Loss\allowbreak{}Infimum}
\noindent\textbf{Paper-local declaration:}\par
{\footnotesize\ttfamily\raggedright Ge\allowbreak{}Et\allowbreak{}Al2024\allowbreak{}Alignment\allowbreak{}Axioms.\allowbreak{}global\allowbreak{}Majority\allowbreak{}Loss\allowbreak{}Infimum\par}
\textbf{Source locator:} {\footnotesize\raggedright cited publication, lines 370-\allowbreak{}386\par}
\paragraph{Verbatim paper-source connection}
\begin{ReviewVerbatim}
3.2   Majority-Based Loss Formulation

Despite the negative results for loss-function-based rules, we may hope for a remedy using slightly
different information. Specifically, we consider a similar loss-based function that rather than getting
penalized for disagreeing with each voter only gets penalized if it disagrees with a majority of voters.
That is, we choose θ minimizing
                                       X
                    Lmaj (θ; π, ell) =         I[wa≻b (π) > 1/2] · ell(rθ(b) − rθ(a) ).
                                       a̸=b∈C

Defining a parameter aggregation function based on this loss suffers from the same caveat as before,
that in some cases no optimal θ exists. Nevertheless, we can apply an analogous fix for a ranking
variant. We say that a linear rank aggregation rule f minimizes ell in the majority formulation if for
all σ = f (π),
                               inf σ Lmaj (θ; π, ell) = inf Lmaj (θ; π, ell).
                               θ:θ∈R                      θ
\end{ReviewVerbatim}



\paragraph{Lean-elaborated paper signature}
\begin{ReviewVerbatim}
type:
{Voter : Type u_1} →
  [Fintype Voter] →
    {n dimension : ℕ} →
      (AppliedModelingLib.SocialChoice.Ranking.Candidate n →
          AppliedModelingLib.Alignment.Axioms.FeatureVector dimension) →
        AppliedModelingLib.Alignment.Axioms.RankingProfile Voter n → (ℝ → ℝ) → ℝ

value:
fun {Voter : Type u_1} [Fintype Voter] {n dimension : ℕ}
    (features :
      AppliedModelingLib.SocialChoice.Ranking.Candidate n → AppliedModelingLib.Alignment.Axioms.FeatureVector dimension)
    (profile : AppliedModelingLib.Alignment.Axioms.RankingProfile Voter n) (loss : ℝ → ℝ) =>
  sInf (Set.range (GeEtAl2024AlignmentAxioms.majorityLoss features profile loss))
\end{ReviewVerbatim}


\paragraph{Recorded source-to-declaration screening}
\textbf{Verdict:} \texttt{matches}
\\
{\raggedright
\textbf{Reason:} The declaration is `sInf` of the range of the expanded majority-\allowbreak{}loss objective over all linear reward parameters, exactly the unrestricted infimum in the routed majority-\allowbreak{}loss formulation.\allowbreak{}
\par}
\reviewmatch{Matches source input}
\noindent\textbf{Reviewer annotation}\par
\reviewerbox
\clearpage
\hypertarget{paper-prerequisite-8892d6fc4aee77f3}{}
\subsection*{\small\ttfamily\raggedright Ge\allowbreak{}Et\allowbreak{}Al2024\allowbreak{}Alignment\allowbreak{}Axioms.\allowbreak{}majority\allowbreak{}Loss}
\noindent\textbf{Paper-local declaration:}\par
{\footnotesize\ttfamily\raggedright Ge\allowbreak{}Et\allowbreak{}Al2024\allowbreak{}Alignment\allowbreak{}Axioms.\allowbreak{}majority\allowbreak{}Loss\par}
\textbf{Source locator:} {\footnotesize\raggedright cited publication, lines 370-\allowbreak{}386\par}
\paragraph{Verbatim paper-source connection}
\begin{ReviewVerbatim}
3.2   Majority-Based Loss Formulation

Despite the negative results for loss-function-based rules, we may hope for a remedy using slightly
different information. Specifically, we consider a similar loss-based function that rather than getting
penalized for disagreeing with each voter only gets penalized if it disagrees with a majority of voters.
That is, we choose θ minimizing
                                       X
                    Lmaj (θ; π, ell) =         I[wa≻b (π) > 1/2] · ell(rθ(b) − rθ(a) ).
                                       a̸=b∈C

Defining a parameter aggregation function based on this loss suffers from the same caveat as before,
that in some cases no optimal θ exists. Nevertheless, we can apply an analogous fix for a ranking
variant. We say that a linear rank aggregation rule f minimizes ell in the majority formulation if for
all σ = f (π),
                               inf σ Lmaj (θ; π, ell) = inf Lmaj (θ; π, ell).
                               θ:θ∈R                      θ
\end{ReviewVerbatim}



\paragraph{Lean-elaborated paper signature}
\begin{ReviewVerbatim}
type:
{Voter : Type u_1} →
  [Fintype Voter] →
    {n dimension : ℕ} →
      (AppliedModelingLib.SocialChoice.Ranking.Candidate n →
          AppliedModelingLib.Alignment.Axioms.FeatureVector dimension) →
        AppliedModelingLib.Alignment.Axioms.RankingProfile Voter n →
          (ℝ → ℝ) → AppliedModelingLib.Alignment.Axioms.LinearRewardParameter dimension → ℝ

value:
fun {Voter : Type u_1} [Fintype Voter] {n dimension : ℕ}
    (features :
      AppliedModelingLib.SocialChoice.Ranking.Candidate n → AppliedModelingLib.Alignment.Axioms.FeatureVector dimension)
    (profile : AppliedModelingLib.Alignment.Axioms.RankingProfile Voter n) (loss : ℝ → ℝ)
    (parameter : AppliedModelingLib.Alignment.Axioms.LinearRewardParameter dimension) =>
  ∑ pair : AppliedModelingLib.SocialChoice.Ranking.Candidate n × AppliedModelingLib.SocialChoice.Ranking.Candidate n,
    if
        Fintype.card Voter <
          2 *
            {voter : Voter |
                AppliedModelingLib.Alignment.Axioms.StrictlyPrefers (profile voter) pair.1 pair.2}.card then
      loss
        (AppliedModelingLib.Alignment.Axioms.linearReward parameter features pair.2 -
          AppliedModelingLib.Alignment.Axioms.linearReward parameter features pair.1)
    else 0
\end{ReviewVerbatim}


\paragraph{Recorded source-to-declaration screening}
\textbf{Verdict:} \texttt{matches}
\\
{\raggedright
\textbf{Reason:} The expanded objective sums a loss term exactly when a strict voter majority ranks the ordered pair first-\allowbreak{}over-\allowbreak{}second, at the reward difference second minus first.\allowbreak{} This is the routed formula for Lmaj;\allowbreak{} self-\allowbreak{}pairs contribute zero because strict self-\allowbreak{}preference is false.\allowbreak{}
\par}
\reviewmatch{Matches source input}
\noindent\textbf{Reviewer annotation}\par
\reviewerbox
\clearpage
\hypertarget{paper-prerequisite-465cdc3ee9c9c6a0}{}
\subsection*{\small\ttfamily\raggedright Ge\allowbreak{}Et\allowbreak{}Al2024\allowbreak{}Alignment\allowbreak{}Axioms.\allowbreak{}restricted\allowbreak{}Majority\allowbreak{}Loss\allowbreak{}Infimum}
\noindent\textbf{Paper-local declaration:}\par
{\footnotesize\ttfamily\raggedright Ge\allowbreak{}Et\allowbreak{}Al2024\allowbreak{}Alignment\allowbreak{}Axioms.\allowbreak{}restricted\allowbreak{}Majority\allowbreak{}Loss\allowbreak{}Infimum\par}
\textbf{Source locator:} {\footnotesize\raggedright cited publication, lines 370-\allowbreak{}386\par}
\paragraph{Verbatim paper-source connection}
\begin{ReviewVerbatim}
3.2   Majority-Based Loss Formulation

Despite the negative results for loss-function-based rules, we may hope for a remedy using slightly
different information. Specifically, we consider a similar loss-based function that rather than getting
penalized for disagreeing with each voter only gets penalized if it disagrees with a majority of voters.
That is, we choose θ minimizing
                                       X
                    Lmaj (θ; π, ell) =         I[wa≻b (π) > 1/2] · ell(rθ(b) − rθ(a) ).
                                       a̸=b∈C

Defining a parameter aggregation function based on this loss suffers from the same caveat as before,
that in some cases no optimal θ exists. Nevertheless, we can apply an analogous fix for a ranking
variant. We say that a linear rank aggregation rule f minimizes ell in the majority formulation if for
all σ = f (π),
                               inf σ Lmaj (θ; π, ell) = inf Lmaj (θ; π, ell).
                               θ:θ∈R                      θ
\end{ReviewVerbatim}



\paragraph{Lean-elaborated paper signature}
\begin{ReviewVerbatim}
type:
{Voter : Type u_1} →
  [Fintype Voter] →
    {n dimension : ℕ} →
      (AppliedModelingLib.SocialChoice.Ranking.Candidate n →
          AppliedModelingLib.Alignment.Axioms.FeatureVector dimension) →
        AppliedModelingLib.Alignment.Axioms.RankingProfile Voter n →
          (ℝ → ℝ) → AppliedModelingLib.SocialChoice.Ranking.Ranking n → ℝ

value:
fun {Voter : Type u_1} [Fintype Voter] {n dimension : ℕ}
    (features :
      AppliedModelingLib.SocialChoice.Ranking.Candidate n → AppliedModelingLib.Alignment.Axioms.FeatureVector dimension)
    (profile : AppliedModelingLib.Alignment.Axioms.RankingProfile Voter n) (loss : ℝ → ℝ)
    (ranking : AppliedModelingLib.SocialChoice.Ranking.Ranking n) =>
  sInf
    (GeEtAl2024AlignmentAxioms.majorityLoss features profile loss ''
      {parameter : AppliedModelingLib.Alignment.Axioms.LinearRewardParameter dimension |
        AppliedModelingLib.Alignment.Axioms.InducesRanking features parameter ranking})
\end{ReviewVerbatim}


\paragraph{Recorded source-to-declaration screening}
\textbf{Verdict:} \texttt{matches}
\\
{\raggedright
\textbf{Reason:} The declaration takes `sInf` of the majority-\allowbreak{}loss image of precisely the parameters that induce the supplied ranking.\allowbreak{} This is the source's infimum restricted to the parameter region R\textasciicircum{}sigma.\allowbreak{}
\par}
\reviewmatch{Matches source input}
\noindent\textbf{Reviewer annotation}\par
\reviewerbox
\clearpage
\hypertarget{paper-prerequisite-efbd7cf1dcd57ea6}{}
\subsection*{\small\ttfamily\raggedright Ge\allowbreak{}Et\allowbreak{}Al2024\allowbreak{}Alignment\allowbreak{}Axioms.\allowbreak{}standard\allowbreak{}Loss\allowbreak{}Formulation\allowbreak{}Spec}
\noindent\textbf{Paper-local declaration:}\par
{\footnotesize\ttfamily\raggedright Ge\allowbreak{}Et\allowbreak{}Al2024\allowbreak{}Alignment\allowbreak{}Axioms.\allowbreak{}standard\allowbreak{}Loss\allowbreak{}Formulation\allowbreak{}Spec\par}
\textbf{Source locator:} {\footnotesize\raggedright cited publication, lines 237-\allowbreak{}259\par}
\paragraph{Verbatim paper-source connection}
\begin{ReviewVerbatim}
3         Loss-Based Rules
3.1       Standard Loss Formulation

We begin our study of linear social choice by considering a quite broad yet natural class of rules
that capture how RLHF is currently being done. Their core idea is the following: when considering
parameter vector θ, for each voter i that ranks a pair of candidates a ≻i b, we should incur some
loss for giving b a higher reward than a. To formalize this, let ell : R → R be a loss function, which
we assume is nonnegative. We can then choose a parameter vector minimizing
                                         X
                           L(θ; π, ell) =       na≻b (π) · ell(rθ (b) − rθ (a)).
                                             a̸=b∈C

Note that the BTL model fits within this framework using ell(x) = ln(1 + ex ), i.e., binary cross-
entropy loss.5 One caveat to this approach, however, is that an optimal θ need not be well-defined:
it is possible that no minimum is attained. Fortunately, since we only care about rankings induced
by optimal parameter vectors, we can conveniently remedy this by saying the output is any ranking
that is induced by parameter vectors that are arbitrarily close to optimal. More formally, we say that
a linear rank aggregation rule f minimizes ell if for all σ = f (π),
                                       inf L(θ; π, ell) = inf L(θ; π, ell).                                       (1)
                                     θ:θ∈Rσ                    θ

Even if no minimum is attained, there is always a choice of feasible ranking σ such that Equation (1)
is satisfied.
\end{ReviewVerbatim}



\paragraph{Lean-elaborated paper signature}
\begin{ReviewVerbatim}
type:
{Voter : Type u_1} →
  [Fintype Voter] →
    {n dimension : ℕ} →
      (features :
          AppliedModelingLib.SocialChoice.Ranking.Candidate n →
            AppliedModelingLib.Alignment.Axioms.FeatureVector dimension) →
        (ℝ → ℝ) →
          AppliedModelingLib.Alignment.Axioms.LinearRankAggregationRule Voter n
              (AppliedModelingLib.Alignment.Axioms.LinearFeasibleRanking features) →
            Prop

value:
fun {Voter : Type u_1} [Fintype Voter] {n dimension : ℕ}
    (features :
      AppliedModelingLib.SocialChoice.Ranking.Candidate n → AppliedModelingLib.Alignment.Axioms.FeatureVector dimension)
    (loss : ℝ → ℝ)
    (rule :
      AppliedModelingLib.Alignment.Axioms.LinearRankAggregationRule Voter n
        (AppliedModelingLib.Alignment.Axioms.LinearFeasibleRanking features)) =>
  have objective :
    AppliedModelingLib.Alignment.Axioms.RankingProfile Voter n →
      AppliedModelingLib.Alignment.Axioms.LinearRewardParameter dimension → ℝ :=
    fun (profile : AppliedModelingLib.Alignment.Axioms.RankingProfile Voter n)
      (parameter : AppliedModelingLib.Alignment.Axioms.LinearRewardParameter dimension) =>
    ∑ voter : Voter,
      ∑ pair :
        AppliedModelingLib.SocialChoice.Ranking.Candidate n × AppliedModelingLib.SocialChoice.Ranking.Candidate n,
        if AppliedModelingLib.Alignment.Axioms.StrictlyPrefers (profile voter) pair.1 pair.2 then
          loss
            (AppliedModelingLib.Alignment.Axioms.linearReward parameter features pair.2 -
              AppliedModelingLib.Alignment.Axioms.linearReward parameter features pair.1)
        else 0;
  (∀ (input : ℝ), 0 ≤ loss input) ∧
    ∀ (profile : AppliedModelingLib.Alignment.Axioms.RankingProfile Voter n),
      AppliedModelingLib.Alignment.Axioms.FeasibleProfile
          (AppliedModelingLib.Alignment.Axioms.LinearFeasibleRanking features) profile →
        sInf
            (objective profile ''
              {parameter : AppliedModelingLib.Alignment.Axioms.LinearRewardParameter dimension |
                AppliedModelingLib.Alignment.Axioms.InducesRanking features parameter (rule.1 profile)}) =
          sInf (Set.range (objective profile))
\end{ReviewVerbatim}


\paragraph{Recorded source-to-declaration screening}
\textbf{Verdict:} \texttt{matches}
\\
{\raggedright
\textbf{Reason:} The expanded target requires a nonnegative loss, sums loss over every strict submitted comparison at rtheta(b)-\allowbreak{}rtheta(a), and equates the infimum restricted to parameters inducing the output ranking with the unrestricted infimum.\allowbreak{} These are the routed standard-\allowbreak{}loss formulation's stated components.\allowbreak{}
\par}
\reviewmatch{Matches source input}
\noindent\textbf{Reviewer annotation}\par
\reviewerbox

\clearpage
\hypertarget{source-claim-74b7c9c683220bb0}{}
\hypertarget{source-section-a2cb018f24443c81}{}
\section*{Main-text results}
\subsection*{Review row 1}
\textbf{Source locator:} {\footnotesize\raggedright cited publication, lines 204-\allowbreak{}211\par}
\paragraph{Verbatim source input}
\begin{ReviewVerbatim}
Definition 2.2 (Pairwise Majority Consistency (PMC)). A ranking σ is called a PMC ranking for
profile π if for all a, b ∈ C, a ≻σ b if and only if a majority of voters rank a ≻σi b, i.e., wa≻b > 1/2.
A linear rank aggregation rule satisfies PMC if, when a PMC ranking σ exists for the input profile
π and σ is feasible, then f (π) = σ.

Note that a PMC ranking for each π need not exist, but when one does, it is unique. The words “σ
is feasible” allude to the possibility that no non-degenerate parameter vector θ induces the unique
PMC ranking. Indeed, we have such an example; see Appendix B for details.

[Next verbatim source excerpt]

2     The Linear Social Choice Model
Let C be a set of m distinct prompt/responses, referred to as candidates, and let V = {1, . . . , n} be
a set of n human participants, known as voters. We denote by Rd the d-dimensional real space in
which both candidate feature vectors and chosen parameter vectors lie.
Each candidate c ∈ C is associated with a distinct feature vector xc ∈ Rd . A parameter vector
θ ∈ Rd induces a linear reward function rθ : C 7→ R defined by taking the dot product with feature
vectors rθ (c) = ⟨θ, xc ⟩. We will primarily be interested in how these parameterized functions rank
the candidates by reward. Let Ra≻b = {θ | rθ (a) ≥ rθ (b)} be the region where the reward of a is
at least as large as that of b. Note that Ra≻b and Rb≻a split Rd into two half spaces, separated by
the hyperplane orthogonal to xa − xb . Parameter vectors θ on the hyperplane have rθ (a) = rθ (b),
while rankings in the interior of either half-space strictly rank one over the other.
For a ranking σ over the candidates, we say that θ induces σ, denoted θ ▷ σ, if a ≻σ b implies
rθ (a) ≥ rθ (b). Let Rσ = {θ | θ ▷ σ} be the set of vectors θ thatTinduce it. Note that this can
be written as the intersection of corresponding half spaces Rσ = a,b:a≻σ b Ra≻b . Further, the


                                                    3

[form-feed in source extraction]
collection of {Rσ } essentially form a partition of Rd , covering the space and intersecting only at
their boundaries.
We call a θ non-degenerate if it is fully on one side of each of the separating hyperplanes, i.e.,
rθ (a) ̸= rθ (b) for all a, b ∈ C. Non-degenerate parameter vectors lie in the interior of some Rσ ,
and thus induce exactly one ranking. We call σ feasible if Rσ has a nonempty interior, i.e., is induced
by some nondegenerate θ.3
Each voter i ∈ V submits a ranking over the candidate σi . We assume that the feature space is
rich enough that voter preferences can be captured via non-degenerate parameter vectors. In other
words, we assume that each σi is feasible. We refer to the vector of voter rankings π = (σi )i∈V as
a profile. Further, for two candidates a, b, we write na≻b (π) := |{i ∈ V | a ≻σi b}| for the number
of voters that prefer a to b, and wa≻b (π) = na≻b (π)/n for the proportion of such voters. When the
profile π is clear from context, we may shorten these to na≻b and wa≻b , respectively.
We define a parameter aggregation rule as a function that takes as input a profile π and outputs a
parameter vector θ∗ . Our goal is to design parameter aggregation rules such that rθ∗ satisfies desir-
able properties with respect to the voter preferences. However, as the properties we care about will
only be with respect to how rθ∗ ranks the candidates, it will be more convenient to work with what
we call linear rank aggregation rules that take as input a profile π and output a feasible ranking σ.
There is a natural way to interpret a parameter aggregation rule as a linear rank aggregation rule,
namely, output any feasible ranking induced by θ∗ . The exact properties of the parameter aggrega-
tion rule could in principle be sensitive to the tie-breaking of non-degenerate outputs, however, all
of our results will be robust to such tie-breaking.4
We pay special attention to a prominent family of rules from social choice theory referred to as C1
rules [14], whose outputs depend only on majority relationships, i.e., they only need to know for
each pair of candidates (a, b) whether the majority prefers a or b.
In our study, we examine several axioms borrowed from social choice theory to evaluate the reason-
ableness (fairness) of our aggregation mechanisms. These axioms include:
Definition 2.1 (Pareto Optimality). A linear rank aggregation rule f satisfies Pareto optimality if,
whenever every voter prefers candidate a over candidate b on π, i.e., wa≻b (π) = 1, then candidate
a is ranked higher than candidate b in the output ranking, i.e., a ≻f (π) b.
Definition 2.2 (Pairwise Majority Consistency (PMC)). A ranking σ is called a PMC ranking for
profile π if for all a, b ∈ C, a ≻σ b if and only if a majority of voters rank a ≻σi b, i.e., wa≻b > 1/2.
A linear rank aggregation rule satisfies PMC if, when a PMC ranking σ exists for the input profile
π and σ is feasible, then f (π) = σ.

Note that a PMC ranking for each π need not exist, but when one does, it is unique. The words “σ
is feasible” allude to the possibility that no non-degenerate parameter vector θ induces the unique
PMC ranking. Indeed, we have such an example; see Appendix B for details.
\end{ReviewVerbatim}



\paragraph{Expanded PaperInterface specification}
\begin{ReviewVerbatim}
(∀ (Voter : Type) [inst : Fintype Voter] (n : ℕ) (feasible : AppliedModelingLib.SocialChoice.Ranking.Ranking n → Prop)
    (rule : AppliedModelingLib.Alignment.Axioms.LinearRankAggregationRule Voter n feasible),
    (∀ (profile : AppliedModelingLib.Alignment.Axioms.RankingProfile Voter n)
        (majorityRanking : AppliedModelingLib.SocialChoice.Ranking.Ranking n),
        AppliedModelingLib.Alignment.Axioms.FeasibleProfile feasible profile →
          feasible majorityRanking →
            (∀ (first second : AppliedModelingLib.SocialChoice.Ranking.Candidate n),
                AppliedModelingLib.Alignment.Axioms.StrictlyPrefers majorityRanking first second ↔
                  Fintype.card Voter <
                    2 *
                      {voter : Voter |
                          AppliedModelingLib.Alignment.Axioms.StrictlyPrefers (profile voter) first second}.card) →
              rule.1 profile = majorityRanking) ↔
      ∀ (profile : AppliedModelingLib.Alignment.Axioms.RankingProfile Voter n)
        (majorityRanking : AppliedModelingLib.SocialChoice.Ranking.Ranking n),
        AppliedModelingLib.Alignment.Axioms.FeasibleProfile feasible profile →
          feasible majorityRanking →
            (∀ (first second : AppliedModelingLib.SocialChoice.Ranking.Candidate n),
                AppliedModelingLib.Alignment.Axioms.StrictlyPrefers majorityRanking first second ↔
                  Fintype.card Voter <
                    2 *
                      {voter : Voter |
                          AppliedModelingLib.Alignment.Axioms.StrictlyPrefers (profile voter) first second}.card) →
              rule.1 profile = majorityRanking) ∧
  (¬∃ (ranking : AppliedModelingLib.SocialChoice.Ranking.Ranking 7),
        ∀ (first second : AppliedModelingLib.SocialChoice.Ranking.Candidate 7),
          AppliedModelingLib.Alignment.Axioms.StrictlyPrefers ranking first second ↔
            Fintype.card (Fin 5) <
              2 *
                {voter : Fin 5 |
                    AppliedModelingLib.Alignment.Axioms.StrictlyPrefers
                      ((fun (x : Fin 5) =>
                          Fin.casesOn x fun (val : ℕ) (isLt : val < 5) =>
                            Nat.casesOn (motive := fun (x : ℕ) =>
                              (isLt : x < 5) →
                                (fun (x : Fin 5) => AppliedModelingLib.SocialChoice.Ranking.Ranking 7) ⟨x, isLt⟩)
                              val
                              (fun (isLt : Nat.zero < 5) =>
                                (fun (isLt : 0 < 5) =>
                                    ((((Equiv.trans (Equiv.swap 1 8) (Equiv.swap 1 4)).trans (Equiv.swap 1 3)).trans
                                              (Equiv.swap 1 2)).trans
                                          (Equiv.swap 5 7)).trans
                                      {
                                        toFun :=
                                          fun (candidate : AppliedModelingLib.SocialChoice.Ranking.Candidate 7) =>
                                          if hplus : ↑candidate < 4 then ⟨↑((Equiv.swap 0 1) ⟨↑candidate, hplus⟩), ⋯⟩
                                          else
                                            if hminus : ↑candidate < 8 then
                                              ⟨↑((Equiv.swap 0 1) ⟨↑candidate - 4, ⋯⟩) + 4, ⋯⟩
                                            else ⟨8, GeEtAl2024AlignmentAxioms.c1Star._proof_2⟩,
                                        invFun :=
                                          fun (candidate : AppliedModelingLib.SocialChoice.Ranking.Candidate 7) =>
                                          if hplus : ↑candidate < 4 then
                                            ⟨↑((Equiv.symm (Equiv.swap 0 1)) ⟨↑candidate, hplus⟩), ⋯⟩
                                          else
                                            if hminus : ↑candidate < 8 then
                                              ⟨↑((Equiv.symm (Equiv.swap 0 1)) ⟨↑candidate - 4, ⋯⟩) + 4, ⋯⟩
                                            else ⟨8, GeEtAl2024AlignmentAxioms.c1Star._proof_2⟩,
                                        left_inv := ⋯, right_inv := ⋯ })
                                  isLt)
                              (fun (n : ℕ) (isLt : n.succ < 5) =>
                                Nat.casesOn (motive := fun (x : ℕ) =>
                                  (isLt : x.succ < 5) →
                                    (fun (x : Fin 5) => AppliedModelingLib.SocialChoice.Ranking.Ranking 7)
                                      ⟨x.succ, isLt⟩)
                                  n
                                  (fun (isLt : Nat.zero.succ < 5) =>
                                    (fun (isLt : 1 < 5) =>
                                        ((((Equiv.trans (Equiv.swap 1 8) (Equiv.swap 1 4)).trans (Equiv.swap 1 3)).trans
                                                  (Equiv.swap 1 2)).trans
                                              (Equiv.swap 5 7)).trans
                                          {
                                            toFun :=
                                              fun (candidate : AppliedModelingLib.SocialChoice.Ranking.Candidate 7) =>
                                              if hplus : ↑candidate < 4 then
                                                ⟨↑((Equiv.trans (Equiv.swap 0 2) (Equiv.swap 0 1)) ⟨↑candidate, hplus⟩),
                                                  ⋯⟩
                                              else
                                                if hminus : ↑candidate < 8 then
                                                  ⟨↑((Equiv.trans (Equiv.swap 0 2) (Equiv.swap 0 1))
                                                          ⟨↑candidate - 4, ⋯⟩) +
                                                      4,
                                                    ⋯⟩
                                                else ⟨8, GeEtAl2024AlignmentAxioms.c1Star._proof_2⟩,
                                            invFun :=
                                              fun (candidate : AppliedModelingLib.SocialChoice.Ranking.Candidate 7) =>
                                              if hplus : ↑candidate < 4 then
                                                ⟨↑((Equiv.trans (Equiv.swap 0 2) (Equiv.swap 0 1)).symm
                                                      ⟨↑candidate, hplus⟩),
                                                  ⋯⟩
                                              else
                                                if hminus : ↑candidate < 8 then
                                                  ⟨↑((Equiv.trans (Equiv.swap 0 2) (Equiv.swap 0 1)).symm
                                                          ⟨↑candidate - 4, ⋯⟩) +
                                                      4,
                                                    ⋯⟩
                                                else ⟨8, GeEtAl2024AlignmentAxioms.c1Star._proof_2⟩,
                                            left_inv := ⋯, right_inv := ⋯ })
                                      isLt)
                                  (fun (n : ℕ) (isLt : n.succ.succ < 5) =>
                                    Nat.casesOn (motive := fun (x : ℕ) =>
                                      (isLt : x.succ.succ < 5) →
                                        (fun (x : Fin 5) => AppliedModelingLib.SocialChoice.Ranking.Ranking 7)
                                          ⟨x.succ.succ, isLt⟩)
                                      n
                                      (fun (isLt : Nat.zero.succ.succ < 5) =>
                                        (fun (isLt : 2 < 5) =>
                                            ((((Equiv.trans (Equiv.swap 1 8) (Equiv.swap 1 4)).trans
                                                          (Equiv.swap 1 3)).trans
                                                      (Equiv.swap 1 2)).trans
                                                  (Equiv.swap 5 7)).trans
                                              {
                                                toFun :=
                                                  fun
                                                    (candidate : AppliedModelingLib.SocialChoice.Ranking.Candidate 7) =>
                                                  if hplus : ↑candidate < 4 then
                                                    ⟨↑((Equiv.trans (Equiv.swap 0 2) (Equiv.swap 0 3))
                                                          ⟨↑candidate, hplus⟩),
                                                      ⋯⟩
                                                  else
                                                    if hminus : ↑candidate < 8 then
                                                      ⟨↑((Equiv.trans (Equiv.swap 0 2) (Equiv.swap 0 3))
                                                              ⟨↑candidate - 4, ⋯⟩) +
                                                          4,
                                                        ⋯⟩
                                                    else ⟨8, GeEtAl2024AlignmentAxioms.c1Star._proof_2⟩,
                                                invFun :=
                                                  fun
                                                    (candidate : AppliedModelingLib.SocialChoice.Ranking.Candidate 7) =>
                                                  if hplus : ↑candidate < 4 then
                                                    ⟨↑((Equiv.trans (Equiv.swap 0 2) (Equiv.swap 0 3)).symm
                                                          ⟨↑candidate, hplus⟩),
                                                      ⋯⟩
                                                  else
                                                    if hminus : ↑candidate < 8 then
                                                      ⟨↑((Equiv.trans (Equiv.swap 0 2) (Equiv.swap 0 3)).symm
                                                              ⟨↑candidate - 4, ⋯⟩) +
                                                          4,
                                                        ⋯⟩
                                                    else ⟨8, GeEtAl2024AlignmentAxioms.c1Star._proof_2⟩,
                                                left_inv := ⋯, right_inv := ⋯ })
                                          isLt)
                                      (fun (n : ℕ) (isLt : n.succ.succ.succ < 5) =>
                                        Nat.casesOn (motive := fun (x : ℕ) =>
                                          (isLt : x.succ.succ.succ < 5) →
                                            (fun (x : Fin 5) => AppliedModelingLib.SocialChoice.Ranking.Ranking 7)
                                              ⟨x.succ.succ.succ, isLt⟩)
                                          n
                                          (fun (isLt : Nat.zero.succ.succ.succ < 5) =>
                                            (fun (isLt : 3 < 5) =>
                                                ((((Equiv.trans (Equiv.swap 1 8) (Equiv.swap 1 4)).trans
                                                              (Equiv.swap 1 3)).trans
                                                          (Equiv.swap 1 2)).trans
                                                      (Equiv.swap 5 7)).trans
                                                  {
                                                    toFun :=
                                                      fun
                                                        (candidate :
                                                          AppliedModelingLib.SocialChoice.Ranking.Candidate 7) =>
                                                      if hplus : ↑candidate < 4 then
                                                        ⟨↑(((Equiv.trans (Equiv.swap 0 3) (Equiv.swap 0 2)).trans
                                                                (Equiv.swap 0 1))
                                                              ⟨↑candidate, hplus⟩),
                                                          ⋯⟩
                                                      else
                                                        if hminus : ↑candidate < 8 then
                                                          ⟨↑(((Equiv.trans (Equiv.swap 0 3) (Equiv.swap 0 2)).trans
                                                                    (Equiv.swap 0 1))
                                                                  ⟨↑candidate - 4, ⋯⟩) +
                                                              4,
                                                            ⋯⟩
                                                        else ⟨8, GeEtAl2024AlignmentAxioms.c1Star._proof_2⟩,
                                                    invFun :=
                                                      fun
                                                        (candidate :
                                                          AppliedModelingLib.SocialChoice.Ranking.Candidate 7) =>
                                                      if hplus : ↑candidate < 4 then
                                                        ⟨↑(((Equiv.trans (Equiv.swap 0 3) (Equiv.swap 0 2)).trans
                                                                  (Equiv.swap 0 1)).symm
                                                              ⟨↑candidate, hplus⟩),
                                                          ⋯⟩
                                                      else
                                                        if hminus : ↑candidate < 8 then
                                                          ⟨↑(((Equiv.trans (Equiv.swap 0 3) (Equiv.swap 0 2)).trans
                                                                      (Equiv.swap 0 1)).symm
                                                                  ⟨↑candidate - 4, ⋯⟩) +
                                                              4,
                                                            ⋯⟩
                                                        else ⟨8, GeEtAl2024AlignmentAxioms.c1Star._proof_2⟩,
                                                    left_inv := ⋯, right_inv := ⋯ })
                                              isLt)
                                          (fun (n : ℕ) (isLt : n.succ.succ.succ.succ < 5) =>
                                            Nat.casesOn (motive := fun (x : ℕ) =>
                                              (isLt : x.succ.succ.succ.succ < 5) →
                                                (fun (x : Fin 5) => AppliedModelingLib.SocialChoice.Ranking.Ranking 7)
                                                  ⟨x.succ.succ.succ.succ, isLt⟩)
                                              n
                                              (fun (isLt : Nat.zero.succ.succ.succ.succ < 5) =>
                                                (fun (isLt : 4 < 5) =>
                                                    ((((Equiv.trans (Equiv.swap 1 8) (Equiv.swap 1 4)).trans
                                                                  (Equiv.swap 1 3)).trans
                                                              (Equiv.swap 1 2)).trans
                                                          (Equiv.swap 5 7)).trans
                                                      {
                                                        toFun :=
                                                          fun
                                                            (candidate :
                                                              AppliedModelingLib.SocialChoice.Ranking.Candidate 7) =>
                                                          if hplus : ↑candidate < 4 then
                                                            ⟨↑(((Equiv.trans (Equiv.swap 0 3) (Equiv.swap 0 2)).trans
                                                                    (Equiv.swap 0 1))
                                                                  ⟨↑candidate, hplus⟩),
                                                              ⋯⟩
                                                          else
                                                            if hminus : ↑candidate < 8 then
                                                              ⟨↑(((Equiv.trans (Equiv.swap 0 3) (Equiv.swap 0 2)).trans
                                                                        (Equiv.swap 0 1))
                                                                      ⟨↑candidate - 4, ⋯⟩) +
                                                                  4,
                                                                ⋯⟩
                                                            else ⟨8, GeEtAl2024AlignmentAxioms.c1Star._proof_2⟩,
                                                        invFun :=
                                                          fun
                                                            (candidate :
                                                              AppliedModelingLib.SocialChoice.Ranking.Candidate 7) =>
                                                          if hplus : ↑candidate < 4 then
                                                            ⟨↑(((Equiv.trans (Equiv.swap 0 3) (Equiv.swap 0 2)).trans
                                                                      (Equiv.swap 0 1)).symm
                                                                  ⟨↑candidate, hplus⟩),
                                                              ⋯⟩
                                                          else
                                                            if hminus : ↑candidate < 8 then
                                                              ⟨↑(((Equiv.trans (Equiv.swap 0 3) (Equiv.swap 0 2)).trans
                                                                          (Equiv.swap 0 1)).symm
                                                                      ⟨↑candidate - 4, ⋯⟩) +
                                                                  4,
                                                                ⋯⟩
                                                            else ⟨8, GeEtAl2024AlignmentAxioms.c1Star._proof_2⟩,
                                                        left_inv := ⋯, right_inv := ⋯ })
                                                  isLt)
                                              (fun (n : ℕ) (isLt : n.succ.succ.succ.succ.succ < 5) => ⋯.elim) isLt)
                                          isLt)
                                      isLt)
                                  isLt)
                              isLt)
                        voter)
                      first second}.card) ∧
    ∀ (Voter : Type) [inst : Fintype Voter] (n : ℕ)
      (profile : AppliedModelingLib.Alignment.Axioms.RankingProfile Voter n)
      (first second : AppliedModelingLib.SocialChoice.Ranking.Ranking n),
      (∀ (first_1 second : AppliedModelingLib.SocialChoice.Ranking.Candidate n),
          AppliedModelingLib.Alignment.Axioms.StrictlyPrefers first first_1 second ↔
            Fintype.card Voter <
              2 *
                {voter : Voter |
                    AppliedModelingLib.Alignment.Axioms.StrictlyPrefers (profile voter) first_1 second}.card) →
        (∀ (first second_1 : AppliedModelingLib.SocialChoice.Ranking.Candidate n),
            AppliedModelingLib.Alignment.Axioms.StrictlyPrefers second first second_1 ↔
              Fintype.card Voter <
                2 *
                  {voter : Voter |
                      AppliedModelingLib.Alignment.Axioms.StrictlyPrefers (profile voter) first second_1}.card) →
          first = second
\end{ReviewVerbatim}



\paragraph{Lean proof endpoint (built separately)}\begin{ReviewVerbatim}
GeEtAl2024AlignmentAxioms.definition2_2_pairwiseMajorityConsistencyAndConsequences
\end{ReviewVerbatim}

\paragraph{Recorded source-to-Spec screening}
\textbf{Verdict:} \texttt{matches}
\\
{\raggedright
\textbf{Reason:} The expanded target matches Definition 2.\allowbreak{}2 and the attached observations clause by clause:\allowbreak{} PMC is exactly the condition that, on a feasible input profile, any feasible ranking whose strict comparisons agree iff the finite voter count is a strict majority must be the rule output;\allowbreak{} the displayed five-\allowbreak{}voter, nine-\allowbreak{}candidate cyclic profile supplies a concrete witness for the stated possibility that no PMC ranking exists;\allowbreak{} and the final universally quantified clause states uniqueness whenever two PMC rankings do exist.\allowbreak{} The displayed Ranking, Candidate, StrictlyPrefers, FeasibleProfile, and LinearRankAggregationRule meanings preserve the source domains and add no substantive premise or conclusion.\allowbreak{}
\par}
\reviewmatch{Matches source input}
\noindent\textbf{Reviewer annotation}\par
\reviewerbox
\clearpage
\hypertarget{source-claim-eb42e4a359411adb}{}
\section*{Review row 2}
\textbf{Source locator:} {\footnotesize\raggedright cited publication, lines 267-\allowbreak{}269\par}
\paragraph{Verbatim source input}
\begin{ReviewVerbatim}
Theorem 3.1. If a linear rank aggregation rule f optimizes a loss function ell that satisfies
inf x ell(x) < ell(0) and is either nondecreasing and weakly convex, or strictly convex (and possibly
nonmonotone), then f fails PMC and PO.

[Next verbatim source excerpt]


[form-feed in source extraction]
3         Loss-Based Rules
3.1       Standard Loss Formulation

We begin our study of linear social choice by considering a quite broad yet natural class of rules
that capture how RLHF is currently being done. Their core idea is the following: when considering
parameter vector θ, for each voter i that ranks a pair of candidates a ≻i b, we should incur some
loss for giving b a higher reward than a. To formalize this, let ell : R → R be a loss function, which
we assume is nonnegative. We can then choose a parameter vector minimizing
                                         X
                           L(θ; π, ell) =       na≻b (π) · ell(rθ (b) − rθ (a)).
                                             a̸=b∈C

Note that the BTL model fits within this framework using ell(x) = ln(1 + ex ), i.e., binary cross-
entropy loss.5 One caveat to this approach, however, is that an optimal θ need not be well-defined:
it is possible that no minimum is attained. Fortunately, since we only care about rankings induced
by optimal parameter vectors, we can conveniently remedy this by saying the output is any ranking
that is induced by parameter vectors that are arbitrarily close to optimal. More formally, we say that
a linear rank aggregation rule f minimizes ell if for all σ = f (π),
                                       inf L(θ; π, ell) = inf L(θ; π, ell).                                       (1)
                                     θ:θ∈Rσ                    θ

Even if no minimum is attained, there is always a choice of feasible ranking σ such that Equation (1)
is satisfied.
With this definition in hand, we proceed to our first main result, which spells rather bad news for
this class of rules: Any loss-based aggregation rule using a nondecreasing and convex loss function
(of which BTL is one, and hinge loss is another) will fail our two core axioms, PMC and PO. This
paints a negative picture for current RLHF methods with respect to their social choice guarantees.
Note that we will exclude the discussion of loss functions with a global minimum at zero, like ReLU,
because the loss minimizer will be zero, making all rankings vacuous consequently. And we have
focused on convex loss functions due to their practical optimization ease.
Theorem 3.1. If a linear rank aggregation rule f optimizes a loss function ell that satisfies
inf x ell(x) < ell(0) and is either nondecreasing and weakly convex, or strictly convex (and possibly
nonmonotone), then f fails PMC and PO.

[Next verbatim source excerpt]

Proof. Fix a loss function ell satisfying the theorem conditions. Note that since ell is convex, we may
also assume it is continuous [24, Corollary 10.1.1]. Furthermore, since inf x ell(x) < ell(0), we know
that there exists x ̸= 0 such that ell(x) < ell(0). The case where x > 0 is relatively simple (as
such loss functions lead to unnatural behavior), and we handle it at the end of the proof. For now,
we assume that there exists x < 0 such that ell(x) < ell(0). Note that this also implies that for all
y ≥ 0, ell is lower bounded by the affine linear function connecting (x, ell(x)) and (0, ell(0)), and thus,
limx→∞ ell(x) = ∞.
We begin with a small instance of just three candidates C core = {a, b, c} to gain some traction on
how ell behaves. We will later extend this instance with additional candidates to demonstrate a profile
where PO and PMC fail. The candidates will have feature vectors xa := (2, 1), xb := (1, 1), and
xc := (0, 0), respectively. Furthermore, a p-fraction of voters (for p to be chosen later) will rank
a ≻ b ≻ c, while the remaining (1 − p)-fraction will have inverted preferences, ranking c ≻ b ≻ a.6
Let                                              X
                              Lcore (θ) :=                  wx≻y ell(rθ (y) − rθ (x))
                                              x̸=y∈C core

      5
     Typically, BTL is presented as choosing θ maximizing the likelihood of seeing the pairwise com-
                                                                erθ (a)
parisons we observed, assuming that Pr[a ≻ b] = erθ (a)           +erθ (b)
                                                                           . That is, we choose θ maximizing
       (                ) na≻b
            erθ (a)
Q
  a̸=b erθ (a) +erθ (b)         . By taking the log and swapping the sign, we see that this is equivalent to mini-
                    (                     )
        P                   rθ (b)−rθ (a)
mizing a̸=b log 1 + e                      .
    6
      It so happens that these rankings are feasible, but for now we will not worry about this as loss function-
based rules still make sense regardless of whether the inputs are feasible. For the final example, we will ensure
that the rankings are feasible.


                                                        5

[form-feed in source extraction]
with wx≻y ∈ {1 − p, p} be the loss function on this instance (scaling nx≻y down to wx≻y leads to
an equivalent formulation). Let g(x) = p · ell(−x) + (1 − p)ell(x). Note that we can rewrite Lcom as
                    Lcore (θ) = g(rθ (a) − rθ (b)) + g(rθ (a) − rθ (c)) + g(rθ (b) − rθ (c)).
Note that rθ (c) = 0 for all θ, so we can simplify this to
                              Lcore (θ) = g(rθ (a) − rθ (b)) + g(rθ (a)) + g(rθ (b)).
We will consider an unconstrained version of this problem where we are free to choose rewards
ra , rb ∈ R arbitrarily, and later show by which vectors θ these optimal values can be induced. That
is, we will first find ra , rb ∈ R minimizing
                                Lunconstr (ra , rb ) := g(ra − rb ) + g(ra ) + g(rb ).
Let OP T core = {θ | Lcore (θ) = inf θ′ Lcore (θ′ )} and OP T unconstr {(ra , rb ) | Lunconstr (ra , rb ) =
inf ra′ ,rb′ Lunconstr (ra′ , rb′ )} be the set of minimizers for these two loss functions. In Appendix A, we
establish the following results about these optimal sets.
Lemma 3.2. There exists a rational p ∈ (1/2, 1] and values A1 < A2 with A2 > 0 such that
OP T unconstr is nonempty and for all (ra , rb ) ∈ OP T unconstr , ra > A2 and rb ≤ A1 .
Lemma 3.3. Suppose Lemma 3.2 holds for values p, A1 and A2 , then, for this same choice of p,
OP T core is nonempty and there exist A3 and A4 with A3 > 0 such that for all (θ1 , θ2 ) ∈ OP T core ,
θ1 > A3 and θ2 < A4 .

We will now explicitly construct a family of instances with candidate feature vectors parameterized
by a value ε ∈ R such that for sufficiently small ε > 0, the output of f fails the two axioms. Fix
p, A3 and A4 from Lemma 3.3, and choose δ with 0 < δ < 1 such that δA4 − A3 < 0 (δ < A3 /A4
works if A4 > 0, and otherwise, any 0 < δ < 1 will do).
Each instance will have six candidates, which we will think of as two groups of three, C = C core ∪
C copies . The first group C core = {a, b, c} will be the same as the three-candidate instance from
above, while the second group C copies = {a′ , b′ , c′ } will be new. The candidates a, b, c will still be
located at xa := (2, 1), xb := (1, 1), and xc := (0, 0), respectively. The candidates a′ , b′ , c′ will
be located near their undecorated counterparts at xa′ := xa + (−ε, 0), xb′ := xb + (−ε, 0) and
xc′ := xc + (−ε, δ · ε).
Next, we describe the voter preferences. A p-fraction of voters will have the ranking a ≻ a′ ≻ b ≻
b′ ≻ c′ ≻ c, and the remaining (1−p)-fraction of voters will have ranking c′ ≻ c ≻ b′ ≻ b ≻ a′ ≻ a.
As long as 0 < ε < 1 (which will be the case for our final chosen ε), these are both feasible rankings.
The former is induced by the nondegenerate feature vector (1, 1)7 and the latter by (−1, 0).8
For each ε ∈ R, let                                 X
                                      Lε (θ) =             wx≻y ell(rθ (y) − rθ (x))
                                                  x̸=y∈C

with wx≻y ∈ {0, 1 − p, p, 1} be the loss function we are optimizing using candidate locations
parameterized by ε.
We will show that for sufficiently small ε > 0, inf θ∈Rc′ ≻c Lε (θ) > inf θ Lε (θ). This means that
f must output a ranking with c ≻ c′ . Observe that this is a PO violation because all voters agree
that c′ ≻ c. Furthermore, this is a PMC violation because a majority of voters have the ranking
a ≻ a′ ≻ b ≻ b′ ≻ c′ ≻ c, yet this is not the output.
Let OP T (ε) be the set of vectors optimizing Lε . The rest of the proof will follow from the following
two lemmas, whose proofs are in Appendix A.

Lemma 3.4. OP T (0) ⊆ Rc′ ≻c .

Lemma 3.5. Suppose OP T (0) ⊆ Rc′ ≻c , then, for sufficiently small ε > 0, inf θ:θ∈Rc′ ≻c Lε (θ) >
inf θ Lε (θ).
   7
       This induces rewards 3, 3 − ε, 2, 2 − ε, (1 − δ) · ε and 0 for a, a′ , b, b′ , c′ , and c, respectively.
   8
       This induces rewards ε, 0, −1 + ε, −1, −2 + ε and −2. for c′ , c, b′ , b, a′ , and a, respectively.


                                                              6

[form-feed in source extraction]
Finally, we handle the case that exists x > 0 such that ell(x) < ell(0). Note that by convexity, this
implies that for all y < 0, ell(y) > ell(0) > ell(x), so inf y≤0 ell(y) < inf y ell(y). Now, consider an
instance with two candidates {a, b} located at xa = (1, 0) and xb = (0, 1), and a single voter
ranking a ≻ b (feasible via the parameter vector (1, 0)). It is possible to achieve a loss of ell(x), e.g.,
by outputting the parameter vector (0, x). On the other hand, any θ inducing the ranking a ≻ b will
be lower bounded by ell(0) > ell(x) from above. Hence, f must output b ≻ a, which is both a PO and
PMC violation.
\end{ReviewVerbatim}



\paragraph{Expanded PaperInterface specification}
\begin{ReviewVerbatim}
∀ (loss : ℝ → ℝ),
  (∀ (input : ℝ), 0 ≤ loss input) →
    Monotone loss ∧ ConvexOn ℝ Set.univ loss ∨ StrictConvexOn ℝ Set.univ loss →
      sInf (Set.range loss) < loss 0 →
        (∃ (positiveInput : ℝ),
            0 < positiveInput ∧
              loss positiveInput < loss 0 ∧
                ∀
                  (rule :
                    AppliedModelingLib.Alignment.Axioms.LinearRankAggregationRule Unit 0
                      (AppliedModelingLib.Alignment.Axioms.LinearFeasibleRanking
                        fun (candidate : AppliedModelingLib.SocialChoice.Ranking.Candidate 0) (coordinate : Fin 2) =>
                        if candidate = 0 then if coordinate = 0 then 1 else 0 else if coordinate = 0 then 0 else 1)),
                  (∀ (profile : AppliedModelingLib.Alignment.Axioms.RankingProfile Unit 0),
                      AppliedModelingLib.Alignment.Axioms.FeasibleProfile
                          (AppliedModelingLib.Alignment.Axioms.LinearFeasibleRanking
                            fun (candidate : AppliedModelingLib.SocialChoice.Ranking.Candidate 0)
                              (coordinate : Fin 2) =>
                            if candidate = 0 then if coordinate = 0 then 1 else 0 else if coordinate = 0 then 0 else 1)
                          profile →
                        sInf
                            ((fun (parameter : AppliedModelingLib.Alignment.Axioms.LinearRewardParameter 2) =>
                                ∑ voter : Unit,
                                  ∑ pair :
                                    AppliedModelingLib.SocialChoice.Ranking.Candidate 0 ×
                                      AppliedModelingLib.SocialChoice.Ranking.Candidate 0,
                                    if
                                        AppliedModelingLib.Alignment.Axioms.StrictlyPrefers (profile voter) pair.1
                                          pair.2 then
                                      loss
                                        (AppliedModelingLib.Alignment.Axioms.linearReward parameter
                                            (fun (candidate : AppliedModelingLib.SocialChoice.Ranking.Candidate 0)
                                                (coordinate : Fin 2) =>
                                              if candidate = 0 then if coordinate = 0 then 1 else 0
                                              else if coordinate = 0 then 0 else 1)
                                            pair.2 -
                                          AppliedModelingLib.Alignment.Axioms.linearReward parameter
                                            (fun (candidate : AppliedModelingLib.SocialChoice.Ranking.Candidate 0)
                                                (coordinate : Fin 2) =>
                                              if candidate = 0 then if coordinate = 0 then 1 else 0
                                              else if coordinate = 0 then 0 else 1)
                                            pair.1)
                                    else 0) ''
                              {parameter : AppliedModelingLib.Alignment.Axioms.LinearRewardParameter 2 |
                                AppliedModelingLib.Alignment.Axioms.InducesRanking
                                  (fun (candidate : AppliedModelingLib.SocialChoice.Ranking.Candidate 0)
                                      (coordinate : Fin 2) =>
                                    if candidate = 0 then if coordinate = 0 then 1 else 0
                                    else if coordinate = 0 then 0 else 1)
                                  parameter (rule.1 profile)}) =
                          sInf
                            (Set.range fun (parameter : AppliedModelingLib.Alignment.Axioms.LinearRewardParameter 2) =>
                              ∑ voter : Unit,
                                ∑ pair :
                                  AppliedModelingLib.SocialChoice.Ranking.Candidate 0 ×
                                    AppliedModelingLib.SocialChoice.Ranking.Candidate 0,
                                  if
                                      AppliedModelingLib.Alignment.Axioms.StrictlyPrefers (profile voter) pair.1
                                        pair.2 then
                                    loss
                                      (AppliedModelingLib.Alignment.Axioms.linearReward parameter
                                          (fun (candidate : AppliedModelingLib.SocialChoice.Ranking.Candidate 0)
                                              (coordinate : Fin 2) =>
                                            if candidate = 0 then if coordinate = 0 then 1 else 0
                                            else if coordinate = 0 then 0 else 1)
                                          pair.2 -
                                        AppliedModelingLib.Alignment.Axioms.linearReward parameter
                                          (fun (candidate : AppliedModelingLib.SocialChoice.Ranking.Candidate 0)
                                              (coordinate : Fin 2) =>
                                            if candidate = 0 then if coordinate = 0 then 1 else 0
                                            else if coordinate = 0 then 0 else 1)
                                          pair.1)
                                  else 0)) →
                    ¬AppliedModelingLib.Alignment.Axioms.ParetoOptimal
                          (AppliedModelingLib.Alignment.Axioms.LinearFeasibleRanking
                            fun (candidate : AppliedModelingLib.SocialChoice.Ranking.Candidate 0)
                              (coordinate : Fin 2) =>
                            if candidate = 0 then if coordinate = 0 then 1 else 0 else if coordinate = 0 then 0 else 1)
                          rule ∧
                      ¬∀ (profile : AppliedModelingLib.Alignment.Axioms.RankingProfile Unit 0)
                          (majorityRanking : AppliedModelingLib.SocialChoice.Ranking.Ranking 0),
                          AppliedModelingLib.Alignment.Axioms.FeasibleProfile
                              (AppliedModelingLib.Alignment.Axioms.LinearFeasibleRanking
                                fun (candidate : AppliedModelingLib.SocialChoice.Ranking.Candidate 0)
                                  (coordinate : Fin 2) =>
                                if candidate = 0 then if coordinate = 0 then 1 else 0
                                else if coordinate = 0 then 0 else 1)
                              profile →
                            AppliedModelingLib.Alignment.Axioms.LinearFeasibleRanking
                                (fun (candidate : AppliedModelingLib.SocialChoice.Ranking.Candidate 0)
                                    (coordinate : Fin 2) =>
                                  if candidate = 0 then if coordinate = 0 then 1 else 0
                                  else if coordinate = 0 then 0 else 1)
                                majorityRanking →
                              (∀ (first second : AppliedModelingLib.SocialChoice.Ranking.Candidate 0),
                                  AppliedModelingLib.Alignment.Axioms.StrictlyPrefers majorityRanking first second ↔
                                    Fintype.card Unit <
                                      2 *
                                        {voter : Unit |
                                            AppliedModelingLib.Alignment.Axioms.StrictlyPrefers (profile voter) first
                                              second}.card) →
                                rule.1 profile = majorityRanking) ∨
          ∃ (weight : ℚ),
            1 / 2 < weight ∧
              weight < 1 ∧
                ∃ (δ : ℝ),
                  0 < δ ∧
                    δ < 1 ∧
                      ∃ (majorityCount : ℕ),
                        ∃ minorityCount < majorityCount,
                          0 < majorityCount + minorityCount ∧
                            ∃ (ε : ℝ),
                              0 < ε ∧
                                ε < 1 ∧
                                  ∀
                                    (rule :
                                      AppliedModelingLib.Alignment.Axioms.LinearRankAggregationRule
                                        (Fin (majorityCount + minorityCount)) 4
                                        (AppliedModelingLib.Alignment.Axioms.LinearFeasibleRanking
                                          fun (candidate : AppliedModelingLib.SocialChoice.Ranking.Candidate 4)
                                            (coordinate : Fin 2) =>
                                          if candidate = 0 then if coordinate = 0 then 2 else 1
                                          else
                                            if candidate = 1 then if coordinate = 0 then 2 - ε else 1
                                            else
                                              if candidate = 2 then if coordinate = 0 then 1 else 1
                                              else
                                                if candidate = 3 then if coordinate = 0 then 1 - ε else 1
                                                else
                                                  if candidate = 4 then if coordinate = 0 then -ε else δ * ε else 0)),
                                    (∀
                                        (profile :
                                          AppliedModelingLib.Alignment.Axioms.RankingProfile
                                            (Fin (majorityCount + minorityCount)) 4),
                                        AppliedModelingLib.Alignment.Axioms.FeasibleProfile
                                            (AppliedModelingLib.Alignment.Axioms.LinearFeasibleRanking
                                              fun (candidate : AppliedModelingLib.SocialChoice.Ranking.Candidate 4)
                                                (coordinate : Fin 2) =>
                                              if candidate = 0 then if coordinate = 0 then 2 else 1
                                              else
                                                if candidate = 1 then if coordinate = 0 then 2 - ε else 1
                                                else
                                                  if candidate = 2 then if coordinate = 0 then 1 else 1
                                                  else
                                                    if candidate = 3 then if coordinate = 0 then 1 - ε else 1
                                                    else
                                                      if candidate = 4 then if coordinate = 0 then -ε else δ * ε else 0)
                                            profile →
                                          sInf
                                              ((fun
                                                    (parameter :
                                                      AppliedModelingLib.Alignment.Axioms.LinearRewardParameter 2) =>
                                                  ∑ voter : Fin (majorityCount + minorityCount),
                                                    ∑ pair :
                                                      AppliedModelingLib.SocialChoice.Ranking.Candidate 4 ×
                                                        AppliedModelingLib.SocialChoice.Ranking.Candidate 4,
                                                      if
                                                          AppliedModelingLib.Alignment.Axioms.StrictlyPrefers
                                                            (profile voter) pair.1 pair.2 then
                                                        loss
                                                          (AppliedModelingLib.Alignment.Axioms.linearReward parameter
                                                              (fun
                                                                  (candidate :
                                                                    AppliedModelingLib.SocialChoice.Ranking.Candidate 4)
                                                                  (coordinate : Fin 2) =>
                                                                if candidate = 0 then if coordinate = 0 then 2 else 1
                                                                else
                                                                  if candidate = 1 then
                                                                    if coordinate = 0 then 2 - ε else 1
                                                                  else
                                                                    if candidate = 2 then
                                                                      if coordinate = 0 then 1 else 1
                                                                    else
                                                                      if candidate = 3 then
                                                                        if coordinate = 0 then 1 - ε else 1
                                                                      else
                                                                        if candidate = 4 then
                                                                          if coordinate = 0 then -ε else δ * ε
                                                                        else 0)
                                                              pair.2 -
                                                            AppliedModelingLib.Alignment.Axioms.linearReward parameter
                                                              (fun
                                                                  (candidate :
                                                                    AppliedModelingLib.SocialChoice.Ranking.Candidate 4)
                                                                  (coordinate : Fin 2) =>
                                                                if candidate = 0 then if coordinate = 0 then 2 else 1
                                                                else
                                                                  if candidate = 1 then
                                                                    if coordinate = 0 then 2 - ε else 1
                                                                  else
                                                                    if candidate = 2 then
                                                                      if coordinate = 0 then 1 else 1
                                                                    else
                                                                      if candidate = 3 then
                                                                        if coordinate = 0 then 1 - ε else 1
                                                                      else
                                                                        if candidate = 4 then
                                                                          if coordinate = 0 then -ε else δ * ε
                                                                        else 0)
                                                              pair.1)
                                                      else 0) ''
                                                {parameter :
                                                  AppliedModelingLib.Alignment.Axioms.LinearRewardParameter 2 |
                                                  AppliedModelingLib.Alignment.Axioms.InducesRanking
                                                    (fun
                                                        (candidate :
                                                          AppliedModelingLib.SocialChoice.Ranking.Candidate 4)
                                                        (coordinate : Fin 2) =>
                                                      if candidate = 0 then if coordinate = 0 then 2 else 1
                                                      else
                                                        if candidate = 1 then if coordinate = 0 then 2 - ε else 1
                                                        else
                                                          if candidate = 2 then if coordinate = 0 then 1 else 1
                                                          else
                                                            if candidate = 3 then if coordinate = 0 then 1 - ε else 1
                                                            else
                                                              if candidate = 4 then if coordinate = 0 then -ε else δ * ε
                                                              else 0)
                                                    parameter (rule.1 profile)}) =
                                            sInf
                                              (Set.range
                                                fun
                                                  (parameter :
                                                    AppliedModelingLib.Alignment.Axioms.LinearRewardParameter 2) =>
                                                ∑ voter : Fin (majorityCount + minorityCount),
                                                  ∑ pair :
                                                    AppliedModelingLib.SocialChoice.Ranking.Candidate 4 ×
                                                      AppliedModelingLib.SocialChoice.Ranking.Candidate 4,
                                                    if
                                                        AppliedModelingLib.Alignment.Axioms.StrictlyPrefers
                                                          (profile voter) pair.1 pair.2 then
                                                      loss
                                                        (AppliedModelingLib.Alignment.Axioms.linearReward parameter
                                                            (fun
                                                                (candidate :
                                                                  AppliedModelingLib.SocialChoice.Ranking.Candidate 4)
                                                                (coordinate : Fin 2) =>
                                                              if candidate = 0 then if coordinate = 0 then 2 else 1
                                                              else
                                                                if candidate = 1 then
                                                                  if coordinate = 0 then 2 - ε else 1
                                                                else
                                                                  if candidate = 2 then if coordinate = 0 then 1 else 1
                                                                  else
                                                                    if candidate = 3 then
                                                                      if coordinate = 0 then 1 - ε else 1
                                                                    else
                                                                      if candidate = 4 then
                                                                        if coordinate = 0 then -ε else δ * ε
                                                                      else 0)
                                                            pair.2 -
                                                          AppliedModelingLib.Alignment.Axioms.linearReward parameter
                                                            (fun
                                                                (candidate :
                                                                  AppliedModelingLib.SocialChoice.Ranking.Candidate 4)
                                                                (coordinate : Fin 2) =>
                                                              if candidate = 0 then if coordinate = 0 then 2 else 1
                                                              else
                                                                if candidate = 1 then
                                                                  if coordinate = 0 then 2 - ε else 1
                                                                else
                                                                  if candidate = 2 then if coordinate = 0 then 1 else 1
                                                                  else
                                                                    if candidate = 3 then
                                                                      if coordinate = 0 then 1 - ε else 1
                                                                    else
                                                                      if candidate = 4 then
                                                                        if coordinate = 0 then -ε else δ * ε
                                                                      else 0)
                                                            pair.1)
                                                    else 0)) →
                                      ¬AppliedModelingLib.Alignment.Axioms.ParetoOptimal
                                            (AppliedModelingLib.Alignment.Axioms.LinearFeasibleRanking
                                              fun (candidate : AppliedModelingLib.SocialChoice.Ranking.Candidate 4)
                                                (coordinate : Fin 2) =>
                                              if candidate = 0 then if coordinate = 0 then 2 else 1
                                              else
                                                if candidate = 1 then if coordinate = 0 then 2 - ε else 1
                                                else
                                                  if candidate = 2 then if coordinate = 0 then 1 else 1
                                                  else
                                                    if candidate = 3 then if coordinate = 0 then 1 - ε else 1
                                                    else
                                                      if candidate = 4 then if coordinate = 0 then -ε else δ * ε else 0)
                                            rule ∧
                                        ¬∀
                                            (profile :
                                              AppliedModelingLib.Alignment.Axioms.RankingProfile
                                                (Fin (majorityCount + minorityCount)) 4)
                                            (majorityRanking : AppliedModelingLib.SocialChoice.Ranking.Ranking 4),
                                            AppliedModelingLib.Alignment.Axioms.FeasibleProfile
                                                (AppliedModelingLib.Alignment.Axioms.LinearFeasibleRanking
                                                  fun (candidate : AppliedModelingLib.SocialChoice.Ranking.Candidate 4)
                                                    (coordinate : Fin 2) =>
                                                  if candidate = 0 then if coordinate = 0 then 2 else 1
                                                  else
                                                    if candidate = 1 then if coordinate = 0 then 2 - ε else 1
                                                    else
                                                      if candidate = 2 then if coordinate = 0 then 1 else 1
                                                      else
                                                        if candidate = 3 then if coordinate = 0 then 1 - ε else 1
                                                        else
                                                          if candidate = 4 then if coordinate = 0 then -ε else δ * ε
                                                          else 0)
                                                profile →
                                              AppliedModelingLib.Alignment.Axioms.LinearFeasibleRanking
                                                  (fun (candidate : AppliedModelingLib.SocialChoice.Ranking.Candidate 4)
                                                      (coordinate : Fin 2) =>
                                                    if candidate = 0 then if coordinate = 0 then 2 else 1
                                                    else
                                                      if candidate = 1 then if coordinate = 0 then 2 - ε else 1
                                                      else
                                                        if candidate = 2 then if coordinate = 0 then 1 else 1
                                                        else
                                                          if candidate = 3 then if coordinate = 0 then 1 - ε else 1
                                                          else
                                                            if candidate = 4 then if coordinate = 0 then -ε else δ * ε
                                                            else 0)
                                                  majorityRanking →
                                                (∀ (first second : AppliedModelingLib.SocialChoice.Ranking.Candidate 4),
                                                    AppliedModelingLib.Alignment.Axioms.StrictlyPrefers majorityRanking
                                                        first second ↔
                                                      Fintype.card (Fin (majorityCount + minorityCount)) <
                                                        2 *
                                                          {voter : Fin (majorityCount + minorityCount) |
                                                              AppliedModelingLib.Alignment.Axioms.StrictlyPrefers
                                                                (profile voter) first second}.card) →
                                                  rule.1 profile = majorityRanking
\end{ReviewVerbatim}



\paragraph{Lean proof endpoint (built separately)}\begin{ReviewVerbatim}
GeEtAl2024AlignmentAxioms.theorem3_1_lossBasedImpossibility
\end{ReviewVerbatim}

\paragraph{Recorded source-to-Spec screening}
\textbf{Verdict:} \texttt{matches}
\\
{\raggedright
\textbf{Reason:} Under the source's nonnegative-\allowbreak{}loss convention, the target has exactly the stated shape alternatives and strict infimum gap.\allowbreak{} Its positive-\allowbreak{}input two-\allowbreak{}candidate branch and negative-\allowbreak{}input six-\allowbreak{}candidate branch each give a feasible standard-\allowbreak{}loss-\allowbreak{}minimizing rule that fails both PO and PMC, which is the source theorem's finite failure conclusion.\allowbreak{}
\par}
\reviewmatch{Matches source input}
\noindent\textbf{Reviewer annotation}\par
\reviewerbox
\clearpage
\hypertarget{source-claim-31b67361459405c7}{}
\section*{Review row 3}
\textbf{Source locator:} {\footnotesize\raggedright cited publication, lines 318-\allowbreak{}319\par}
\paragraph{Verbatim source input}
\begin{ReviewVerbatim}
Lemma 3.2. There exists a rational p ∈ (1/2, 1] and values A1 < A2 with A2 > 0 such that
OP T unconstr is nonempty and for all (ra , rb ) ∈ OP T unconstr , ra > A2 and rb ≤ A1 .

[Next verbatim source excerpt]

We begin with a small instance of just three candidates C core = {a, b, c} to gain some traction on
how ell behaves. We will later extend this instance with additional candidates to demonstrate a profile
where PO and PMC fail. The candidates will have feature vectors xa := (2, 1), xb := (1, 1), and
xc := (0, 0), respectively. Furthermore, a p-fraction of voters (for p to be chosen later) will rank
a ≻ b ≻ c, while the remaining (1 − p)-fraction will have inverted preferences, ranking c ≻ b ≻ a.6
Let                                              X
                              Lcore (θ) :=                  wx≻y ell(rθ (y) − rθ (x))
                                              x̸=y∈C core

      5
     Typically, BTL is presented as choosing θ maximizing the likelihood of seeing the pairwise com-
                                                                erθ (a)
parisons we observed, assuming that Pr[a ≻ b] = erθ (a)           +erθ (b)
                                                                           . That is, we choose θ maximizing
       (                ) na≻b
            erθ (a)
Q
  a̸=b erθ (a) +erθ (b)         . By taking the log and swapping the sign, we see that this is equivalent to mini-
                    (                     )
        P                   rθ (b)−rθ (a)
mizing a̸=b log 1 + e                      .
    6
      It so happens that these rankings are feasible, but for now we will not worry about this as loss function-
based rules still make sense regardless of whether the inputs are feasible. For the final example, we will ensure
that the rankings are feasible.


                                                        5

[form-feed in source extraction]
with wx≻y ∈ {1 − p, p} be the loss function on this instance (scaling nx≻y down to wx≻y leads to
an equivalent formulation). Let g(x) = p · ell(−x) + (1 − p)ell(x). Note that we can rewrite Lcom as
                    Lcore (θ) = g(rθ (a) − rθ (b)) + g(rθ (a) − rθ (c)) + g(rθ (b) − rθ (c)).
Note that rθ (c) = 0 for all θ, so we can simplify this to
                              Lcore (θ) = g(rθ (a) − rθ (b)) + g(rθ (a)) + g(rθ (b)).
We will consider an unconstrained version of this problem where we are free to choose rewards
ra , rb ∈ R arbitrarily, and later show by which vectors θ these optimal values can be induced. That
is, we will first find ra , rb ∈ R minimizing
                                Lunconstr (ra , rb ) := g(ra − rb ) + g(ra ) + g(rb ).
Let OP T core = {θ | Lcore (θ) = inf θ′ Lcore (θ′ )} and OP T unconstr {(ra , rb ) | Lunconstr (ra , rb ) =
inf ra′ ,rb′ Lunconstr (ra′ , rb′ )} be the set of minimizers for these two loss functions. In Appendix A, we
establish the following results about these optimal sets.
Lemma 3.2. There exists a rational p ∈ (1/2, 1] and values A1 < A2 with A2 > 0 such that
OP T unconstr is nonempty and for all (ra , rb ) ∈ OP T unconstr , ra > A2 and rb ≤ A1 .

[Next verbatim source excerpt]

A.1   Proof of Lemma 3.2

Proof. We begin with some observations on g. First, we have that since ell is nonnegative, g must also
be nonnegative. This along with the fact that limx→∞ ell(x) = ∞, we have that both limx→∞ g(x) =
∞ and limx→−∞ g(x) = ∞. Together, these imply that Lunconstr attains a minimum. Indeed,
Lunconstr (0, 0) = 3g(0), and there is some bound B such that for all x > B and x < −B,
g(x) > 3g(0). We can therefore restrict the optimization problem to ra , rb ∈ [−B, B] without
changing the solutions. Since Lunconstr is continuous and [−B, B]2 is compact, a minimum is
attained.
Next, note that g is convex because compositions of convex functions with monotonic functions and
convex combinations of convex functions are convex [24]. From this, we claim that if there is an
optimal solution (ra , rb ), then (ra , ra /2) is also an optimal solution. Indeed, fix such an (ra , rb ),
                    Lunconstr (ra , ra /2) = g(ra − ra /2) + g(ra ) + g(ra /2)
                                           = g(ra ) + 2g(ra /2)
                                           = g(ra ) + 2g(1/2(ra − rb ) + 1/2rb )
                                           ≤ g(ra ) + 2(1/2g(ra − rb ) + 1/2g(rb ))
                                          = Lunconstr (ra , rb ),
where the inequality comes from convexity. This implies that (ra , ra /2) is also optimal.
By above, we have that if (rb , ra ) is optimal, it must be the case that ra minimizes
                                      h(ra ) := 2g(ra /2) + g(ra ).
Observe that h is again convex by monotonic composition and convex combinations.
Next, we will make use of the following facts about convex functions. Although they need not be
differentiable, right- and left-hand derivatives always exist. For a function k these are defined as
                                    ′            k(x + h) − k(x)
                                   k+ (x) = lim+
                                              h→0       h
                                    ′            k(x + h) − k(x)
                                   k− (x) = lim−                 .
                                           h→0          h
Further, if k is convex, we have that [24]:
                          ′      ′
                     (i) k+ and k− are nondecreasing,
                             ′        ′
                     (ii)   k− (x) ≤ k+ (x) for every x,
                                                         ′            ′
                     (iii) x minimizes k if and only if k− (x) ≤ 0 ≤ k+ (x),
                           ′      ′
                     (iv) k+ and k− are right and left continuous, respectively.

They also follow standard linearity and chain rule properties, which allow for simpler computation.
For example:
                                      ′         ′         ′
       if k(x) = aα(x) + bβ(x), then k+ (x) = aα+ (x) + bβ+ (x),
                              ′         ′                     ′         ′
       if k(x) = α(γx), then k+ (x) = γα+ (γx) if γ ≥ 0, and k+ (x) = γα− (γx) if γ < 0,
(all of these hold for left-hand derivatives by swapping the positions of + and −) [9].
Next, we claim that we can find a valid p (rational with 1/2 < p < 1) and w > 0 such that
 ′
g+ (w) > 0, while h′+ (w) < 0. To that end, expanding the first derivative, we have
                                 ′
                                g+ (x) = −pell′− (−x) + (1 − p)ell′+ (x).


As long as ell′− (−x) + ell′+ (x) > 0, this is strictly more than 0 for p satisfying
                                                   ell′+ (x)
                                         p<                     .                                      (2)
                                              ell− (−x) + ell′+ (x)
                                               ′



                                                    19

[form-feed in source extraction]
For h,
              h′+ (x) = 2 · (1/2)g′+ (x/2) + g′+ (x)


                      = −pell′− (−x/2) + (1 − p)ell′+ (x/2) − pell′− (−x) + (1 − p)ell′+ (x).
                      = −p [ell′− (−x/2) + ell′− (−x)] + (1 − p) [ell′+ (x/2) + ell′+ (x)].


As long as ell′− (−x/2) + ell′− (−x) + ell′+ (x/2) + ell′+ (x) > 0, then this is strictly less than 0 for

                                            ell′+ (x/2) + ell′+ (x)
                           p>                                               .                           (3)
                                 ell− (−x/2) + ell′− (−x) + ell′+ (x/2) + ell′+ (x)
                                  ′

Now, choose w > 0 such that ell′− (−w/2) > ell′− (−w) ≥ 0. This is possible by using the following
procedure. We know ell′− (0) > 0 (as otherwise ell(0) < ell(x) for all x < 0, contradicting our assump-
tion on ell). We will split into cases depending on whether ell is nondecreasing or strictly-convex (at
least one must be true by the theorem assumptions).
First, suppose ell is non-decreasing. This implies that ell′− (x) ≥ 0 for all x. We know that there is
a point x < 0 such that ell′− (x) < ell′− (0) (as otherwise ell would eventually become negative). Let
d = ell′− (x). Take w such that
                                   −w = sup{x | ell′− (x) = d}.
Note that ell′− (−w) = d because ell′− is left continuous (from (iv) above). Since −w < 0, so − w2 >
−w. This implies that ell′− (− w2 ) > d, as otherwise −w would not be the supremum of such points.
In addition, we have that d ≥ 0 because ell is nondecreasing.
Next, suppose ell is strictly convex. then ell′− is strictly increasing. Further, since it is left continuous
and ell′0 (0) > 0, there is a γ > 0 such that for all x ∈ [−γ, 0], ell′− (x) ≥ 0. Therefore, choosing w = γ
will do: ell′− (−γ) ≥ 0 by choice of γ, and −γ/2 > −γ, so ell′− (−γ/2) > ell′− (−γ) since ell′− is strictly
increasing.
For this choice of w, we first claim that the preconditions of denominators being positive hold for
(2) and (3). Indeed, let us write z1 , z2 , z3 , z4 for ell′− (−2w), ell′− (−w), ell′+ (w), ell′+ (2w). We know that
                                           z1 ≤ z2 ≤ z3 ≤ z4
by properties of convexity, and since ell′− (−w/2) > ell′− (−w) ≥ 0, we have that 0 ≤ z1 < z2 . The
denominators are of the form z1 + z4 and z1 + z2 + z3 + z4 , which are now both necessarily positive.
In addition, we also claim that a rational p with 1/2 < p < 1 satisfying both inequalities (2) and (3)
will exist. Note that inequality (2) can now be represented as z1z+z
                                                                  4
                                                                     4
                                                                       , and we have that
                                                  z4
                                                       ≤1
                                               z1 + z4
because 0 ≤ z1 < z4 . Additionally, inequality (3) can be represented as
                                                 z3 + z4
                                            z1 + z2 + z3 + z4
which is at least 1/2 because z3 + z4 > z1 + z2 . Finally, we have
                                        z3        z4        z4
                                             ≤         <
                                     z2 + z3   z2 + z4   z1 + z4
which implies that
                                          z3 + z4           z4
                                                       <         .
                                     z1 + z2 + z3 + z4   z1 + z4
Hence, there exists some rational p in this interval, which is necessarily between 1/2 and 1, as
needed.
To summarize, we have found a valid p and value w > 0 such that h′+ (w) < 0 while g+    ′
                                                                                          (w) > 0.
         ′                                                                        ′
Because h+ is right continuous (from (iv) above), there is some γ > 0 such that h+ (w + γ) < 0 as
well. We will now show for all (ra , rb ) ∈ OP T unconstr , ra > w + γ and rb ≤ w. Indeed, note that
ra must minimize h, so h′+ (ra ) ≥ 0 (from (iii) above), which implies ra > w + c. For rb , suppose


                                                    20

[form-feed in source extraction]
                                                ′
for a contradiction rb > w as well. Note that g+  (w) > 0 means g is increasing to the right of w. Let
                                       ′   ′
d = min(ra , rb ) − w > 0. Consider (ra , rb ) = (ra − d, rb − d). We then have
                           Lunconstr (ra′ , rb′ ) = g(ra′ − rb′ ) + g(ra′ ) + g(rb′ )
                                                = g(ra − rb ) + g(ra′ ) + g(rb′ )
                                                < g(ra − rb ) + g(ra ) + g(rb )
                                                = Lunconstr (ra , rb ).
where the equality holds because ra′ − rb′ = ra − rb and the inequality because g is increasing to the
right of w. Therefore, we reach a contradiction, because then (ra , rb ) would not be optimal. Thus,
we have found values satisfying the lemma statement with A1 = w and A2 = w + γ.
\end{ReviewVerbatim}



\paragraph{Expanded PaperInterface specification}
\begin{ReviewVerbatim}
∀ {loss : ℝ → ℝ} {negativeInput : ℝ},
  (∀ (input : ℝ), 0 ≤ loss input) →
    Monotone loss ∧ ConvexOn ℝ Set.univ loss ∨ StrictConvexOn ℝ Set.univ loss →
      negativeInput < 0 →
        loss negativeInput < loss 0 →
          ∃ (p : ℚ),
            1 / 2 < ↑p ∧
              ↑p ≤ 1 ∧
                ∃ (A1 : ℝ) (A2 : ℝ),
                  A1 < A2 ∧
                    0 < A2 ∧
                      ∃ (rewards : ℝ × ℝ),
                        (∀ (contender : ℝ × ℝ),
                            (fun (rewards : ℝ × ℝ) =>
                                  (fun (x : ℝ) => ↑p * loss (-x) + (1 - ↑p) * loss x) (rewards.1 - rewards.2) +
                                      (fun (x : ℝ) => ↑p * loss (-x) + (1 - ↑p) * loss x) rewards.1 +
                                    (fun (x : ℝ) => ↑p * loss (-x) + (1 - ↑p) * loss x) rewards.2)
                                rewards ≤
                              (fun (rewards : ℝ × ℝ) =>
                                  (fun (x : ℝ) => ↑p * loss (-x) + (1 - ↑p) * loss x) (rewards.1 - rewards.2) +
                                      (fun (x : ℝ) => ↑p * loss (-x) + (1 - ↑p) * loss x) rewards.1 +
                                    (fun (x : ℝ) => ↑p * loss (-x) + (1 - ↑p) * loss x) rewards.2)
                                contender) ∧
                          ∀ (rewards : ℝ × ℝ),
                            (∀ (contender : ℝ × ℝ),
                                (fun (rewards : ℝ × ℝ) =>
                                      (fun (x : ℝ) => ↑p * loss (-x) + (1 - ↑p) * loss x) (rewards.1 - rewards.2) +
                                          (fun (x : ℝ) => ↑p * loss (-x) + (1 - ↑p) * loss x) rewards.1 +
                                        (fun (x : ℝ) => ↑p * loss (-x) + (1 - ↑p) * loss x) rewards.2)
                                    rewards ≤
                                  (fun (rewards : ℝ × ℝ) =>
                                      (fun (x : ℝ) => ↑p * loss (-x) + (1 - ↑p) * loss x) (rewards.1 - rewards.2) +
                                          (fun (x : ℝ) => ↑p * loss (-x) + (1 - ↑p) * loss x) rewards.1 +
                                        (fun (x : ℝ) => ↑p * loss (-x) + (1 - ↑p) * loss x) rewards.2)
                                    contender) →
                              rewards.1 > A2 ∧ rewards.2 ≤ A1
\end{ReviewVerbatim}



\paragraph{Lean proof endpoint (built separately)}\begin{ReviewVerbatim}
GeEtAl2024AlignmentAxioms.lemma3_2_unconstrainedCoreSeparation
\end{ReviewVerbatim}

\paragraph{Recorded source-to-Spec screening}
\textbf{Verdict:} \texttt{matches}
\\
{\raggedright
\textbf{Reason:} Compared the negative-\allowbreak{}input context, the weighted loss g(x)=p*loss(-\allowbreak{}x)+(1-\allowbreak{}p)*loss(x), the three-\allowbreak{}term unconstrained objective g(ra-\allowbreak{}rb)+g(ra)+g(rb), and every bound.\allowbreak{} The target assumes the source loss conditions and a negative improving input, produces rational p in (1/\allowbreak{}2,1], A1<A2 with A2>0, an attained global minimum, and the stated ra>A2 and rb<=A1 bounds for every minimizer.\allowbreak{}
\par}
\reviewmatch{Matches source input}
\noindent\textbf{Reviewer annotation}\par
\reviewerbox
\clearpage
\hypertarget{source-claim-df13a69840ea5ff9}{}
\section*{Review row 4}
\textbf{Source locator:} {\footnotesize\raggedright cited publication, lines 320-\allowbreak{}322\par}
\paragraph{Verbatim source input}
\begin{ReviewVerbatim}
Lemma 3.3. Suppose Lemma 3.2 holds for values p, A1 and A2 , then, for this same choice of p,
OP T core is nonempty and there exist A3 and A4 with A3 > 0 such that for all (θ1 , θ2 ) ∈ OP T core ,
θ1 > A3 and θ2 < A4 .

[Next verbatim source excerpt]

We begin with a small instance of just three candidates C core = {a, b, c} to gain some traction on
how ell behaves. We will later extend this instance with additional candidates to demonstrate a profile
where PO and PMC fail. The candidates will have feature vectors xa := (2, 1), xb := (1, 1), and
xc := (0, 0), respectively. Furthermore, a p-fraction of voters (for p to be chosen later) will rank
a ≻ b ≻ c, while the remaining (1 − p)-fraction will have inverted preferences, ranking c ≻ b ≻ a.6
Let                                              X
                              Lcore (θ) :=                  wx≻y ell(rθ (y) − rθ (x))
                                              x̸=y∈C core

      5
     Typically, BTL is presented as choosing θ maximizing the likelihood of seeing the pairwise com-
                                                                erθ (a)
parisons we observed, assuming that Pr[a ≻ b] = erθ (a)           +erθ (b)
                                                                           . That is, we choose θ maximizing
       (                ) na≻b
            erθ (a)
Q
  a̸=b erθ (a) +erθ (b)         . By taking the log and swapping the sign, we see that this is equivalent to mini-
                    (                     )
        P                   rθ (b)−rθ (a)
mizing a̸=b log 1 + e                      .
    6
      It so happens that these rankings are feasible, but for now we will not worry about this as loss function-
based rules still make sense regardless of whether the inputs are feasible. For the final example, we will ensure
that the rankings are feasible.


                                                        5

[form-feed in source extraction]
with wx≻y ∈ {1 − p, p} be the loss function on this instance (scaling nx≻y down to wx≻y leads to
an equivalent formulation). Let g(x) = p · ell(−x) + (1 − p)ell(x). Note that we can rewrite Lcom as
                    Lcore (θ) = g(rθ (a) − rθ (b)) + g(rθ (a) − rθ (c)) + g(rθ (b) − rθ (c)).
Note that rθ (c) = 0 for all θ, so we can simplify this to
                              Lcore (θ) = g(rθ (a) − rθ (b)) + g(rθ (a)) + g(rθ (b)).
We will consider an unconstrained version of this problem where we are free to choose rewards
ra , rb ∈ R arbitrarily, and later show by which vectors θ these optimal values can be induced. That
is, we will first find ra , rb ∈ R minimizing
                                Lunconstr (ra , rb ) := g(ra − rb ) + g(ra ) + g(rb ).
Let OP T core = {θ | Lcore (θ) = inf θ′ Lcore (θ′ )} and OP T unconstr {(ra , rb ) | Lunconstr (ra , rb ) =
inf ra′ ,rb′ Lunconstr (ra′ , rb′ )} be the set of minimizers for these two loss functions. In Appendix A, we
establish the following results about these optimal sets.
Lemma 3.2. There exists a rational p ∈ (1/2, 1] and values A1 < A2 with A2 > 0 such that
OP T unconstr is nonempty and for all (ra , rb ) ∈ OP T unconstr , ra > A2 and rb ≤ A1 .
Lemma 3.3. Suppose Lemma 3.2 holds for values p, A1 and A2 , then, for this same choice of p,
OP T core is nonempty and there exist A3 and A4 with A3 > 0 such that for all (θ1 , θ2 ) ∈ OP T core ,
θ1 > A3 and θ2 < A4 .

[Next verbatim source excerpt]

A.2   Proof of Lemma 3.3

Proof. Fix p inducing a nonempty OP T unconstr satisfying Lemma 3.2 with values A1 and A2 .
We first claim that for any (ra , rb ) (regardless of optimality), it is possible to find a θ such that
rθ (a) = ra and rθ (b) = rb . Indeed, note that
                                     [[2, 1], [1, 1]] · [θ1, θ2] = [rθ(a), rθ(b)].

Since M = [[2, 1],
           [1, 1]] is invertible with inverse
                                         M^−1 = [[1, −1],
                                                 [−1, 2]],
for any ra , rb , we can simply set
                                             θ = M^−1 [ra, rb].

Thus, OP T com is nonempty, and is simply the image of OP T unconstr
under M^−1. Now, fix
(ra , rb ) ∈ OP T unconstr ,
and let θ = M^−1 [ra, rb].
By assumption, we have ra > A2 and rb ≤ A1 .









This implies that θ1 = ra − rb > A2 − A1 , while θ2 = 2rb − ra < 2A1 − A2 . Thus, setting
A3 = A2 − A1 and A4 = 2A1 − A2 satisfy the desired properties.
\end{ReviewVerbatim}



\paragraph{Expanded PaperInterface specification}
\begin{ReviewVerbatim}
∀ {g : ℝ → ℝ} {A1 A2 : ℝ},
  A1 < A2 →
    (∃ (rewards : ℝ × ℝ),
        ∀ (contender : ℝ × ℝ),
          (fun (rewards : ℝ × ℝ) => g (rewards.1 - rewards.2) + g rewards.1 + g rewards.2) rewards ≤
            (fun (rewards : ℝ × ℝ) => g (rewards.1 - rewards.2) + g rewards.1 + g rewards.2) contender) →
      (∀ (rewards : ℝ × ℝ),
          (∀ (contender : ℝ × ℝ),
              (fun (rewards : ℝ × ℝ) => g (rewards.1 - rewards.2) + g rewards.1 + g rewards.2) rewards ≤
                (fun (rewards : ℝ × ℝ) => g (rewards.1 - rewards.2) + g rewards.1 + g rewards.2) contender) →
            rewards.1 > A2 ∧ rewards.2 ≤ A1) →
        ∃ (A3 : ℝ) (A4 : ℝ),
          0 < A3 ∧
            ∃ (parameter : ℝ × ℝ),
              (∀ (contender : ℝ × ℝ),
                  (fun (parameter : ℝ × ℝ) =>
                        g
                              ((2 * parameter.1 + parameter.2, parameter.1 + parameter.2).1 -
                                (2 * parameter.1 + parameter.2, parameter.1 + parameter.2).2) +
                            g (2 * parameter.1 + parameter.2, parameter.1 + parameter.2).1 +
                          g (2 * parameter.1 + parameter.2, parameter.1 + parameter.2).2)
                      parameter ≤
                    (fun (parameter : ℝ × ℝ) =>
                        g
                              ((2 * parameter.1 + parameter.2, parameter.1 + parameter.2).1 -
                                (2 * parameter.1 + parameter.2, parameter.1 + parameter.2).2) +
                            g (2 * parameter.1 + parameter.2, parameter.1 + parameter.2).1 +
                          g (2 * parameter.1 + parameter.2, parameter.1 + parameter.2).2)
                      contender) ∧
                ∀ (parameter : ℝ × ℝ),
                  (∀ (contender : ℝ × ℝ),
                      (fun (parameter : ℝ × ℝ) =>
                            g
                                  ((2 * parameter.1 + parameter.2, parameter.1 + parameter.2).1 -
                                    (2 * parameter.1 + parameter.2, parameter.1 + parameter.2).2) +
                                g (2 * parameter.1 + parameter.2, parameter.1 + parameter.2).1 +
                              g (2 * parameter.1 + parameter.2, parameter.1 + parameter.2).2)
                          parameter ≤
                        (fun (parameter : ℝ × ℝ) =>
                            g
                                  ((2 * parameter.1 + parameter.2, parameter.1 + parameter.2).1 -
                                    (2 * parameter.1 + parameter.2, parameter.1 + parameter.2).2) +
                                g (2 * parameter.1 + parameter.2, parameter.1 + parameter.2).1 +
                              g (2 * parameter.1 + parameter.2, parameter.1 + parameter.2).2)
                          contender) →
                    parameter.1 > A3 ∧ parameter.2 < A4
\end{ReviewVerbatim}



\paragraph{Lean proof endpoint (built separately)}\begin{ReviewVerbatim}
GeEtAl2024AlignmentAxioms.lemma3_3_coreMinimizers
\end{ReviewVerbatim}

\paragraph{Recorded source-to-Spec screening}
\textbf{Verdict:} \texttt{matches}
\\
{\raggedright
\textbf{Reason:} Compared the conditional optimizer transfer in Lemma 3.\allowbreak{}3.\allowbreak{} The target assumes nonempty minimizers of the exact unconstrained three-\allowbreak{}term objective and the Lemma-\allowbreak{}3.\allowbreak{}2 reward bounds, expands the core rewards as (2*theta1+theta2, theta1+theta2), and concludes nonempty core minimizers together with A3>0 and uniform theta1>A3, theta2<A4 bounds.\allowbreak{} Its abstraction over g is used on the source weighted loss and does not change the transfer claim.\allowbreak{}
\par}
\reviewmatch{Matches source input}
\noindent\textbf{Reviewer annotation}\par
\reviewerbox
\clearpage
\hypertarget{source-claim-1682b709e6333170}{}
\section*{Review row 5}
\textbf{Source locator:} {\footnotesize\raggedright cited publication, lines 351\par}
\paragraph{Verbatim source input}
\begin{ReviewVerbatim}
Lemma 3.4. OP T (0) ⊆ Rc′ ≻c .

[Next verbatim source excerpt]

We will now explicitly construct a family of instances with candidate feature vectors parameterized
by a value ε ∈ R such that for sufficiently small ε > 0, the output of f fails the two axioms. Fix
p, A3 and A4 from Lemma 3.3, and choose δ with 0 < δ < 1 such that δA4 − A3 < 0 (δ < A3 /A4
works if A4 > 0, and otherwise, any 0 < δ < 1 will do).
Each instance will have six candidates, which we will think of as two groups of three, C = C core ∪
C copies . The first group C core = {a, b, c} will be the same as the three-candidate instance from
above, while the second group C copies = {a′ , b′ , c′ } will be new. The candidates a, b, c will still be
located at xa := (2, 1), xb := (1, 1), and xc := (0, 0), respectively. The candidates a′ , b′ , c′ will
be located near their undecorated counterparts at xa′ := xa + (−ε, 0), xb′ := xb + (−ε, 0) and
xc′ := xc + (−ε, δ · ε).
Next, we describe the voter preferences. A p-fraction of voters will have the ranking a ≻ a′ ≻ b ≻
b′ ≻ c′ ≻ c, and the remaining (1−p)-fraction of voters will have ranking c′ ≻ c ≻ b′ ≻ b ≻ a′ ≻ a.
As long as 0 < ε < 1 (which will be the case for our final chosen ε), these are both feasible rankings.
The former is induced by the nondegenerate feature vector (1, 1)7 and the latter by (−1, 0).8
For each ε ∈ R, let                                 X
                                      Lε (θ) =             wx≻y ell(rθ (y) − rθ (x))
                                                  x̸=y∈C

with wx≻y ∈ {0, 1 − p, p, 1} be the loss function we are optimizing using candidate locations
parameterized by ε.
We will show that for sufficiently small ε > 0, inf θ∈Rc′ ≻c Lε (θ) > inf θ Lε (θ). This means that
f must output a ranking with c ≻ c′ . Observe that this is a PO violation because all voters agree
that c′ ≻ c. Furthermore, this is a PMC violation because a majority of voters have the ranking
a ≻ a′ ≻ b ≻ b′ ≻ c′ ≻ c, yet this is not the output.
Let OP T (ε) be the set of vectors optimizing Lε . The rest of the proof will follow from the following
two lemmas, whose proofs are in Appendix A.

Lemma 3.4. OP T (0) ⊆ Rc′ ≻c .

Lemma 3.5. Suppose OP T (0) ⊆ Rc′ ≻c , then, for sufficiently small ε > 0, inf θ:θ∈Rc′ ≻c Lε (θ) >

[Next verbatim source excerpt]

A.3   Proof of Lemma 3.4

Proof. We first claim that OP T (0) = OP T core . This will follow from showing
                                      L0 (θ) = 4 · Lcore (θ) + 3 · ell(0).
Indeed, in L0 , the copied candidates are in exactly the same location as their counterparts. Hence,
each term in Lcore appears 4 times, one for each combination of original and copy. In addition
to these, there are the 6 terms for each ordered pair of (a, a′ ), (b, b′ ), and (c, c′ ). Note that, each
rθ (x) = rθ (x′ ) for each x ∈ {a, b, c} regardless of θ since they are in the same location. Therefore,
the ell portion is always ell(0), and the corresponding wx≻x′ and wx′ ≻x terms add up to one for each
pair. Hence, the total sum of these terms is 3 · ell(0). Since L0 is equivalent to Lcore up to positive
scaling and translation, they have the same optima.
Finally, fix a θ ∈ OP T (0). Since θ ∈ OP T core , by Lemma 3.3, θ1 > A3 and θ2 < A4 . Thus,
                      rθ (c′ ) = ε · (−θ1 + δθ2 ) < ε(−A3 + δA4 ) < 0 = rθ (c)

by choice of δ. Hence, θ ∈ Rc′ ≻c
\end{ReviewVerbatim}


\paragraph{Formalized review target}
The archival source above is not asserted equivalent to this formalized target.
\begin{ReviewVerbatim}
The Lemma 3.3 coordinate bounds put every core minimizer in the strict cone that ranks each original candidate above its copy for every positive perturbation satisfying the displayed delta bound.
\end{ReviewVerbatim}

\textbf{Archival source anchor:} {\footnotesize\raggedright cited publication, lines 351\par}
\paragraph{Expanded PaperInterface specification}
\begin{ReviewVerbatim}
∀ {g : ℝ → ℝ} {A3 A4 δ : ℝ},
  0 < A3 →
    0 < δ →
      δ * A4 - A3 < 0 →
        (∀ (parameter : ℝ × ℝ),
            (∀ (contender : ℝ × ℝ),
                (fun (parameter : ℝ × ℝ) =>
                      g
                            ((2 * parameter.1 + parameter.2, parameter.1 + parameter.2).1 -
                              (2 * parameter.1 + parameter.2, parameter.1 + parameter.2).2) +
                          g (2 * parameter.1 + parameter.2, parameter.1 + parameter.2).1 +
                        g (2 * parameter.1 + parameter.2, parameter.1 + parameter.2).2)
                    parameter ≤
                  (fun (parameter : ℝ × ℝ) =>
                      g
                            ((2 * parameter.1 + parameter.2, parameter.1 + parameter.2).1 -
                              (2 * parameter.1 + parameter.2, parameter.1 + parameter.2).2) +
                          g (2 * parameter.1 + parameter.2, parameter.1 + parameter.2).1 +
                        g (2 * parameter.1 + parameter.2, parameter.1 + parameter.2).2)
                    contender) →
              A3 < parameter.1 ∧ parameter.2 < A4) →
          ∀ (parameter : ℝ × ℝ),
            (∀ (contender : ℝ × ℝ),
                (fun (parameter : ℝ × ℝ) =>
                      g
                            ((2 * parameter.1 + parameter.2, parameter.1 + parameter.2).1 -
                              (2 * parameter.1 + parameter.2, parameter.1 + parameter.2).2) +
                          g (2 * parameter.1 + parameter.2, parameter.1 + parameter.2).1 +
                        g (2 * parameter.1 + parameter.2, parameter.1 + parameter.2).2)
                    parameter ≤
                  (fun (parameter : ℝ × ℝ) =>
                      g
                            ((2 * parameter.1 + parameter.2, parameter.1 + parameter.2).1 -
                              (2 * parameter.1 + parameter.2, parameter.1 + parameter.2).2) +
                          g (2 * parameter.1 + parameter.2, parameter.1 + parameter.2).1 +
                        g (2 * parameter.1 + parameter.2, parameter.1 + parameter.2).2)
                    contender) →
              0 < parameter.1 ∧ -parameter.1 + δ * parameter.2 < 0
\end{ReviewVerbatim}



\paragraph{Lean proof endpoint (built separately)}\begin{ReviewVerbatim}
GeEtAl2024AlignmentAxioms.lemma3_4_correctedCopyCone
\end{ReviewVerbatim}

\paragraph{Recorded source-to-Spec screening}
\textbf{Verdict:} \texttt{matches}
\\
{\raggedright
\textbf{Reason:} Compared the repaired expansion directly with the approved Lemma-\allowbreak{}3.\allowbreak{}4 correction and the six-\allowbreak{}candidate perturbation.\allowbreak{} The target now assumes 0<A3 and the Lemma-\allowbreak{}3.\allowbreak{}3 bounds, so every core minimizer has theta1>A3>0;\allowbreak{} together with 0<delta, theta2<A4, and delta*A4-\allowbreak{}A3<0 it yields -\allowbreak{}theta1+delta*theta2<0.\allowbreak{} The expanded conjunction therefore places a above a' and b above b' through theta1>0 and c above c' through the second strict inequality for every positive perturbation, exactly the approved all-\allowbreak{}copies cone.\allowbreak{}
\par}
\reviewmatch{Matches formalized target}
\noindent\textbf{Reviewer annotation}\par
\reviewerbox
\clearpage
\hypertarget{source-claim-03ca4893d06d5c73}{}
\section*{Review row 6}
\textbf{Source locator:} {\footnotesize\raggedright cited publication, lines 353-\allowbreak{}354\par}
\paragraph{Verbatim source input}
\begin{ReviewVerbatim}
Lemma 3.5. Suppose OP T (0) ⊆ Rc′ ≻c , then, for sufficiently small ε > 0, inf θ:θ∈Rc′ ≻c Lε (θ) >
inf θ Lε (θ).

[Next verbatim source excerpt]

We will now explicitly construct a family of instances with candidate feature vectors parameterized
by a value ε ∈ R such that for sufficiently small ε > 0, the output of f fails the two axioms. Fix
p, A3 and A4 from Lemma 3.3, and choose δ with 0 < δ < 1 such that δA4 − A3 < 0 (δ < A3 /A4
works if A4 > 0, and otherwise, any 0 < δ < 1 will do).
Each instance will have six candidates, which we will think of as two groups of three, C = C core ∪
C copies . The first group C core = {a, b, c} will be the same as the three-candidate instance from
above, while the second group C copies = {a′ , b′ , c′ } will be new. The candidates a, b, c will still be
located at xa := (2, 1), xb := (1, 1), and xc := (0, 0), respectively. The candidates a′ , b′ , c′ will
be located near their undecorated counterparts at xa′ := xa + (−ε, 0), xb′ := xb + (−ε, 0) and
xc′ := xc + (−ε, δ · ε).
Next, we describe the voter preferences. A p-fraction of voters will have the ranking a ≻ a′ ≻ b ≻
b′ ≻ c′ ≻ c, and the remaining (1−p)-fraction of voters will have ranking c′ ≻ c ≻ b′ ≻ b ≻ a′ ≻ a.
As long as 0 < ε < 1 (which will be the case for our final chosen ε), these are both feasible rankings.
The former is induced by the nondegenerate feature vector (1, 1)7 and the latter by (−1, 0).8
For each ε ∈ R, let                                 X
                                      Lε (θ) =             wx≻y ell(rθ (y) − rθ (x))
                                                  x̸=y∈C

with wx≻y ∈ {0, 1 − p, p, 1} be the loss function we are optimizing using candidate locations
parameterized by ε.
We will show that for sufficiently small ε > 0, inf θ∈Rc′ ≻c Lε (θ) > inf θ Lε (θ). This means that
f must output a ranking with c ≻ c′ . Observe that this is a PO violation because all voters agree
that c′ ≻ c. Furthermore, this is a PMC violation because a majority of voters have the ranking
a ≻ a′ ≻ b ≻ b′ ≻ c′ ≻ c, yet this is not the output.
Let OP T (ε) be the set of vectors optimizing Lε . The rest of the proof will follow from the following
two lemmas, whose proofs are in Appendix A.

Lemma 3.4. OP T (0) ⊆ Rc′ ≻c .

Lemma 3.5. Suppose OP T (0) ⊆ Rc′ ≻c , then, for sufficiently small ε > 0, inf θ:θ∈Rc′ ≻c Lε (θ) >
inf θ Lε (θ).

[Next verbatim source excerpt]

A.4   Proof of Lemma 3.5

Proof. We first show that when optimizing each Lε , it is sufficient to consider only θ coming from
a bounded region. Indeed, observe that Lε (0) = 62 ell(0) for all ε. Since limx→∞ ell(x) = ∞, we


                                                      21

[form-feed in source extraction]
                                                         (6)ell(0)     ε
can find some B > 0 such that for all x > B, ell(x) > 21−p = L1−p        (0)
                                                                           . For a pair of candidates
           com
x ̸= y ∈ C     , in the two terms concerning these candidates, we have
                            wx≻y ell(rθ (y) − rθ (x)) + wy≻x ell(rθ (x) − rθ (y))
                          ≥ (1 − p) (ell(rθ (y) − rθ (x)) + ell(rθ (x) − rθ (y)))
                          ≥ (1 − p)ell(|rθ (y) − rθ (x)|).

Applying this to {a, b} and {b, c},
                      Lε (θ) ≥ (1 − p) (ell(|rθ (a) − rθ (b)| + ell(|rθ (b)) − rθ (c)|))
                             = (1 − p)(ell(|θ1 |) + ell(|θ1 + θ2 |)
This implies that we may restrict our attention to θ in the region

                                Rbounded = {θ | |θ1 | ≤ B, |θ2 | ≤ 2B}.

              / Rbounded , either |θ1 | > B or |θ1 + θ2 | > B. In either case, we have Lε (θ) ≥
Indeed, for θ ∈
(1 − p)ell(B) > Lε (0).
Note that Lε (θ) is continuous not only in θ, but also in ε. Additionally, Rbounded is closed and
bounded, and hence, compact. Therefore, by Berge’s Maximum Theorem, OP T (ε) is nonempty and
upper semi-continuous in ε [3]. As per the definition of upper semi-continuous, since OP T (0) ⊆
                                                                                             ′
Rc′ ≻c , an open set, for sufficiently small ε > 0, OP T (ε) ⊆ Rc′ ≻c . Finally, note that Rc ≻c ∩
Rbounded is compact, so a minimum is attained, and this minimum must therefore be strictly larger
than the values attained by members of OP T (ε).
\end{ReviewVerbatim}


\paragraph{Formalized review target}
The archival source above is not asserted equivalent to this formalized target.
\begin{ReviewVerbatim}
Under the corrected strict minimizer cone, the literal six-candidate source objective has a strict infimum gap on the complete weak bad half-space for all sufficiently small positive perturbations.
\end{ReviewVerbatim}

\textbf{Archival source anchor:} {\footnotesize\raggedright cited publication, lines 353-\allowbreak{}354\par}
\paragraph{Expanded PaperInterface specification}
\begin{ReviewVerbatim}
∀ {loss : ℝ → ℝ} {negativeInput δ : ℝ} {weight : ℚ} {base : ℝ × ℝ},
  1 / 2 < ↑weight →
    ↑weight < 1 →
      (∀ (input : ℝ), 0 ≤ loss input) →
        ConvexOn ℝ Set.univ loss →
          negativeInput < 0 →
            loss negativeInput < loss 0 →
              (∀ (contender : ℝ × ℝ),
                  (fun (parameter : ℝ × ℝ) =>
                        (fun (x : ℝ) => ↑weight * loss (-x) + (1 - ↑weight) * loss x)
                              ((2 * parameter.1 + parameter.2, parameter.1 + parameter.2).1 -
                                (2 * parameter.1 + parameter.2, parameter.1 + parameter.2).2) +
                            (fun (x : ℝ) => ↑weight * loss (-x) + (1 - ↑weight) * loss x)
                              (2 * parameter.1 + parameter.2, parameter.1 + parameter.2).1 +
                          (fun (x : ℝ) => ↑weight * loss (-x) + (1 - ↑weight) * loss x)
                            (2 * parameter.1 + parameter.2, parameter.1 + parameter.2).2)
                      base ≤
                    (fun (parameter : ℝ × ℝ) =>
                        (fun (x : ℝ) => ↑weight * loss (-x) + (1 - ↑weight) * loss x)
                              ((2 * parameter.1 + parameter.2, parameter.1 + parameter.2).1 -
                                (2 * parameter.1 + parameter.2, parameter.1 + parameter.2).2) +
                            (fun (x : ℝ) => ↑weight * loss (-x) + (1 - ↑weight) * loss x)
                              (2 * parameter.1 + parameter.2, parameter.1 + parameter.2).1 +
                          (fun (x : ℝ) => ↑weight * loss (-x) + (1 - ↑weight) * loss x)
                            (2 * parameter.1 + parameter.2, parameter.1 + parameter.2).2)
                      contender) →
                (∀ (parameter : ℝ × ℝ),
                    (∀ (contender : ℝ × ℝ),
                        (fun (parameter : ℝ × ℝ) =>
                              (fun (x : ℝ) => ↑weight * loss (-x) + (1 - ↑weight) * loss x)
                                    ((2 * parameter.1 + parameter.2, parameter.1 + parameter.2).1 -
                                      (2 * parameter.1 + parameter.2, parameter.1 + parameter.2).2) +
                                  (fun (x : ℝ) => ↑weight * loss (-x) + (1 - ↑weight) * loss x)
                                    (2 * parameter.1 + parameter.2, parameter.1 + parameter.2).1 +
                                (fun (x : ℝ) => ↑weight * loss (-x) + (1 - ↑weight) * loss x)
                                  (2 * parameter.1 + parameter.2, parameter.1 + parameter.2).2)
                            parameter ≤
                          (fun (parameter : ℝ × ℝ) =>
                              (fun (x : ℝ) => ↑weight * loss (-x) + (1 - ↑weight) * loss x)
                                    ((2 * parameter.1 + parameter.2, parameter.1 + parameter.2).1 -
                                      (2 * parameter.1 + parameter.2, parameter.1 + parameter.2).2) +
                                  (fun (x : ℝ) => ↑weight * loss (-x) + (1 - ↑weight) * loss x)
                                    (2 * parameter.1 + parameter.2, parameter.1 + parameter.2).1 +
                                (fun (x : ℝ) => ↑weight * loss (-x) + (1 - ↑weight) * loss x)
                                  (2 * parameter.1 + parameter.2, parameter.1 + parameter.2).2)
                            contender) →
                      -parameter.1 + δ * parameter.2 < 0) →
                  ∃ (bound : ℝ) (radius : ℝ),
                    0 < bound ∧
                      0 < radius ∧
                        ∀ (ε : ℝ),
                          0 < ε →
                            ε < radius →
                              sInf
                                  ((fun (parameter : ℝ × ℝ) =>
                                      (↑weight *
                                          ∑ pair :
                                            AppliedModelingLib.SocialChoice.Ranking.Candidate 4 ×
                                              AppliedModelingLib.SocialChoice.Ranking.Candidate 4,
                                            if
                                                AppliedModelingLib.Alignment.Axioms.StrictlyPrefers
                                                  (Equiv.refl (AppliedModelingLib.SocialChoice.Ranking.Candidate 4))
                                                  pair.1 pair.2 then
                                              loss
                                                (AppliedModelingLib.Alignment.Axioms.linearReward
                                                    (fun (coordinate : Fin 2) =>
                                                      if coordinate = 0 then parameter.1 else parameter.2)
                                                    (fun
                                                        (candidate :
                                                          AppliedModelingLib.SocialChoice.Ranking.Candidate 4)
                                                        (coordinate : Fin 2) =>
                                                      if candidate = 0 then if coordinate = 0 then 2 else 1
                                                      else
                                                        if candidate = 1 then if coordinate = 0 then 2 - ε else 1
                                                        else
                                                          if candidate = 2 then if coordinate = 0 then 1 else 1
                                                          else
                                                            if candidate = 3 then if coordinate = 0 then 1 - ε else 1
                                                            else
                                                              if candidate = 4 then if coordinate = 0 then -ε else δ * ε
                                                              else 0)
                                                    pair.2 -
                                                  AppliedModelingLib.Alignment.Axioms.linearReward
                                                    (fun (coordinate : Fin 2) =>
                                                      if coordinate = 0 then parameter.1 else parameter.2)
                                                    (fun
                                                        (candidate :
                                                          AppliedModelingLib.SocialChoice.Ranking.Candidate 4)
                                                        (coordinate : Fin 2) =>
                                                      if candidate = 0 then if coordinate = 0 then 2 else 1
                                                      else
                                                        if candidate = 1 then if coordinate = 0 then 2 - ε else 1
                                                        else
                                                          if candidate = 2 then if coordinate = 0 then 1 else 1
                                                          else
                                                            if candidate = 3 then if coordinate = 0 then 1 - ε else 1
                                                            else
                                                              if candidate = 4 then if coordinate = 0 then -ε else δ * ε
                                                              else 0)
                                                    pair.1)
                                            else 0) +
                                        (1 - ↑weight) *
                                          ∑ pair :
                                            AppliedModelingLib.SocialChoice.Ranking.Candidate 4 ×
                                              AppliedModelingLib.SocialChoice.Ranking.Candidate 4,
                                            if
                                                AppliedModelingLib.Alignment.Axioms.StrictlyPrefers
                                                  (Equiv.trans (Equiv.swap 0 1) Fin.revPerm) pair.1 pair.2 then
                                              loss
                                                (AppliedModelingLib.Alignment.Axioms.linearReward
                                                    (fun (coordinate : Fin 2) =>
                                                      if coordinate = 0 then parameter.1 else parameter.2)
                                                    (fun
                                                        (candidate :
                                                          AppliedModelingLib.SocialChoice.Ranking.Candidate 4)
                                                        (coordinate : Fin 2) =>
                                                      if candidate = 0 then if coordinate = 0 then 2 else 1
                                                      else
                                                        if candidate = 1 then if coordinate = 0 then 2 - ε else 1
                                                        else
                                                          if candidate = 2 then if coordinate = 0 then 1 else 1
                                                          else
                                                            if candidate = 3 then if coordinate = 0 then 1 - ε else 1
                                                            else
                                                              if candidate = 4 then if coordinate = 0 then -ε else δ * ε
                                                              else 0)
                                                    pair.2 -
                                                  AppliedModelingLib.Alignment.Axioms.linearReward
                                                    (fun (coordinate : Fin 2) =>
                                                      if coordinate = 0 then parameter.1 else parameter.2)
                                                    (fun
                                                        (candidate :
                                                          AppliedModelingLib.SocialChoice.Ranking.Candidate 4)
                                                        (coordinate : Fin 2) =>
                                                      if candidate = 0 then if coordinate = 0 then 2 else 1
                                                      else
                                                        if candidate = 1 then if coordinate = 0 then 2 - ε else 1
                                                        else
                                                          if candidate = 2 then if coordinate = 0 then 1 else 1
                                                          else
                                                            if candidate = 3 then if coordinate = 0 then 1 - ε else 1
                                                            else
                                                              if candidate = 4 then if coordinate = 0 then -ε else δ * ε
                                                              else 0)
                                                    pair.1)
                                            else 0) ''
                                    {parameter : ℝ × ℝ |
                                      AppliedModelingLib.Alignment.Axioms.linearReward
                                          (fun (coordinate : Fin 2) =>
                                            if coordinate = 0 then parameter.1 else parameter.2)
                                          (fun (candidate : AppliedModelingLib.SocialChoice.Ranking.Candidate 4)
                                              (coordinate : Fin 2) =>
                                            if candidate = 0 then if coordinate = 0 then 2 else 1
                                            else
                                              if candidate = 1 then if coordinate = 0 then 2 - ε else 1
                                              else
                                                if candidate = 2 then if coordinate = 0 then 1 else 1
                                                else
                                                  if candidate = 3 then if coordinate = 0 then 1 - ε else 1
                                                  else
                                                    if candidate = 4 then if coordinate = 0 then -ε else δ * ε else 0)
                                          5 ≤
                                        AppliedModelingLib.Alignment.Axioms.linearReward
                                          (fun (coordinate : Fin 2) =>
                                            if coordinate = 0 then parameter.1 else parameter.2)
                                          (fun (candidate : AppliedModelingLib.SocialChoice.Ranking.Candidate 4)
                                              (coordinate : Fin 2) =>
                                            if candidate = 0 then if coordinate = 0 then 2 else 1
                                            else
                                              if candidate = 1 then if coordinate = 0 then 2 - ε else 1
                                              else
                                                if candidate = 2 then if coordinate = 0 then 1 else 1
                                                else
                                                  if candidate = 3 then if coordinate = 0 then 1 - ε else 1
                                                  else
                                                    if candidate = 4 then if coordinate = 0 then -ε else δ * ε else 0)
                                          4}) >
                                sInf
                                  (Set.range fun (parameter : ℝ × ℝ) =>
                                    (↑weight *
                                        ∑ pair :
                                          AppliedModelingLib.SocialChoice.Ranking.Candidate 4 ×
                                            AppliedModelingLib.SocialChoice.Ranking.Candidate 4,
                                          if
                                              AppliedModelingLib.Alignment.Axioms.StrictlyPrefers
                                                (Equiv.refl (AppliedModelingLib.SocialChoice.Ranking.Candidate 4))
                                                pair.1 pair.2 then
                                            loss
                                              (AppliedModelingLib.Alignment.Axioms.linearReward
                                                  (fun (coordinate : Fin 2) =>
                                                    if coordinate = 0 then parameter.1 else parameter.2)
                                                  (fun (candidate : AppliedModelingLib.SocialChoice.Ranking.Candidate 4)
                                                      (coordinate : Fin 2) =>
                                                    if candidate = 0 then if coordinate = 0 then 2 else 1
                                                    else
                                                      if candidate = 1 then if coordinate = 0 then 2 - ε else 1
                                                      else
                                                        if candidate = 2 then if coordinate = 0 then 1 else 1
                                                        else
                                                          if candidate = 3 then if coordinate = 0 then 1 - ε else 1
                                                          else
                                                            if candidate = 4 then if coordinate = 0 then -ε else δ * ε
                                                            else 0)
                                                  pair.2 -
                                                AppliedModelingLib.Alignment.Axioms.linearReward
                                                  (fun (coordinate : Fin 2) =>
                                                    if coordinate = 0 then parameter.1 else parameter.2)
                                                  (fun (candidate : AppliedModelingLib.SocialChoice.Ranking.Candidate 4)
                                                      (coordinate : Fin 2) =>
                                                    if candidate = 0 then if coordinate = 0 then 2 else 1
                                                    else
                                                      if candidate = 1 then if coordinate = 0 then 2 - ε else 1
                                                      else
                                                        if candidate = 2 then if coordinate = 0 then 1 else 1
                                                        else
                                                          if candidate = 3 then if coordinate = 0 then 1 - ε else 1
                                                          else
                                                            if candidate = 4 then if coordinate = 0 then -ε else δ * ε
                                                            else 0)
                                                  pair.1)
                                          else 0) +
                                      (1 - ↑weight) *
                                        ∑ pair :
                                          AppliedModelingLib.SocialChoice.Ranking.Candidate 4 ×
                                            AppliedModelingLib.SocialChoice.Ranking.Candidate 4,
                                          if
                                              AppliedModelingLib.Alignment.Axioms.StrictlyPrefers
                                                (Equiv.trans (Equiv.swap 0 1) Fin.revPerm) pair.1 pair.2 then
                                            loss
                                              (AppliedModelingLib.Alignment.Axioms.linearReward
                                                  (fun (coordinate : Fin 2) =>
                                                    if coordinate = 0 then parameter.1 else parameter.2)
                                                  (fun (candidate : AppliedModelingLib.SocialChoice.Ranking.Candidate 4)
                                                      (coordinate : Fin 2) =>
                                                    if candidate = 0 then if coordinate = 0 then 2 else 1
                                                    else
                                                      if candidate = 1 then if coordinate = 0 then 2 - ε else 1
                                                      else
                                                        if candidate = 2 then if coordinate = 0 then 1 else 1
                                                        else
                                                          if candidate = 3 then if coordinate = 0 then 1 - ε else 1
                                                          else
                                                            if candidate = 4 then if coordinate = 0 then -ε else δ * ε
                                                            else 0)
                                                  pair.2 -
                                                AppliedModelingLib.Alignment.Axioms.linearReward
                                                  (fun (coordinate : Fin 2) =>
                                                    if coordinate = 0 then parameter.1 else parameter.2)
                                                  (fun (candidate : AppliedModelingLib.SocialChoice.Ranking.Candidate 4)
                                                      (coordinate : Fin 2) =>
                                                    if candidate = 0 then if coordinate = 0 then 2 else 1
                                                    else
                                                      if candidate = 1 then if coordinate = 0 then 2 - ε else 1
                                                      else
                                                        if candidate = 2 then if coordinate = 0 then 1 else 1
                                                        else
                                                          if candidate = 3 then if coordinate = 0 then 1 - ε else 1
                                                          else
                                                            if candidate = 4 then if coordinate = 0 then -ε else δ * ε
                                                            else 0)
                                                  pair.1)
                                          else 0)
\end{ReviewVerbatim}



\paragraph{Lean proof endpoint (built separately)}\begin{ReviewVerbatim}
GeEtAl2024AlignmentAxioms.lemma3_5_correctedPerturbationGap
\end{ReviewVerbatim}

\paragraph{Recorded source-to-Spec screening}
\textbf{Verdict:} \texttt{matches}
\\
{\raggedright
\textbf{Reason:} Compared the approved strict-\allowbreak{}gap correction rather than the false printed weak-\allowbreak{}halfspace statement.\allowbreak{} The target assumes every zero-\allowbreak{}perturbation core minimizer has -\allowbreak{}theta1+delta*theta2<0, uses the literal six-\allowbreak{}candidate features, the majority order a>a'>b>b'>c'>c with weight p and the reverse minority order with weight 1-\allowbreak{}p, and concludes for all sufficiently small positive epsilon that the infimum on the complete weak half-\allowbreak{}space reward(c)<=reward(c') is strictly above the unrestricted infimum.\allowbreak{}
\par}
\reviewmatch{Matches formalized target}
\noindent\textbf{Reviewer annotation}\par
\reviewerbox
\clearpage
\hypertarget{source-claim-bbe7cd3338235aa4}{}
\section*{Review row 7}
\textbf{Source locator:} {\footnotesize\raggedright cited publication, lines 389-\allowbreak{}390\par}
\paragraph{Verbatim source input}
\begin{ReviewVerbatim}
Theorem 3.6. Fix a nondecreasing loss function ell with ell(0) > inf x ell(x). If a linear rank aggrega-
tion rule f minimizes ell in the majority formulation, then f satisfies PMC.

[Next verbatim source excerpt]

3.2   Majority-Based Loss Formulation

Despite the negative results for loss-function-based rules, we may hope for a remedy using slightly
different information. Specifically, we consider a similar loss-based function that rather than getting
penalized for disagreeing with each voter only gets penalized if it disagrees with a majority of voters.
That is, we choose θ minimizing
                                       X
                    Lmaj (θ; π, ell) =         I[wa≻b (π) > 1/2] · ell(rθ(b) − rθ(a) ).
                                       a̸=b∈C

Defining a parameter aggregation function based on this loss suffers from the same caveat as before,
that in some cases no optimal θ exists. Nevertheless, we can apply an analogous fix for a ranking
variant. We say that a linear rank aggregation rule f minimizes ell in the majority formulation if for
all σ = f (π),
                               inf σ Lmaj (θ; π, ell) = inf Lmaj (θ; π, ell).
                               θ:θ∈R                      θ

We first show (in Appendix A.5) that this does indeed help achieve PMC with essentially all loss
functions.
Theorem 3.6. Fix a nondecreasing loss function ell with ell(0) > inf x ell(x). If a linear rank aggrega-
tion rule f minimizes ell in the majority formulation, then f satisfies PMC.

[Next verbatim source excerpt]

A.5   Proof of Theorem 3.6

Proof. Without loss of generality, we may assume that inf x ell(x) = 0, as otherwise we could trans-
late ell without affecting the optimization problem. Fix a profile π with feasible PMC ranking σ, and
let θP M C be a non-degenerate parameter vector that induces σ.
First, we show that inf θ Lmaj (θ; π, ell) = 0. Indeed, note that c · θP M C ∈ Rπ for all c > 0. Further,
note that for any a, b with wa≻b (π) > 1/2, rθP M C (b) − rθP M C (a) < 0. Therefore, by making c
large, the nonzero terms in Lmaj will have an input to ell negative and becoming arbitrarily large in
magnitude. Since ell is nondecreasing, these approach the infemum of 0.
Next, for any σ ′ ̸= σ, inf θ Lmaj (θ; π, ell) ≥ ell(0). Indeed, there must be some pair of candidates a, b
                                               ′
with a ≻σ b and b ≻σ′ a. For any θ ∈ Rσ , rθ (b) ≥ rθ (a), so ell(rθ (b) − rθ (a)) ≥ ell(0), and this
lower bounds the loss function.
\end{ReviewVerbatim}



\paragraph{Expanded PaperInterface specification}
\begin{ReviewVerbatim}
∀ {Voter : Type} [inst : Fintype Voter] {n dimension : ℕ}
  (features :
    AppliedModelingLib.SocialChoice.Ranking.Candidate n → AppliedModelingLib.Alignment.Axioms.FeatureVector dimension)
  (loss : ℝ → ℝ)
  (rule :
    AppliedModelingLib.Alignment.Axioms.LinearRankAggregationRule Voter n
      (AppliedModelingLib.Alignment.Axioms.LinearFeasibleRanking features)),
  (∀ (input : ℝ), 0 ≤ loss input) →
    Monotone loss →
      loss 0 > sInf (Set.range loss) →
        GeEtAl2024AlignmentAxioms.IsMajorityLossMinimizing features loss rule →
          ∀ (profile : AppliedModelingLib.Alignment.Axioms.RankingProfile Voter n)
            (majorityRanking : AppliedModelingLib.SocialChoice.Ranking.Ranking n),
            AppliedModelingLib.Alignment.Axioms.FeasibleProfile
                (AppliedModelingLib.Alignment.Axioms.LinearFeasibleRanking features) profile →
              AppliedModelingLib.Alignment.Axioms.LinearFeasibleRanking features majorityRanking →
                (∀ (first second : AppliedModelingLib.SocialChoice.Ranking.Candidate n),
                    AppliedModelingLib.Alignment.Axioms.StrictlyPrefers majorityRanking first second ↔
                      Fintype.card Voter <
                        2 *
                          {voter : Voter |
                              AppliedModelingLib.Alignment.Axioms.StrictlyPrefers (profile voter) first second}.card) →
                  rule.1 profile = majorityRanking
\end{ReviewVerbatim}



\paragraph{Lean proof endpoint (built separately)}\begin{ReviewVerbatim}
GeEtAl2024AlignmentAxioms.theorem3_6_majorityLoss
\end{ReviewVerbatim}

\paragraph{Recorded source-to-Spec screening}
\textbf{Verdict:} \texttt{matches}
\\
{\raggedright
\textbf{Reason:} The target assumes the source loss conditions:\allowbreak{} nonnegative, nondecreasing, and loss(0) strictly above its infimum, together with majority-\allowbreak{}formulation minimization.\allowbreak{} It concludes pairwise majority consistency, the formal PMC condition stated by Theorem 3.\allowbreak{}6.\allowbreak{}
\par}
\reviewmatch{Matches source input}
\noindent\textbf{Reviewer annotation}\par
\reviewerbox
\clearpage
\hypertarget{source-claim-d5b15047bfa92842}{}
\section*{Review row 8}
\textbf{Source locator:} {\footnotesize\raggedright cited publication, lines 397-\allowbreak{}410\par}
\paragraph{Verbatim source input}
\begin{ReviewVerbatim}
However, despite this good news for PMC, we show that this does not help in achieving PO. In fact,
our negative result extends to every C1 linear rank aggregation rule. Note that if f minimizing ell in
the majority formulation breaks ties consistently (i.e., if multiple feasible rankings are optimal, then
it consistently chooses the same one), then it is C1. We then have the following result, whose proof
is relegated to Appendix A.6.
Theorem 3.7. All C1 linear rank aggregation rules fail PO.

This result is quite unfortunate, because if there were a rule that is both C1 and PO, we would
automatically achieve PMC: Whenever there is a feasible PMC ranking, a C1 rule cannot distinguish
between this profile and a profile where all voters submit this ranking, hence, under the PO criterion,
it must output it. Furthermore, whenever there is a PMC ranking, outputting it is necessarily PO,
as for every pair, a majority of voters agree with the PMC ranking. Interestingly, in the proof, we
construct a profile which has a PMC ranking, yet it is not feasible, and no matter how a C1 linear
rank aggregation rule breaks ties, there is an underlying profile in which this output violates PO.

[Next verbatim source excerpt]

We pay special attention to a prominent family of rules from social choice theory referred to as C1
rules [14], whose outputs depend only on majority relationships, i.e., they only need to know for
each pair of candidates (a, b) whether the majority prefers a or b.
In our study, we examine several axioms borrowed from social choice theory to evaluate the reason-
ableness (fairness) of our aggregation mechanisms. These axioms include:
Definition 2.1 (Pareto Optimality). A linear rank aggregation rule f satisfies Pareto optimality if,
whenever every voter prefers candidate a over candidate b on π, i.e., wa≻b (π) = 1, then candidate
a is ranked higher than candidate b in the output ranking, i.e., a ≻f (π) b.
Definition 2.2 (Pairwise Majority Consistency (PMC)). A ranking σ is called a PMC ranking for
profile π if for all a, b ∈ C, a ≻σ b if and only if a majority of voters rank a ≻σi b, i.e., wa≻b > 1/2.
A linear rank aggregation rule satisfies PMC if, when a PMC ranking σ exists for the input profile
π and σ is feasible, then f (π) = σ.

Note that a PMC ranking for each π need not exist, but when one does, it is unique. The words “σ
is feasible” allude to the possibility that no non-degenerate parameter vector θ induces the unique
PMC ranking. Indeed, we have such an example; see Appendix B for details.

[Next verbatim source excerpt]

A.6   Proof of Theorem 3.7

Proof. We construct an explicit instance and pairwise majority relationships such that no matter
what feasible ranking a rule picks, there is an underlying profile where that output was a PO viola-
tion.
                                                       −
We will have 9 candidates; 8 will be labeled c+                                              ∗
                                              i and ci for i = 1, 2, 3, 4, and one labeled c .
                                   4        ±
They will have feature vectors in R . Each ci will be located at xc± = ±ei where ei is the i’th
                                                                           i
                                                        −
standard basis vector, i.e., c+                                                          ∗
                              2 is at (0, 1, 0, 0) and c4 is at (0, 0, 0, −1). Finally, c will be located at
(1/5, 1/5, 1/5, 1/5).
There will be 5 voters. Their pairwise majority graph will be as follows. Candidate c∗ will pairwise
                                                                −
beat all others. In addition, each c+                                         +
                                    i will pairwise beat each cj . Among the ci candidates, there will
be a cycle c1 ≻ c2 ≻ c3 ≻ c4 ≻ c1 and between the remaining two pairs c+
             +      +     +      +    +                                              +       +
                                                                               1 ≻ c3 and c2 ≻ c4 .
                                                                                                   +

The c−                                                               −    −     −     −      −
      i candidates will be the exact reverse of this, i.e., a cycle c4 ≻ c3 ≻ c2 ≻ c1 ≻ c4 , along
       −     −        −     −
with c3 ≻ c1 and c4 ≻ c2 . A pictorial representation can be found in Figure 1.
A C1 rule must pick a θ solely based on the pairwise majority graph. We will show that regardless
of what θ it outputs, this will lead to a PO violation.


                                                    22

[form-feed in source extraction]
                                 c+
                                  1                c+
                                                    2                          c−
                                                                                1             c−
                                                                                               2



[Next verbatim source excerpt]

However, despite this good news for PMC, we show that this does not help in achieving PO. In fact,
our negative result extends to every C1 linear rank aggregation rule. Note that if f minimizing ell in
the majority formulation breaks ties consistently (i.e., if multiple feasible rankings are optimal, then
it consistently chooses the same one), then it is C1. We then have the following result, whose proof
is relegated to Appendix A.6.
Theorem 3.7. All C1 linear rank aggregation rules fail PO.

This result is quite unfortunate, because if there were a rule that is both C1 and PO, we would
automatically achieve PMC: Whenever there is a feasible PMC ranking, a C1 rule cannot distinguish
between this profile and a profile where all voters submit this ranking, hence, under the PO criterion,
it must output it. Furthermore, whenever there is a PMC ranking, outputting it is necessarily PO,
as for every pair, a majority of voters agree with the PMC ranking. Interestingly, in the proof, we
construct a profile which has a PMC ranking, yet it is not feasible, and no matter how a C1 linear
rank aggregation rule breaks ties, there is an underlying profile in which this output violates PO.
\end{ReviewVerbatim}


\paragraph{Formalized review target}
The archival source above is not asserted equivalent to this formalized target.
\begin{ReviewVerbatim}
Every C1 linear rank aggregation rule fails Pareto optimality; C1 plus Pareto optimality implies PMC; and any existing PMC ranking respects every unanimous comparison. Correcting the attached proof-location sentence, the Appendix-A.6 profile has a cyclic majority relation and no PMC ranking; Appendix B supplies the separate infeasible-PMC example.
\end{ReviewVerbatim}

\textbf{Archival source anchor:} {\footnotesize\raggedright cited publication, lines 404-\allowbreak{}410\par}
\paragraph{Expanded PaperInterface specification}
\begin{ReviewVerbatim}
(∀
    (rule :
      AppliedModelingLib.Alignment.Axioms.LinearRankAggregationRule (Fin 5) 7
        (AppliedModelingLib.Alignment.Axioms.LinearFeasibleRanking
          fun (candidate : AppliedModelingLib.SocialChoice.Ranking.Candidate 7) (coordinate : Fin 4) =>
          if candidate = ⟨↑coordinate, ⋯⟩ then 1
          else
            if candidate = ⟨↑coordinate + 4, ⋯⟩ then -1
            else if candidate = ⟨8, GeEtAl2024AlignmentAxioms.c1Star._proof_2⟩ then 1 / 5 else 0)),
    AppliedModelingLib.Alignment.Axioms.C1LinearRankAggregationRule
        (AppliedModelingLib.Alignment.Axioms.LinearFeasibleRanking
          fun (candidate : AppliedModelingLib.SocialChoice.Ranking.Candidate 7) (coordinate : Fin 4) =>
          if candidate = ⟨↑coordinate, ⋯⟩ then 1
          else
            if candidate = ⟨↑coordinate + 4, ⋯⟩ then -1
            else if candidate = ⟨8, GeEtAl2024AlignmentAxioms.c1Star._proof_2⟩ then 1 / 5 else 0)
        rule →
      ¬AppliedModelingLib.Alignment.Axioms.ParetoOptimal
          (AppliedModelingLib.Alignment.Axioms.LinearFeasibleRanking
            fun (candidate : AppliedModelingLib.SocialChoice.Ranking.Candidate 7) (coordinate : Fin 4) =>
            if candidate = ⟨↑coordinate, ⋯⟩ then 1
            else
              if candidate = ⟨↑coordinate + 4, ⋯⟩ then -1
              else if candidate = ⟨8, GeEtAl2024AlignmentAxioms.c1Star._proof_2⟩ then 1 / 5 else 0)
          rule) ∧
  (∀ (Voter : Type) [inst : Fintype Voter] (n : ℕ) (feasible : AppliedModelingLib.SocialChoice.Ranking.Ranking n → Prop)
      (rule : AppliedModelingLib.Alignment.Axioms.LinearRankAggregationRule Voter n feasible),
      AppliedModelingLib.Alignment.Axioms.C1LinearRankAggregationRule feasible rule →
        AppliedModelingLib.Alignment.Axioms.ParetoOptimal feasible rule →
          ∀ (profile : AppliedModelingLib.Alignment.Axioms.RankingProfile Voter n)
            (majorityRanking : AppliedModelingLib.SocialChoice.Ranking.Ranking n),
            AppliedModelingLib.Alignment.Axioms.FeasibleProfile feasible profile →
              feasible majorityRanking →
                (∀ (first second : AppliedModelingLib.SocialChoice.Ranking.Candidate n),
                    AppliedModelingLib.Alignment.Axioms.StrictlyPrefers majorityRanking first second ↔
                      Fintype.card Voter <
                        2 *
                          {voter : Voter |
                              AppliedModelingLib.Alignment.Axioms.StrictlyPrefers (profile voter) first second}.card) →
                  rule.1 profile = majorityRanking) ∧
    (∀ (Voter : Type) [inst : Fintype Voter] [Nonempty Voter] (n : ℕ)
        (profile : AppliedModelingLib.Alignment.Axioms.RankingProfile Voter n)
        (ranking : AppliedModelingLib.SocialChoice.Ranking.Ranking n),
        (∀ (first second : AppliedModelingLib.SocialChoice.Ranking.Candidate n),
            AppliedModelingLib.Alignment.Axioms.StrictlyPrefers ranking first second ↔
              Fintype.card Voter <
                2 *
                  {voter : Voter |
                      AppliedModelingLib.Alignment.Axioms.StrictlyPrefers (profile voter) first second}.card) →
          ∀ (first second : AppliedModelingLib.SocialChoice.Ranking.Candidate n),
            (∀ (voter : Voter), AppliedModelingLib.Alignment.Axioms.StrictlyPrefers (profile voter) first second) →
              AppliedModelingLib.Alignment.Axioms.StrictlyPrefers ranking first second) ∧
      ¬∃ (ranking : AppliedModelingLib.SocialChoice.Ranking.Ranking 7),
          ∀ (first second : AppliedModelingLib.SocialChoice.Ranking.Candidate 7),
            AppliedModelingLib.Alignment.Axioms.StrictlyPrefers ranking first second ↔
              Fintype.card (Fin 5) <
                2 *
                  {voter : Fin 5 |
                      AppliedModelingLib.Alignment.Axioms.StrictlyPrefers
                        ((fun (x : Fin 5) =>
                            Fin.casesOn x fun (val : ℕ) (isLt : val < 5) =>
                              Nat.casesOn (motive := fun (x : ℕ) =>
                                (isLt : x < 5) →
                                  (fun (x : Fin 5) => AppliedModelingLib.SocialChoice.Ranking.Ranking 7) ⟨x, isLt⟩)
                                val
                                (fun (isLt : Nat.zero < 5) =>
                                  (fun (isLt : 0 < 5) =>
                                      ((((Equiv.trans (Equiv.swap 1 8) (Equiv.swap 1 4)).trans (Equiv.swap 1 3)).trans
                                                (Equiv.swap 1 2)).trans
                                            (Equiv.swap 5 7)).trans
                                        {
                                          toFun :=
                                            fun (candidate : AppliedModelingLib.SocialChoice.Ranking.Candidate 7) =>
                                            if hplus : ↑candidate < 4 then ⟨↑((Equiv.swap 0 1) ⟨↑candidate, hplus⟩), ⋯⟩
                                            else
                                              if hminus : ↑candidate < 8 then
                                                ⟨↑((Equiv.swap 0 1) ⟨↑candidate - 4, ⋯⟩) + 4, ⋯⟩
                                              else ⟨8, GeEtAl2024AlignmentAxioms.c1Star._proof_2⟩,
                                          invFun :=
                                            fun (candidate : AppliedModelingLib.SocialChoice.Ranking.Candidate 7) =>
                                            if hplus : ↑candidate < 4 then
                                              ⟨↑((Equiv.symm (Equiv.swap 0 1)) ⟨↑candidate, hplus⟩), ⋯⟩
                                            else
                                              if hminus : ↑candidate < 8 then
                                                ⟨↑((Equiv.symm (Equiv.swap 0 1)) ⟨↑candidate - 4, ⋯⟩) + 4, ⋯⟩
                                              else ⟨8, GeEtAl2024AlignmentAxioms.c1Star._proof_2⟩,
                                          left_inv := ⋯, right_inv := ⋯ })
                                    isLt)
                                (fun (n : ℕ) (isLt : n.succ < 5) =>
                                  Nat.casesOn (motive := fun (x : ℕ) =>
                                    (isLt : x.succ < 5) →
                                      (fun (x : Fin 5) => AppliedModelingLib.SocialChoice.Ranking.Ranking 7)
                                        ⟨x.succ, isLt⟩)
                                    n
                                    (fun (isLt : Nat.zero.succ < 5) =>
                                      (fun (isLt : 1 < 5) =>
                                          ((((Equiv.trans (Equiv.swap 1 8) (Equiv.swap 1 4)).trans
                                                        (Equiv.swap 1 3)).trans
                                                    (Equiv.swap 1 2)).trans
                                                (Equiv.swap 5 7)).trans
                                            {
                                              toFun :=
                                                fun (candidate : AppliedModelingLib.SocialChoice.Ranking.Candidate 7) =>
                                                if hplus : ↑candidate < 4 then
                                                  ⟨↑((Equiv.trans (Equiv.swap 0 2) (Equiv.swap 0 1))
                                                        ⟨↑candidate, hplus⟩),
                                                    ⋯⟩
                                                else
                                                  if hminus : ↑candidate < 8 then
                                                    ⟨↑((Equiv.trans (Equiv.swap 0 2) (Equiv.swap 0 1))
                                                            ⟨↑candidate - 4, ⋯⟩) +
                                                        4,
                                                      ⋯⟩
                                                  else ⟨8, GeEtAl2024AlignmentAxioms.c1Star._proof_2⟩,
                                              invFun :=
                                                fun (candidate : AppliedModelingLib.SocialChoice.Ranking.Candidate 7) =>
                                                if hplus : ↑candidate < 4 then
                                                  ⟨↑((Equiv.trans (Equiv.swap 0 2) (Equiv.swap 0 1)).symm
                                                        ⟨↑candidate, hplus⟩),
                                                    ⋯⟩
                                                else
                                                  if hminus : ↑candidate < 8 then
                                                    ⟨↑((Equiv.trans (Equiv.swap 0 2) (Equiv.swap 0 1)).symm
                                                            ⟨↑candidate - 4, ⋯⟩) +
                                                        4,
                                                      ⋯⟩
                                                  else ⟨8, GeEtAl2024AlignmentAxioms.c1Star._proof_2⟩,
                                              left_inv := ⋯, right_inv := ⋯ })
                                        isLt)
                                    (fun (n : ℕ) (isLt : n.succ.succ < 5) =>
                                      Nat.casesOn (motive := fun (x : ℕ) =>
                                        (isLt : x.succ.succ < 5) →
                                          (fun (x : Fin 5) => AppliedModelingLib.SocialChoice.Ranking.Ranking 7)
                                            ⟨x.succ.succ, isLt⟩)
                                        n
                                        (fun (isLt : Nat.zero.succ.succ < 5) =>
                                          (fun (isLt : 2 < 5) =>
                                              ((((Equiv.trans (Equiv.swap 1 8) (Equiv.swap 1 4)).trans
                                                            (Equiv.swap 1 3)).trans
                                                        (Equiv.swap 1 2)).trans
                                                    (Equiv.swap 5 7)).trans
                                                {
                                                  toFun :=
                                                    fun
                                                      (candidate :
                                                        AppliedModelingLib.SocialChoice.Ranking.Candidate 7) =>
                                                    if hplus : ↑candidate < 4 then
                                                      ⟨↑((Equiv.trans (Equiv.swap 0 2) (Equiv.swap 0 3))
                                                            ⟨↑candidate, hplus⟩),
                                                        ⋯⟩
                                                    else
                                                      if hminus : ↑candidate < 8 then
                                                        ⟨↑((Equiv.trans (Equiv.swap 0 2) (Equiv.swap 0 3))
                                                                ⟨↑candidate - 4, ⋯⟩) +
                                                            4,
                                                          ⋯⟩
                                                      else ⟨8, GeEtAl2024AlignmentAxioms.c1Star._proof_2⟩,
                                                  invFun :=
                                                    fun
                                                      (candidate :
                                                        AppliedModelingLib.SocialChoice.Ranking.Candidate 7) =>
                                                    if hplus : ↑candidate < 4 then
                                                      ⟨↑((Equiv.trans (Equiv.swap 0 2) (Equiv.swap 0 3)).symm
                                                            ⟨↑candidate, hplus⟩),
                                                        ⋯⟩
                                                    else
                                                      if hminus : ↑candidate < 8 then
                                                        ⟨↑((Equiv.trans (Equiv.swap 0 2) (Equiv.swap 0 3)).symm
                                                                ⟨↑candidate - 4, ⋯⟩) +
                                                            4,
                                                          ⋯⟩
                                                      else ⟨8, GeEtAl2024AlignmentAxioms.c1Star._proof_2⟩,
                                                  left_inv := ⋯, right_inv := ⋯ })
                                            isLt)
                                        (fun (n : ℕ) (isLt : n.succ.succ.succ < 5) =>
                                          Nat.casesOn (motive := fun (x : ℕ) =>
                                            (isLt : x.succ.succ.succ < 5) →
                                              (fun (x : Fin 5) => AppliedModelingLib.SocialChoice.Ranking.Ranking 7)
                                                ⟨x.succ.succ.succ, isLt⟩)
                                            n
                                            (fun (isLt : Nat.zero.succ.succ.succ < 5) =>
                                              (fun (isLt : 3 < 5) =>
                                                  ((((Equiv.trans (Equiv.swap 1 8) (Equiv.swap 1 4)).trans
                                                                (Equiv.swap 1 3)).trans
                                                            (Equiv.swap 1 2)).trans
                                                        (Equiv.swap 5 7)).trans
                                                    {
                                                      toFun :=
                                                        fun
                                                          (candidate :
                                                            AppliedModelingLib.SocialChoice.Ranking.Candidate 7) =>
                                                        if hplus : ↑candidate < 4 then
                                                          ⟨↑(((Equiv.trans (Equiv.swap 0 3) (Equiv.swap 0 2)).trans
                                                                  (Equiv.swap 0 1))
                                                                ⟨↑candidate, hplus⟩),
                                                            ⋯⟩
                                                        else
                                                          if hminus : ↑candidate < 8 then
                                                            ⟨↑(((Equiv.trans (Equiv.swap 0 3) (Equiv.swap 0 2)).trans
                                                                      (Equiv.swap 0 1))
                                                                    ⟨↑candidate - 4, ⋯⟩) +
                                                                4,
                                                              ⋯⟩
                                                          else ⟨8, GeEtAl2024AlignmentAxioms.c1Star._proof_2⟩,
                                                      invFun :=
                                                        fun
                                                          (candidate :
                                                            AppliedModelingLib.SocialChoice.Ranking.Candidate 7) =>
                                                        if hplus : ↑candidate < 4 then
                                                          ⟨↑(((Equiv.trans (Equiv.swap 0 3) (Equiv.swap 0 2)).trans
                                                                    (Equiv.swap 0 1)).symm
                                                                ⟨↑candidate, hplus⟩),
                                                            ⋯⟩
                                                        else
                                                          if hminus : ↑candidate < 8 then
                                                            ⟨↑(((Equiv.trans (Equiv.swap 0 3) (Equiv.swap 0 2)).trans
                                                                        (Equiv.swap 0 1)).symm
                                                                    ⟨↑candidate - 4, ⋯⟩) +
                                                                4,
                                                              ⋯⟩
                                                          else ⟨8, GeEtAl2024AlignmentAxioms.c1Star._proof_2⟩,
                                                      left_inv := ⋯, right_inv := ⋯ })
                                                isLt)
                                            (fun (n : ℕ) (isLt : n.succ.succ.succ.succ < 5) =>
                                              Nat.casesOn (motive := fun (x : ℕ) =>
                                                (isLt : x.succ.succ.succ.succ < 5) →
                                                  (fun (x : Fin 5) => AppliedModelingLib.SocialChoice.Ranking.Ranking 7)
                                                    ⟨x.succ.succ.succ.succ, isLt⟩)
                                                n
                                                (fun (isLt : Nat.zero.succ.succ.succ.succ < 5) =>
                                                  (fun (isLt : 4 < 5) =>
                                                      ((((Equiv.trans (Equiv.swap 1 8) (Equiv.swap 1 4)).trans
                                                                    (Equiv.swap 1 3)).trans
                                                                (Equiv.swap 1 2)).trans
                                                            (Equiv.swap 5 7)).trans
                                                        {
                                                          toFun :=
                                                            fun
                                                              (candidate :
                                                                AppliedModelingLib.SocialChoice.Ranking.Candidate 7) =>
                                                            if hplus : ↑candidate < 4 then
                                                              ⟨↑(((Equiv.trans (Equiv.swap 0 3) (Equiv.swap 0 2)).trans
                                                                      (Equiv.swap 0 1))
                                                                    ⟨↑candidate, hplus⟩),
                                                                ⋯⟩
                                                            else
                                                              if hminus : ↑candidate < 8 then
                                                                ⟨↑(((Equiv.trans (Equiv.swap 0 3)
                                                                              (Equiv.swap 0 2)).trans
                                                                          (Equiv.swap 0 1))
                                                                        ⟨↑candidate - 4, ⋯⟩) +
                                                                    4,
                                                                  ⋯⟩
                                                              else ⟨8, GeEtAl2024AlignmentAxioms.c1Star._proof_2⟩,
                                                          invFun :=
                                                            fun
                                                              (candidate :
                                                                AppliedModelingLib.SocialChoice.Ranking.Candidate 7) =>
                                                            if hplus : ↑candidate < 4 then
                                                              ⟨↑(((Equiv.trans (Equiv.swap 0 3) (Equiv.swap 0 2)).trans
                                                                        (Equiv.swap 0 1)).symm
                                                                    ⟨↑candidate, hplus⟩),
                                                                ⋯⟩
                                                            else
                                                              if hminus : ↑candidate < 8 then
                                                                ⟨↑(((Equiv.trans (Equiv.swap 0 3)
                                                                                (Equiv.swap 0 2)).trans
                                                                            (Equiv.swap 0 1)).symm
                                                                        ⟨↑candidate - 4, ⋯⟩) +
                                                                    4,
                                                                  ⋯⟩
                                                              else ⟨8, GeEtAl2024AlignmentAxioms.c1Star._proof_2⟩,
                                                          left_inv := ⋯, right_inv := ⋯ })
                                                    isLt)
                                                (fun (n : ℕ) (isLt : n.succ.succ.succ.succ.succ < 5) => ⋯.elim) isLt)
                                            isLt)
                                        isLt)
                                    isLt)
                                isLt)
                          voter)
                        first second}.card
\end{ReviewVerbatim}



\paragraph{Lean proof endpoint (built separately)}\begin{ReviewVerbatim}
GeEtAl2024AlignmentAxioms.theorem3_7_C1FailsParetoAndConsequences
\end{ReviewVerbatim}

\paragraph{Recorded source-to-Spec screening}
\textbf{Verdict:} \texttt{matches}
\\
{\raggedright
\textbf{Reason:} The expanded target matches the approval-\allowbreak{}bound corrected target and the archival theorem bundle:\allowbreak{} on the exact Appendix-\allowbreak{}A.\allowbreak{}6 instance it uses nine candidates with the eight feature vectors plus or minus the four coordinate vectors and the ninth vector (1/\allowbreak{}5,1/\allowbreak{}5,1/\allowbreak{}5,1/\allowbreak{}5), and it quantifies every rule on that instance, with dependence only on the directed strict-\allowbreak{}majority relation implying failure of the source Pareto condition.\allowbreak{} Its next two universal conjuncts state exactly that C1 plus Pareto optimality implies PMC and that an existing PMC ranking respects every unanimous comparison;\allowbreak{} Nonempty Voter is the necessary source-\allowbreak{}domain condition for unanimity to imply a strict majority.\allowbreak{} Finally, the displayed five-\allowbreak{}voter profile realizes the Appendix-\allowbreak{}A.\allowbreak{}6 majority graph, including the plus cycle, reverse minus cycle, cross-\allowbreak{}block victories, and the ninth candidate beating all others, so no ranking realizes that cyclic relation.\allowbreak{} This is precisely the recorded correction, which disclaims the archival sentence claiming an infeasible PMC ranking there and assigns that separate example to Appendix B.\allowbreak{}
\par}
\reviewmatch{Matches formalized target}
\noindent\textbf{Reviewer annotation}\par
\reviewerbox
\clearpage
\hypertarget{source-claim-11bcc2903523a4df}{}
\section*{Review row 9}
\textbf{Source locator:} {\footnotesize\raggedright cited publication, lines 459\par}
\paragraph{Verbatim source input}
\begin{ReviewVerbatim}
Theorem 4.3. LCPO satisfies PO, PMC, majority consistency and winner monotonicity.

[Next verbatim source excerpt]

4     Social Choice-Based Rule
In light of the above negative results, in this section, we ask whether there are linear rank aggre-
gation rules that concurrently satisfy our two core axioms, PO and PMC. We answer this question
affirmatively by presenting a new method based on a prominent rule from voting theory.
The Copeland rule assigns a Copeland score to each alternative equal to the number of other alter-
natives it beats in a pairwise competition, i.e., the score for a is |{b | wa≻b > 1/2}|. It then ranks
the candidates in descending order according to their Copeland scores (breaking ties arbitrarily). It


                                                     7

[form-feed in source extraction]
is known that Copeland satisfies PO, PMC, and additional axiomatic properties. However, in linear
social choice, since not every ranking is feasible, we cannot always output the Copeland ranking.
We, therefore, define a new linear rank aggregation rule, which we call leximax Copeland. This
rule chooses a feasible ranking as follows. It ranks first the candidate with the highest Copeland
score that can be feasibly ranked first under some parameter vector θ. Subject to this first position,
it ranks second the candidate with the highest Copeland score which can be feasibly ranked second,
and continues this process for subsequent positions.
Copeland’s rule is a C1 rule because it only requires the majority relationships between the can-
didates. Analogously, leximax Copeland is also a C1 linear rank aggregation rule. Therefore, by
Theorem 3.7, it does not satisfy the PO criterion. To address this issue, we define a variant called
leximax Copeland subject to PO (LCPO), which incorporates the PO criterion. Under LCPO, for
every pair of alternatives where one dominates the other, the rule restricts rankings to place the
dominating alternative above the dominated one.
The rule remains well-defined since the set of feasible rankings when enforcing the PO criterion is
non-empty, as whenever a dominates b, all the rankings in the input profile rank a above b. Note that
if the Copeland ranking is feasible, then this rule outputs that ranking, since unrestricted Copeland
satisfies PO.
In addition to PO and PMC, we wish to show that LCPO satisfies two additional properties, which
we define presently.
Definition 4.1 (majority consistency). A linear rank aggregation rule satisfies majority consistency
if when a candidate a is ranked first by a majority of voters in the input profile, a is ranked first in
the output ranking.

Majority consistency ensures that the collective decision reflects the preference of the majority when
there is a clear favorite. This principle aligns with PMC, but specifically focuses on the majority’s
favorite alternative. However, as we discussed above, a PMC ranking does not necessarily exist, and
even when it exists, it is not necessarily feasible. By contrast, when a majority winner exists, this
candidate is necessarily ranked first by a majority of voters in the input profile, who themselves (by
assumption) submit feasible rankings. Therefore, we need not handle the case where it is impossible
to rank the majority winner first.
Definition 4.2 (winner monotonicity). A linear rank aggregation rule satisfies winner monotonicity
if, when a candidate a is ranked first in the output ranking, elevating a in any voter’s preference
does not cause a to lose their top position in the updated aggregate ranking.

Winner monotonicity ensures that improving a leading candidate’s position among individual voters
will not result in that candidate’s demotion.
We now state and prove the main result of this section.
Theorem 4.3. LCPO satisfies PO, PMC, majority consistency and winner monotonicity.

[Next verbatim source excerpt]

Proof. LCPO trivially satisfies PO since it always outputs a ranking that respects the PO criterion.
Moreover, since Copeland satisfies PMC, and whenever Copeland’s ranking is in the domain, lexi-
max Copeland subject to PO returns this ranking, it clearly satisfies PMC.
Note that if an alternative a is ranked first by at least half of the voters, then a has the highest
Copeland score, meaning that leximax Copeland subject to PO will rank this candidate first if this is
possible. We see that this is indeed possible, by noticing that there is at least one feasible ranking in
the input profile where a is ranked first, and any such input ranking is feasible (by assumption) and
satisfies the PO requirement. Therefore, majority consistency is satisfied.
It remains to show that LCPO satisfies winner monotonicity. Suppose that on input profile π, the
rule outputs a ranking σ where candidate a is ranked first. Now, consider a profile π ′ which is similar
to π with the only exception being a ranking in which a is placed in a higher position. Let S be the
set of agents that ranked above a in Copeland’s ranking under π and let S ′ be the set of agents that
ranked above a in Copeland’s ranking under π ′ . Note that S ′ ⊆ S, since when moving from π to π ′ ,
only the Copeland score of a can increase, and therefore it is not possible for a candidate b to beat a
under π ′ but not under π.


                                                   8

[form-feed in source extraction]
Now, suppose that R and R′ are the set of rankings that satisfy the PO criterion with respect to π and
π ′ , respectively. We show that R′ ⊆ R. First, note that since a is ranked first under π, no alternative
dominates a in π, as otherwise the PO criterion would be violated. Therefore, we get that no other
alternative dominates a in π ′ as well. Moreover, note that if b dominates c in π, then this remains
true in π ′ as well. On the other hand, it is possible that a dominates an alternative b in π ′ but not in
π. From all of the above, we conclude that R′ ⊆ R.
Since a is ranked first in σ, we get that for every candidate b ∈ S, there is no ranking in R in which
b is ranked first, since otherwise, LCPO would output such a ranking. This also means that for every
candidate b in S ′ , there is no ranking in R′ in which b is ranked first, since R′ ⊆ R and S ′ ⊆ S.
Moreover, note that every ranking in R in which a is ranked first is also in R′ since it satisfies all the
PO restrictions of π ′ . Therefore, under π ′ , LCPO outputs a ranking in which a is ranked first.

Leximax Copeland subject to PO can be implemented in polynomial time by solving O(|C|2 ) rela-
tively small linear programs. Specifically, given an input profile, we sequentially choose the candi-
date that is ranked in position r + 1 as follows. We denote by σr the partial ranking, where the first
r positions have been fixed. For each candidate c that has not been ranked yet, we want to check
if there is a parameter vector that adheres to the partial ranking σr , respects the Pareto optimality
criterion and ranks c at position r + 1. Since all these constraints can be expressed as pairwise
comparisons, we can use a linear program such as the one described in Footnote 4 to check if such
a feasible ranking exists. Among the candidates meeting this criterion, we select the one with the
highest Copeland score for position r.
\end{ReviewVerbatim}



\paragraph{Expanded PaperInterface specification}
\begin{ReviewVerbatim}
∀ {Voter : Type} [inst : Fintype Voter] [inst_1 : Nonempty Voter] {n : ℕ}
  (feasible : AppliedModelingLib.SocialChoice.Ranking.Ranking n → Prop)
  (fallback : AppliedModelingLib.SocialChoice.Ranking.Ranking n) (hfallback : feasible fallback),
  AppliedModelingLib.Alignment.Axioms.ParetoOptimal feasible
      (GeEtAl2024AlignmentAxioms.fixedTieLCPO feasible fallback hfallback) ∧
    (∀ (profile : AppliedModelingLib.Alignment.Axioms.RankingProfile Voter n)
        (majorityRanking : AppliedModelingLib.SocialChoice.Ranking.Ranking n),
        AppliedModelingLib.Alignment.Axioms.FeasibleProfile feasible profile →
          feasible majorityRanking →
            (∀ (first second : AppliedModelingLib.SocialChoice.Ranking.Candidate n),
                AppliedModelingLib.Alignment.Axioms.StrictlyPrefers majorityRanking first second ↔
                  Fintype.card Voter <
                    2 *
                      {voter : Voter |
                          AppliedModelingLib.Alignment.Axioms.StrictlyPrefers (profile voter) first second}.card) →
              (GeEtAl2024AlignmentAxioms.fixedTieLCPO feasible fallback hfallback).1 profile = majorityRanking) ∧
      AppliedModelingLib.Alignment.Axioms.MajorityConsistent feasible
          (GeEtAl2024AlignmentAxioms.fixedTieLCPO feasible fallback hfallback) ∧
        AppliedModelingLib.Alignment.Axioms.WinnerMonotonic feasible
          (GeEtAl2024AlignmentAxioms.fixedTieLCPO feasible fallback hfallback)
\end{ReviewVerbatim}



\paragraph{Lean proof endpoint (built separately)}\begin{ReviewVerbatim}
GeEtAl2024AlignmentAxioms.theorem4_3_constructedFixedTieLCPO
\end{ReviewVerbatim}

\paragraph{Recorded source-to-Spec screening}
\textbf{Verdict:} \texttt{matches}
\\
{\raggedright
\textbf{Reason:} Compared LCPO's feasible-\allowbreak{}and-\allowbreak{}Pareto search domain, sequential Copeland maximization, and all four axioms.\allowbreak{} On every feasible profile the expanded fixedTieLCPO maximizes the lexicographic vector of Copeland scores with a profile-\allowbreak{}independent candidate tie break over exactly the feasible rankings respecting unanimous comparisons;\allowbreak{} fallback behavior is confined to infeasible profiles.\allowbreak{} The conclusion is Pareto optimality, PMC, strict-\allowbreak{}majority top-\allowbreak{}choice consistency, and the source one-\allowbreak{}voter winner-\allowbreak{}elevation monotonicity.\allowbreak{}
\par}
\reviewmatch{Matches source input}
\noindent\textbf{Reviewer annotation}\par
\reviewerbox
\clearpage
\hypertarget{source-claim-42097890307741b0}{}
\hypertarget{source-section-1cf4eccafb9ce99a}{}
\section*{Appendix results}
\subsection*{Review row 10}
\textbf{Source locator:} {\footnotesize\raggedright cited publication, lines 1333-\allowbreak{}1385\par}
\paragraph{Verbatim source input}
\begin{ReviewVerbatim}
B    PMC Infeasibility Example

Consider the case with d = 3 and seven candidates: one special candidate a∗ located at
(1/4, 1/4, 1/4), and six others c±                                    ±
                                 i located at standard basis vectors ei . We have three voters with
parameter vectors (1, 2ε, ε), (2ε, 1, ε), and (2ε, ε, 1), where ε < 1/5 is a small positive number.
These voters have the following induced rankings:


                                                        23

[form-feed in source extraction]
                                          Rank          v1       v2   v3
                                              1         c+
                                                         1       c+
                                                                  2   c+
                                                                       3
                                                         ∗        ∗
                                              2         a        a    a∗
                                              3         c+
                                                         2       c+
                                                                  1   c+
                                                                       1
                                              4         c+
                                                         3       c+
                                                                  3   c+
                                                                       2
                                              5         c−
                                                         3       c−
                                                                  3   c−
                                                                       2
                                              6         c−
                                                         2       c−
                                                                  1   c−
                                                                       1
                                              7         c−
                                                         1       c−
                                                                  2   c−
                                                                       3

                                                −   − −
We argue the ranking a∗ ≻ c+       +      +
                             1 ≻ c2 ≻ c3 ≻ c3 ≻ c2 ≻ c1 is a PMC ranking but no linear reward
                       ∗
function can position a at the top of this ranking.
For any reward vector θ = (θ1 , θ2 , θ3 ) ∈ R3 , the reward for a∗ is :
                                                     1
                                         raθ∗ =        (θ1 + θ2 + θ3 )
                                                     4
Given the placement of c±                                        ±
                          i at the standard basis vectors, each ci achieves a reward equivalent to one
of the absolute values of the components of θ, thus surpassing raθ∗ since
                                                  raθ∗ < max |θi |.
                                                             i


[Next verbatim source excerpt]

2     The Linear Social Choice Model
Let C be a set of m distinct prompt/responses, referred to as candidates, and let V = {1, . . . , n} be
a set of n human participants, known as voters. We denote by Rd the d-dimensional real space in
which both candidate feature vectors and chosen parameter vectors lie.
Each candidate c ∈ C is associated with a distinct feature vector xc ∈ Rd . A parameter vector
θ ∈ Rd induces a linear reward function rθ : C 7→ R defined by taking the dot product with feature
vectors rθ (c) = ⟨θ, xc ⟩. We will primarily be interested in how these parameterized functions rank
the candidates by reward. Let Ra≻b = {θ | rθ (a) ≥ rθ (b)} be the region where the reward of a is
at least as large as that of b. Note that Ra≻b and Rb≻a split Rd into two half spaces, separated by
the hyperplane orthogonal to xa − xb . Parameter vectors θ on the hyperplane have rθ (a) = rθ (b),
while rankings in the interior of either half-space strictly rank one over the other.
For a ranking σ over the candidates, we say that θ induces σ, denoted θ ▷ σ, if a ≻σ b implies
rθ (a) ≥ rθ (b). Let Rσ = {θ | θ ▷ σ} be the set of vectors θ thatTinduce it. Note that this can
be written as the intersection of corresponding half spaces Rσ = a,b:a≻σ b Ra≻b . Further, the


                                                    3

[form-feed in source extraction]
collection of {Rσ } essentially form a partition of Rd , covering the space and intersecting only at
their boundaries.
We call a θ non-degenerate if it is fully on one side of each of the separating hyperplanes, i.e.,
rθ (a) ̸= rθ (b) for all a, b ∈ C. Non-degenerate parameter vectors lie in the interior of some Rσ ,
and thus induce exactly one ranking. We call σ feasible if Rσ has a nonempty interior, i.e., is induced
by some nondegenerate θ.3
Each voter i ∈ V submits a ranking over the candidate σi . We assume that the feature space is
rich enough that voter preferences can be captured via non-degenerate parameter vectors. In other
words, we assume that each σi is feasible. We refer to the vector of voter rankings π = (σi )i∈V as
a profile. Further, for two candidates a, b, we write na≻b (π) := |{i ∈ V | a ≻σi b}| for the number
of voters that prefer a to b, and wa≻b (π) = na≻b (π)/n for the proportion of such voters. When the
profile π is clear from context, we may shorten these to na≻b and wa≻b , respectively.
We define a parameter aggregation rule as a function that takes as input a profile π and outputs a
parameter vector θ∗ . Our goal is to design parameter aggregation rules such that rθ∗ satisfies desir-
able properties with respect to the voter preferences. However, as the properties we care about will
only be with respect to how rθ∗ ranks the candidates, it will be more convenient to work with what
we call linear rank aggregation rules that take as input a profile π and output a feasible ranking σ.
There is a natural way to interpret a parameter aggregation rule as a linear rank aggregation rule,
namely, output any feasible ranking induced by θ∗ . The exact properties of the parameter aggrega-
tion rule could in principle be sensitive to the tie-breaking of non-degenerate outputs, however, all
of our results will be robust to such tie-breaking.4
We pay special attention to a prominent family of rules from social choice theory referred to as C1
rules [14], whose outputs depend only on majority relationships, i.e., they only need to know for
each pair of candidates (a, b) whether the majority prefers a or b.
In our study, we examine several axioms borrowed from social choice theory to evaluate the reason-
ableness (fairness) of our aggregation mechanisms. These axioms include:
Definition 2.1 (Pareto Optimality). A linear rank aggregation rule f satisfies Pareto optimality if,
whenever every voter prefers candidate a over candidate b on π, i.e., wa≻b (π) = 1, then candidate
a is ranked higher than candidate b in the output ranking, i.e., a ≻f (π) b.
Definition 2.2 (Pairwise Majority Consistency (PMC)). A ranking σ is called a PMC ranking for
profile π if for all a, b ∈ C, a ≻σ b if and only if a majority of voters rank a ≻σi b, i.e., wa≻b > 1/2.
A linear rank aggregation rule satisfies PMC if, when a PMC ranking σ exists for the input profile
π and σ is feasible, then f (π) = σ.

Note that a PMC ranking for each π need not exist, but when one does, it is unique. The words “σ
is feasible” allude to the possibility that no non-degenerate parameter vector θ induces the unique
PMC ranking. Indeed, we have such an example; see Appendix B for details.

[Next verbatim source excerpt]

B    PMC Infeasibility Example

Consider the case with d = 3 and seven candidates: one special candidate a∗ located at
(1/4, 1/4, 1/4), and six others c±                                    ±
                                 i located at standard basis vectors ei . We have three voters with
parameter vectors (1, 2ε, ε), (2ε, 1, ε), and (2ε, ε, 1), where ε < 1/5 is a small positive number.
These voters have the following induced rankings:


                                                        23

[form-feed in source extraction]
                                          Rank          v1       v2   v3
                                              1         c+
                                                         1       c+
                                                                  2   c+
                                                                       3
                                                         ∗        ∗
                                              2         a        a    a∗
                                              3         c+
                                                         2       c+
                                                                  1   c+
                                                                       1
                                              4         c+
                                                         3       c+
                                                                  3   c+
                                                                       2
                                              5         c−
                                                         3       c−
                                                                  3   c−
                                                                       2
                                              6         c−
                                                         2       c−
                                                                  1   c−
                                                                       1
                                              7         c−
                                                         1       c−
                                                                  2   c−
                                                                       3

                                                −   − −
We argue the ranking a∗ ≻ c+       +      +
                             1 ≻ c2 ≻ c3 ≻ c3 ≻ c2 ≻ c1 is a PMC ranking but no linear reward
                       ∗
function can position a at the top of this ranking.
For any reward vector θ = (θ1 , θ2 , θ3 ) ∈ R3 , the reward for a∗ is :
                                                     1
                                         raθ∗ =        (θ1 + θ2 + θ3 )
                                                     4
Given the placement of c±                                        ±
                          i at the standard basis vectors, each ci achieves a reward equivalent to one
of the absolute values of the components of θ, thus surpassing raθ∗ since
                                                  raθ∗ < max |θi |.
                                                             i
\end{ReviewVerbatim}



\paragraph{Expanded PaperInterface specification}
\begin{ReviewVerbatim}
(AppliedModelingLib.Alignment.Axioms.FeasibleProfile
    (AppliedModelingLib.Alignment.Axioms.LinearFeasibleRanking
      fun (candidate : AppliedModelingLib.SocialChoice.Ranking.Candidate 5) (coordinate : Fin 3) =>
      if candidate = ⟨↑coordinate, ⋯⟩ then 1
      else
        if candidate = ⟨↑coordinate + 3, ⋯⟩ then -1
        else if candidate = ⟨6, GeEtAl2024AlignmentAxioms.appendixBStar._proof_2⟩ then 1 / 4 else 0)
    fun (x : Fin 3) =>
    Fin.casesOn x fun (val : ℕ) (isLt : val < 3) =>
      Nat.casesOn (motive := fun (x : ℕ) =>
        (isLt : x < 3) → (fun (x : Fin 3) => AppliedModelingLib.SocialChoice.Ranking.Ranking 5) ⟨x, isLt⟩) val
        (fun (isLt : Nat.zero < 3) =>
          (fun (isLt : 0 < 3) =>
              ((((Equiv.refl (AppliedModelingLib.SocialChoice.Ranking.Candidate 5)).trans (Equiv.swap 1 6)).trans
                        (Equiv.swap 2 1)).trans
                    (Equiv.swap 3 2)).trans
                (Equiv.swap 4 5))
            isLt)
        (fun (n : ℕ) (isLt : n.succ < 3) =>
          Nat.casesOn (motive := fun (x : ℕ) =>
            (isLt : x.succ < 3) → (fun (x : Fin 3) => AppliedModelingLib.SocialChoice.Ranking.Ranking 5) ⟨x.succ, isLt⟩)
            n
            (fun (isLt : Nat.zero.succ < 3) =>
              (fun (isLt : 1 < 3) =>
                  ((((((Equiv.refl (AppliedModelingLib.SocialChoice.Ranking.Candidate 5)).trans (Equiv.swap 0 1)).trans
                                    (Equiv.swap 0 6)).trans
                                (Equiv.swap 2 0)).trans
                            (Equiv.swap 3 2)).trans
                        (Equiv.swap 4 5)).trans
                    (Equiv.swap 4 3))
                isLt)
            (fun (n : ℕ) (isLt : n.succ.succ < 3) =>
              Nat.casesOn (motive := fun (x : ℕ) =>
                (isLt : x.succ.succ < 3) →
                  (fun (x : Fin 3) => AppliedModelingLib.SocialChoice.Ranking.Ranking 5) ⟨x.succ.succ, isLt⟩)
                n
                (fun (isLt : Nat.zero.succ.succ < 3) =>
                  (fun (isLt : 2 < 3) =>
                      ((((Equiv.refl (AppliedModelingLib.SocialChoice.Ranking.Candidate 5)).trans
                                    (Equiv.swap 0 2)).trans
                                (Equiv.swap 1 6)).trans
                            (Equiv.swap 3 1)).trans
                        (Equiv.swap 5 3))
                    isLt)
                (fun (n : ℕ) (isLt : n.succ.succ.succ < 3) => ⋯.elim) isLt)
            isLt)
        isLt) ∧
  (∀ (first second : AppliedModelingLib.SocialChoice.Ranking.Candidate 5),
      AppliedModelingLib.Alignment.Axioms.StrictlyPrefers
          ((((((Equiv.refl (AppliedModelingLib.SocialChoice.Ranking.Candidate 5)).trans (Equiv.swap 0 6)).trans
                        (Equiv.swap 1 0)).trans
                    (Equiv.swap 2 1)).trans
                (Equiv.swap 3 2)).trans
            (Equiv.swap 4 5))
          first second ↔
        Fintype.card (Fin 3) <
          2 *
            {voter : Fin 3 |
                AppliedModelingLib.Alignment.Axioms.StrictlyPrefers
                  ((fun (x : Fin 3) =>
                      Fin.casesOn x fun (val : ℕ) (isLt : val < 3) =>
                        Nat.casesOn (motive := fun (x : ℕ) =>
                          (isLt : x < 3) →
                            (fun (x : Fin 3) => AppliedModelingLib.SocialChoice.Ranking.Ranking 5) ⟨x, isLt⟩)
                          val
                          (fun (isLt : Nat.zero < 3) =>
                            (fun (isLt : 0 < 3) =>
                                ((((Equiv.refl (AppliedModelingLib.SocialChoice.Ranking.Candidate 5)).trans
                                              (Equiv.swap 1 6)).trans
                                          (Equiv.swap 2 1)).trans
                                      (Equiv.swap 3 2)).trans
                                  (Equiv.swap 4 5))
                              isLt)
                          (fun (n : ℕ) (isLt : n.succ < 3) =>
                            Nat.casesOn (motive := fun (x : ℕ) =>
                              (isLt : x.succ < 3) →
                                (fun (x : Fin 3) => AppliedModelingLib.SocialChoice.Ranking.Ranking 5) ⟨x.succ, isLt⟩)
                              n
                              (fun (isLt : Nat.zero.succ < 3) =>
                                (fun (isLt : 1 < 3) =>
                                    ((((((Equiv.refl (AppliedModelingLib.SocialChoice.Ranking.Candidate 5)).trans
                                                          (Equiv.swap 0 1)).trans
                                                      (Equiv.swap 0 6)).trans
                                                  (Equiv.swap 2 0)).trans
                                              (Equiv.swap 3 2)).trans
                                          (Equiv.swap 4 5)).trans
                                      (Equiv.swap 4 3))
                                  isLt)
                              (fun (n : ℕ) (isLt : n.succ.succ < 3) =>
                                Nat.casesOn (motive := fun (x : ℕ) =>
                                  (isLt : x.succ.succ < 3) →
                                    (fun (x : Fin 3) => AppliedModelingLib.SocialChoice.Ranking.Ranking 5)
                                      ⟨x.succ.succ, isLt⟩)
                                  n
                                  (fun (isLt : Nat.zero.succ.succ < 3) =>
                                    (fun (isLt : 2 < 3) =>
                                        ((((Equiv.refl (AppliedModelingLib.SocialChoice.Ranking.Candidate 5)).trans
                                                      (Equiv.swap 0 2)).trans
                                                  (Equiv.swap 1 6)).trans
                                              (Equiv.swap 3 1)).trans
                                          (Equiv.swap 5 3))
                                      isLt)
                                  (fun (n : ℕ) (isLt : n.succ.succ.succ < 3) => ⋯.elim) isLt)
                              isLt)
                          isLt)
                    voter)
                  first second}.card) ∧
    (∀ (ranking : AppliedModelingLib.SocialChoice.Ranking.Ranking 5),
        (∀ (first second : AppliedModelingLib.SocialChoice.Ranking.Candidate 5),
            AppliedModelingLib.Alignment.Axioms.StrictlyPrefers ranking first second ↔
              Fintype.card (Fin 3) <
                2 *
                  {voter : Fin 3 |
                      AppliedModelingLib.Alignment.Axioms.StrictlyPrefers
                        ((fun (x : Fin 3) =>
                            Fin.casesOn x fun (val : ℕ) (isLt : val < 3) =>
                              Nat.casesOn (motive := fun (x : ℕ) =>
                                (isLt : x < 3) →
                                  (fun (x : Fin 3) => AppliedModelingLib.SocialChoice.Ranking.Ranking 5) ⟨x, isLt⟩)
                                val
                                (fun (isLt : Nat.zero < 3) =>
                                  (fun (isLt : 0 < 3) =>
                                      ((((Equiv.refl (AppliedModelingLib.SocialChoice.Ranking.Candidate 5)).trans
                                                    (Equiv.swap 1 6)).trans
                                                (Equiv.swap 2 1)).trans
                                            (Equiv.swap 3 2)).trans
                                        (Equiv.swap 4 5))
                                    isLt)
                                (fun (n : ℕ) (isLt : n.succ < 3) =>
                                  Nat.casesOn (motive := fun (x : ℕ) =>
                                    (isLt : x.succ < 3) →
                                      (fun (x : Fin 3) => AppliedModelingLib.SocialChoice.Ranking.Ranking 5)
                                        ⟨x.succ, isLt⟩)
                                    n
                                    (fun (isLt : Nat.zero.succ < 3) =>
                                      (fun (isLt : 1 < 3) =>
                                          ((((((Equiv.refl (AppliedModelingLib.SocialChoice.Ranking.Candidate 5)).trans
                                                                (Equiv.swap 0 1)).trans
                                                            (Equiv.swap 0 6)).trans
                                                        (Equiv.swap 2 0)).trans
                                                    (Equiv.swap 3 2)).trans
                                                (Equiv.swap 4 5)).trans
                                            (Equiv.swap 4 3))
                                        isLt)
                                    (fun (n : ℕ) (isLt : n.succ.succ < 3) =>
                                      Nat.casesOn (motive := fun (x : ℕ) =>
                                        (isLt : x.succ.succ < 3) →
                                          (fun (x : Fin 3) => AppliedModelingLib.SocialChoice.Ranking.Ranking 5)
                                            ⟨x.succ.succ, isLt⟩)
                                        n
                                        (fun (isLt : Nat.zero.succ.succ < 3) =>
                                          (fun (isLt : 2 < 3) =>
                                              ((((Equiv.refl
                                                                (AppliedModelingLib.SocialChoice.Ranking.Candidate
                                                                  5)).trans
                                                            (Equiv.swap 0 2)).trans
                                                        (Equiv.swap 1 6)).trans
                                                    (Equiv.swap 3 1)).trans
                                                (Equiv.swap 5 3))
                                            isLt)
                                        (fun (n : ℕ) (isLt : n.succ.succ.succ < 3) => ⋯.elim) isLt)
                                    isLt)
                                isLt)
                          voter)
                        first second}.card) →
          ranking =
            (((((Equiv.refl (AppliedModelingLib.SocialChoice.Ranking.Candidate 5)).trans (Equiv.swap 0 6)).trans
                          (Equiv.swap 1 0)).trans
                      (Equiv.swap 2 1)).trans
                  (Equiv.swap 3 2)).trans
              (Equiv.swap 4 5)) ∧
      ¬AppliedModelingLib.Alignment.Axioms.LinearFeasibleRanking
          (fun (candidate : AppliedModelingLib.SocialChoice.Ranking.Candidate 5) (coordinate : Fin 3) =>
            if candidate = ⟨↑coordinate, ⋯⟩ then 1
            else
              if candidate = ⟨↑coordinate + 3, ⋯⟩ then -1
              else if candidate = ⟨6, GeEtAl2024AlignmentAxioms.appendixBStar._proof_2⟩ then 1 / 4 else 0)
          ((((((Equiv.refl (AppliedModelingLib.SocialChoice.Ranking.Candidate 5)).trans (Equiv.swap 0 6)).trans
                        (Equiv.swap 1 0)).trans
                    (Equiv.swap 2 1)).trans
                (Equiv.swap 3 2)).trans
            (Equiv.swap 4 5))
\end{ReviewVerbatim}



\paragraph{Lean proof endpoint (built separately)}\begin{ReviewVerbatim}
GeEtAl2024AlignmentAxioms.appendixBExplicitInfeasiblePMC
\end{ReviewVerbatim}

\paragraph{Recorded source-to-Spec screening}
\textbf{Verdict:} \texttt{matches}
\\
{\raggedright
\textbf{Reason:} Compared the Appendix-\allowbreak{}B seven-\allowbreak{}candidate construction, all three submitted rankings, the asserted unique pairwise-\allowbreak{}majority ranking, and its infeasibility.\allowbreak{} Candidate 5 has seven elements;\allowbreak{} the expanded features are the three positive and three negative standard-\allowbreak{}basis vectors plus (1/\allowbreak{}4,1/\allowbreak{}4,1/\allowbreak{}4), the expanded permutations reproduce the displayed voter and PMC orders, and the conclusion includes submitted-\allowbreak{}profile feasibility, majority realization, uniqueness, and failure of nondegenerate linear feasibility.\allowbreak{}
\par}
\reviewmatch{Matches source input}
\noindent\textbf{Reviewer annotation}\par
\reviewerbox
\clearpage
\hypertarget{source-claim-19bc4515dc72f018}{}
\section*{Review row 11}
\textbf{Source locator:} {\footnotesize\raggedright cited publication, lines 1398-\allowbreak{}1399\par}
\paragraph{Verbatim source input}
\begin{ReviewVerbatim}
Theorem C.2. Both Copeland (in traditional social choice) and LCPO (in linear social choice) fail
separability.

[Next verbatim source excerpt]

Proof. Consider the following two profiles on 7 candidates with five and three voters each:

                    Rank     v1    v2    v3        v4    v5           Rank     v6   v7   v8
                      1       a     b    b         c     d                 1   a    b    c
                      2       g     a    a         e     c                 2   d    a    a
                      3       d     c    e         f     f                 3   b    e    e
                      4       e     e    d         g     g                 4   f    d    b
                      5       f     d    f         b     b                 5   c    c    f
                      6       b     g    g         a     e                 6   g    g    g
                      7       c     f    c         d     a                 7   e    f    d

In the first profile, the Copeland scores are 5, 4, 3, 3, 3, 2, 1 for candidates a, b, c, d, e, f, g, respec-
tively. Similarly, in the second profile, they are 6, 5, 3, 3, 3, 1, 0. So under any consistent tie-breaking
rule, both of these profiles would output the ranking a ≻ b ≻ c ≻ d ≻ e ≻ f ≻ g.
However, if we combine these two profiles, then the score of a is 5 while the score of b is 6, and
thus, b will be ranked above a, violating separability.
To see that this also holds for LCPO, note that when every ranking is feasible, LCPO coincides with
Copeland. We can simply have 7 candidates in R7 all located at unit vectors, and voters with inputs
from the above profiles.
\end{ReviewVerbatim}



\paragraph{Expanded PaperInterface specification}
\begin{ReviewVerbatim}
(∀
    (rule :
      (voterCount : ℕ) →
        AppliedModelingLib.Alignment.Axioms.LinearRankAggregationRule (Fin voterCount) 5
          fun (x : AppliedModelingLib.SocialChoice.Ranking.Ranking 5) => True),
    ((∀ (voterCount : ℕ) (profile : AppliedModelingLib.Alignment.Axioms.RankingProfile (Fin voterCount) 5)
          (first second : AppliedModelingLib.SocialChoice.Ranking.Candidate 5),
          {opponent : AppliedModelingLib.SocialChoice.Ranking.Candidate 5 |
                  Fintype.card (Fin voterCount) <
                    2 *
                      {voter : Fin voterCount |
                          AppliedModelingLib.Alignment.Axioms.StrictlyPrefers (profile voter) second
                            opponent}.card}.card <
              {opponent : AppliedModelingLib.SocialChoice.Ranking.Candidate 5 |
                  Fintype.card (Fin voterCount) <
                    2 *
                      {voter : Fin voterCount |
                          AppliedModelingLib.Alignment.Axioms.StrictlyPrefers (profile voter) first
                            opponent}.card}.card →
            AppliedModelingLib.Alignment.Axioms.StrictlyPrefers ((rule voterCount).1 profile) first second) ∧
        ∀ (leftCount rightCount : ℕ) (left : AppliedModelingLib.Alignment.Axioms.RankingProfile (Fin leftCount) 5)
          (right : AppliedModelingLib.Alignment.Axioms.RankingProfile (Fin rightCount) 5),
          (∀ (first second : AppliedModelingLib.SocialChoice.Ranking.Candidate 5),
              {opponent : AppliedModelingLib.SocialChoice.Ranking.Candidate 5 |
                      Fintype.card (Fin leftCount) <
                        2 *
                          {voter : Fin leftCount |
                              AppliedModelingLib.Alignment.Axioms.StrictlyPrefers (left voter) first
                                opponent}.card}.card <
                  {opponent : AppliedModelingLib.SocialChoice.Ranking.Candidate 5 |
                      Fintype.card (Fin leftCount) <
                        2 *
                          {voter : Fin leftCount |
                              AppliedModelingLib.Alignment.Axioms.StrictlyPrefers (left voter) second
                                opponent}.card}.card ↔
                {opponent : AppliedModelingLib.SocialChoice.Ranking.Candidate 5 |
                      Fintype.card (Fin rightCount) <
                        2 *
                          {voter : Fin rightCount |
                              AppliedModelingLib.Alignment.Axioms.StrictlyPrefers (right voter) first
                                opponent}.card}.card <
                  {opponent : AppliedModelingLib.SocialChoice.Ranking.Candidate 5 |
                      Fintype.card (Fin rightCount) <
                        2 *
                          {voter : Fin rightCount |
                              AppliedModelingLib.Alignment.Axioms.StrictlyPrefers (right voter) second
                                opponent}.card}.card) →
            (rule leftCount).1 left = (rule rightCount).1 right) →
      ¬AppliedModelingLib.Alignment.Axioms.RankingSeparability
          (fun (x : AppliedModelingLib.SocialChoice.Ranking.Ranking 5) => True) rule) ∧
  ∀
    (copelandRule lcpoRule :
      (voterCount : ℕ) →
        AppliedModelingLib.Alignment.Axioms.LinearRankAggregationRule (Fin voterCount) 5
          fun (x : AppliedModelingLib.SocialChoice.Ranking.Ranking 5) => True),
    (((∀ (voterCount : ℕ) (profile : AppliedModelingLib.Alignment.Axioms.RankingProfile (Fin voterCount) 5)
            (first second : AppliedModelingLib.SocialChoice.Ranking.Candidate 5),
            {opponent : AppliedModelingLib.SocialChoice.Ranking.Candidate 5 |
                    Fintype.card (Fin voterCount) <
                      2 *
                        {voter : Fin voterCount |
                            AppliedModelingLib.Alignment.Axioms.StrictlyPrefers (profile voter) second
                              opponent}.card}.card <
                {opponent : AppliedModelingLib.SocialChoice.Ranking.Candidate 5 |
                    Fintype.card (Fin voterCount) <
                      2 *
                        {voter : Fin voterCount |
                            AppliedModelingLib.Alignment.Axioms.StrictlyPrefers (profile voter) first
                              opponent}.card}.card →
              AppliedModelingLib.Alignment.Axioms.StrictlyPrefers ((copelandRule voterCount).1 profile) first second) ∧
          ∀ (leftCount rightCount : ℕ) (left : AppliedModelingLib.Alignment.Axioms.RankingProfile (Fin leftCount) 5)
            (right : AppliedModelingLib.Alignment.Axioms.RankingProfile (Fin rightCount) 5),
            (∀ (first second : AppliedModelingLib.SocialChoice.Ranking.Candidate 5),
                {opponent : AppliedModelingLib.SocialChoice.Ranking.Candidate 5 |
                        Fintype.card (Fin leftCount) <
                          2 *
                            {voter : Fin leftCount |
                                AppliedModelingLib.Alignment.Axioms.StrictlyPrefers (left voter) first
                                  opponent}.card}.card <
                    {opponent : AppliedModelingLib.SocialChoice.Ranking.Candidate 5 |
                        Fintype.card (Fin leftCount) <
                          2 *
                            {voter : Fin leftCount |
                                AppliedModelingLib.Alignment.Axioms.StrictlyPrefers (left voter) second
                                  opponent}.card}.card ↔
                  {opponent : AppliedModelingLib.SocialChoice.Ranking.Candidate 5 |
                        Fintype.card (Fin rightCount) <
                          2 *
                            {voter : Fin rightCount |
                                AppliedModelingLib.Alignment.Axioms.StrictlyPrefers (right voter) first
                                  opponent}.card}.card <
                    {opponent : AppliedModelingLib.SocialChoice.Ranking.Candidate 5 |
                        Fintype.card (Fin rightCount) <
                          2 *
                            {voter : Fin rightCount |
                                AppliedModelingLib.Alignment.Axioms.StrictlyPrefers (right voter) second
                                  opponent}.card}.card) →
              (copelandRule leftCount).1 left = (copelandRule rightCount).1 right) ∧
        (∀ (voterCount : ℕ) (profile : AppliedModelingLib.Alignment.Axioms.RankingProfile (Fin voterCount) 5),
            AppliedModelingLib.Alignment.Axioms.FeasibleProfile
                (fun (x : AppliedModelingLib.SocialChoice.Ranking.Ranking 5) => True) profile →
              ((fun (x : AppliedModelingLib.SocialChoice.Ranking.Ranking 5) => True) ((lcpoRule voterCount).1 profile) ∧
                  ∀ (first second : AppliedModelingLib.SocialChoice.Ranking.Candidate 5),
                    (∀ (voter : Fin voterCount),
                        AppliedModelingLib.Alignment.Axioms.StrictlyPrefers (profile voter) first second) →
                      AppliedModelingLib.Alignment.Axioms.StrictlyPrefers ((lcpoRule voterCount).1 profile) first
                        second) ∧
                ∀ (position candidate : AppliedModelingLib.SocialChoice.Ranking.Candidate 5),
                  (∃ (contender : AppliedModelingLib.SocialChoice.Ranking.Ranking 5),
                      ((fun (x : AppliedModelingLib.SocialChoice.Ranking.Ranking 5) => True) contender ∧
                          ∀ (first second : AppliedModelingLib.SocialChoice.Ranking.Candidate 5),
                            (∀ (voter : Fin voterCount),
                                AppliedModelingLib.Alignment.Axioms.StrictlyPrefers (profile voter) first second) →
                              AppliedModelingLib.Alignment.Axioms.StrictlyPrefers contender first second) ∧
                        (∀ earlier < position, contender earlier = ((lcpoRule voterCount).1 profile) earlier) ∧
                          contender position = candidate) →
                    {opponent : AppliedModelingLib.SocialChoice.Ranking.Candidate 5 |
                          Fintype.card (Fin voterCount) <
                            2 *
                              {voter : Fin voterCount |
                                  AppliedModelingLib.Alignment.Axioms.StrictlyPrefers (profile voter) candidate
                                    opponent}.card}.card ≤
                      {opponent : AppliedModelingLib.SocialChoice.Ranking.Candidate 5 |
                          Fintype.card (Fin voterCount) <
                            2 *
                              {voter : Fin voterCount |
                                  AppliedModelingLib.Alignment.Axioms.StrictlyPrefers (profile voter)
                                    (((lcpoRule voterCount).1 profile) position) opponent}.card}.card) ∧
          ∀ (voterCount : ℕ) (profile : AppliedModelingLib.Alignment.Axioms.RankingProfile (Fin voterCount) 5),
            ((fun (x : AppliedModelingLib.SocialChoice.Ranking.Ranking 5) => True)
                  ((copelandRule voterCount).1 profile) ∧
                ∀ (first second : AppliedModelingLib.SocialChoice.Ranking.Candidate 5),
                  (∀ (voter : Fin voterCount),
                      AppliedModelingLib.Alignment.Axioms.StrictlyPrefers (profile voter) first second) →
                    AppliedModelingLib.Alignment.Axioms.StrictlyPrefers ((copelandRule voterCount).1 profile) first
                      second) →
              (lcpoRule voterCount).1 profile = (copelandRule voterCount).1 profile) →
      ¬AppliedModelingLib.Alignment.Axioms.RankingSeparability
          (fun (x : AppliedModelingLib.SocialChoice.Ranking.Ranking 5) => True) lcpoRule
\end{ReviewVerbatim}



\paragraph{Lean proof endpoint (built separately)}\begin{ReviewVerbatim}
GeEtAl2024AlignmentAxioms.theoremC_2_copelandAndLCPOFailSeparability
\end{ReviewVerbatim}

\paragraph{Recorded source-to-Spec screening}
\textbf{Verdict:} \texttt{matches}
\\
{\raggedright
\textbf{Reason:} The two target branches encode the source proof's consistent tie convention:\allowbreak{} Copeland preserves strict score order and uses one output for equal weak score orders, while LCPO uses that convention whenever the Copeland outcome is Pareto-\allowbreak{}feasible.\allowbreak{} On the literal seven-\allowbreak{}candidate all-\allowbreak{}feasible construction, each branch concludes failure of separability as in Theorem C.\allowbreak{}2.\allowbreak{}
\par}
\reviewmatch{Matches source input}
\noindent\textbf{Reviewer annotation}\par
\reviewerbox
\clearpage
\hypertarget{source-claim-db88b9868ac40a8b}{}
\section*{Review row 12}
\textbf{Source locator:} {\footnotesize\raggedright cited publication, lines 1450\par}
\paragraph{Verbatim source input}
\begin{ReviewVerbatim}
Theorem C.3. Linear Kemeny satisfies PMC and separability.

[Next verbatim source excerpt]


[form-feed in source extraction]
C.2   Linear Kemeny Rule

Next, we consider a different rule from social choice theory, the Kemeny rule, which can be trans-
formed to the linear setting while maintaining the separability can be achieved along with PMC.
Given an input profile π, the Kemeny rule returns a ranking σ ∗ that minimizes the total number of
pairwise disagreements with voters rankings, i.e.
                                            X X
                              σ ∗ ∈ arg min                1(b ≻σ a)
                                        σ∈S m
                                                 i∈V (a,b):a≻σi b
        m
where S contains all possible permutations of the m candidates. This expression can be equiva-
lently written as                        X
                          σ ∗ ∈ arg min        na≻b (π) · 1(b ≻σ a).
                                      σ∈S m
                                                a̸=b∈C
Here, we define the linear Kemeny rule, which outputs a parameter vector θ∗ that induces a ranking
that minimizes the total number of pairwise disagreements with voters’ rankings, i.e.,
                                        X
                          θ∗ ∈ arg min       na≻b (π) · 1(rθ (b) > rθ (a)).
                                 θ∈Rd    a̸=b∈C
Note that this rule conforms to the standard loss formulation, where the loss function is binary: it is
0 if two rankings agree with respect to the relative ranking of a pair of candidates and 1 otherwise.
Since binary loss is not convex, it does not fit in the impossibility result of Theorem 3.1.
Note that Kemeny is generally NP-hard to compute; however, even for linear Kemeny, there is
at least an exponential time aglorithm by brute-force computing the score of every ranking, and
determining whether or not it is feasible.
Theorem C.3. Linear Kemeny satisfies PMC and separability.

Proof. PMC holds because the linear Kemeny score minimizes disagreements even among non-
transitive rankings, making the PMC ranking the optimal choice whenever it is feasible. Separability
is evident as the Kemeny score of a ranking over two datasets is simply the sum of the scores in each
dataset. If the same ranking minimizes the score in both datasets independently, it will also minimize
the score in their combination.
\end{ReviewVerbatim}



\paragraph{Expanded PaperInterface specification}
\begin{ReviewVerbatim}
∀ {n : ℕ} (feasible : AppliedModelingLib.SocialChoice.Ranking.Ranking n → Prop)
  (fallback : AppliedModelingLib.SocialChoice.Ranking.Ranking n) (hfallback : feasible fallback) (voterCount : ℕ),
  (∀ (profile : AppliedModelingLib.Alignment.Axioms.RankingProfile (Fin voterCount) n)
      (majorityRanking : AppliedModelingLib.SocialChoice.Ranking.Ranking n),
      AppliedModelingLib.Alignment.Axioms.FeasibleProfile feasible profile →
        feasible majorityRanking →
          (∀ (first second : AppliedModelingLib.SocialChoice.Ranking.Candidate n),
              AppliedModelingLib.Alignment.Axioms.StrictlyPrefers majorityRanking first second ↔
                Fintype.card (Fin voterCount) <
                  2 *
                    {voter : Fin voterCount |
                        AppliedModelingLib.Alignment.Axioms.StrictlyPrefers (profile voter) first second}.card) →
            (AppliedModelingLib.Alignment.Axioms.canonicalKemenyRule feasible fallback hfallback voterCount).1 profile =
              majorityRanking) ∧
    AppliedModelingLib.Alignment.Axioms.RankingSeparability feasible
      (AppliedModelingLib.Alignment.Axioms.canonicalKemenyRule feasible fallback hfallback)
\end{ReviewVerbatim}



\paragraph{Lean proof endpoint (built separately)}\begin{ReviewVerbatim}
GeEtAl2024AlignmentAxioms.theoremC_3_linearKemeny
\end{ReviewVerbatim}

\paragraph{Recorded source-to-Spec screening}
\textbf{Verdict:} \texttt{matches}
\\
{\raggedright
\textbf{Reason:} Compared the linear-\allowbreak{}Kemeny objective, feasible-\allowbreak{}domain restriction, PMC claim, and separability claim.\allowbreak{} The canonical rule minimizes the exact voter-\allowbreak{}by-\allowbreak{}ordered-\allowbreak{}pair reversal count among feasible rankings and applies one profile-\allowbreak{}independent tie key;\allowbreak{} the conclusion returns every feasible PMC ranking and says that equal selected minimizers for two profiles remain the selected minimizer after literal concatenation, matching the additive Kemeny argument in Theorem C.\allowbreak{}3.\allowbreak{}
\par}
\reviewmatch{Matches source input}
\noindent\textbf{Reviewer annotation}\par
\reviewerbox
\clearpage
\hypertarget{source-claim-aba4949a527a8006}{}
\section*{Review row 13}
\textbf{Source locator:} {\footnotesize\raggedright cited publication, lines 1457\par}
\paragraph{Verbatim source input}
\begin{ReviewVerbatim}
Theorem C.4. Linear Kemeny does not satisfy PO or majority consistency.

[Next verbatim source excerpt]

Proof. Consider the scenario with 20 candidates whose feature vectors are represented in the table
below:

                             Candidate        Feature Vector
                                  1           (2000000, 0, 0, 0, 0, 0, 0)
                                  2           (0, 2000000, 0, 0, 0, 0, 0)
                                  3           (0, 200000, 0, 0, 0, 0, 0)
                                  4           (0, 100000, 100000, 0, 0, 0, 0)
                                  5           (0, 0, 200000, 0, 0, 0, 0)
                                  6           (0, 0, 20000, 0, 0, 0, 0)
                                  7           (0, 0, 10000, 10000, 0, 0, 0)
                                  8           (0, 0, 0, 20000, 0, 0, 0)
                                  9           (0, 0, 0, 2000, 0, 0, 0)
                                 10           (0, 0, 0, 1000, 1000, 0, 0)
                                 11           (0, 0, 0, 0, 2000, 0, 0)
                                 12           (0, 0, 0, 0, 200, 0, 0)
                                 13           (0, 0, 0, 0, 100, 100, 0)
                                 14           (0, 0, 0, 0, 0, 200, 0)
                                 15           (0, 0, 0, 0, 0, 20, 0)
                                 16           (0, 0, 0, 0, 0, 10, 10)
                                 17           (0, 0, 0, 0, 0, 0, 20)
                                 18           (0, 0, 0, 0, 0, 0, 2)
                                 19           (1, 0, 0, 0, 0, 0, 1)
                                 20           (2, 0, 0, 0, 0, 0, 0)


                                                     25

[form-feed in source extraction]
Each vector is constructed such that candidates are prioritized based on the magnitude of their
first non-zero entry, leading to a natural ordering within grouped subsets: {1, 2} ≻ {3, 4, 5} ≻
{6, 7, 8} ≻ {9, 10, 11} ≻ {12, 13, 14} ≻ {15, 16, 17} ≻ {18, 19, 20}. We will have six voters,
with rankings induced by the parameter vectors described below.

                                      Voter    Parameter Vector
                                        v1     (2, 1, 7, 6, 5, 4, 3)
                                        v2     (3, 2, 1, 7, 6, 5, 4)
                                        v3     (4, 3, 2, 1, 7, 6, 5)
                                        v4     (5, 4, 3, 2, 1, 7, 6)
                                        v5     (6, 5, 4, 3, 2, 1, 7)
                                        v6     (7, 6, 5, 4, 3, 2, 1)

Each voter ranks candidates within the group in increasing order except for a single reversed group.
For instance, v1 ranks candidates as 1 ≻ 2 ≻ 5 ≻ 4 ≻ 3 and so forth. Under this setup, no param-
eter vector θ can create a ranking that satisfies all voters’ preferences due to the cyclic nature and
individual group preferences.
We can check that linear Kemeny rule outputs a ranking in which candidate 2 is ranked above 1.
From this analysis, we also conclude that the rules does not satisfy majority consistency either.
\end{ReviewVerbatim}



\paragraph{Expanded PaperInterface specification}
\begin{ReviewVerbatim}
∀
  (rule :
    AppliedModelingLib.Alignment.Axioms.LinearRankAggregationRule (Fin 6) 18
      (AppliedModelingLib.Alignment.Axioms.LinearFeasibleRanking
        ![![2000000, 0, 0, 0, 0, 0, 0], ![0, 2000000, 0, 0, 0, 0, 0], ![0, 200000, 0, 0, 0, 0, 0],
          ![0, 100000, 100000, 0, 0, 0, 0], ![0, 0, 200000, 0, 0, 0, 0], ![0, 0, 20000, 0, 0, 0, 0],
          ![0, 0, 10000, 10000, 0, 0, 0], ![0, 0, 0, 20000, 0, 0, 0], ![0, 0, 0, 2000, 0, 0, 0],
          ![0, 0, 0, 1000, 1000, 0, 0], ![0, 0, 0, 0, 2000, 0, 0], ![0, 0, 0, 0, 200, 0, 0], ![0, 0, 0, 0, 100, 100, 0],
          ![0, 0, 0, 0, 0, 200, 0], ![0, 0, 0, 0, 0, 20, 0], ![0, 0, 0, 0, 0, 10, 10], ![0, 0, 0, 0, 0, 0, 20],
          ![0, 0, 0, 0, 0, 0, 2], ![1, 0, 0, 0, 0, 0, 1], ![2, 0, 0, 0, 0, 0, 0]])),
  AppliedModelingLib.Alignment.Axioms.IsKemenySelector
      (AppliedModelingLib.Alignment.Axioms.LinearFeasibleRanking
        ![![2000000, 0, 0, 0, 0, 0, 0], ![0, 2000000, 0, 0, 0, 0, 0], ![0, 200000, 0, 0, 0, 0, 0],
          ![0, 100000, 100000, 0, 0, 0, 0], ![0, 0, 200000, 0, 0, 0, 0], ![0, 0, 20000, 0, 0, 0, 0],
          ![0, 0, 10000, 10000, 0, 0, 0], ![0, 0, 0, 20000, 0, 0, 0], ![0, 0, 0, 2000, 0, 0, 0],
          ![0, 0, 0, 1000, 1000, 0, 0], ![0, 0, 0, 0, 2000, 0, 0], ![0, 0, 0, 0, 200, 0, 0], ![0, 0, 0, 0, 100, 100, 0],
          ![0, 0, 0, 0, 0, 200, 0], ![0, 0, 0, 0, 0, 20, 0], ![0, 0, 0, 0, 0, 10, 10], ![0, 0, 0, 0, 0, 0, 20],
          ![0, 0, 0, 0, 0, 0, 2], ![1, 0, 0, 0, 0, 0, 1], ![2, 0, 0, 0, 0, 0, 0]])
      rule →
    ¬AppliedModelingLib.Alignment.Axioms.ParetoOptimal
          (AppliedModelingLib.Alignment.Axioms.LinearFeasibleRanking
            ![![2000000, 0, 0, 0, 0, 0, 0], ![0, 2000000, 0, 0, 0, 0, 0], ![0, 200000, 0, 0, 0, 0, 0],
              ![0, 100000, 100000, 0, 0, 0, 0], ![0, 0, 200000, 0, 0, 0, 0], ![0, 0, 20000, 0, 0, 0, 0],
              ![0, 0, 10000, 10000, 0, 0, 0], ![0, 0, 0, 20000, 0, 0, 0], ![0, 0, 0, 2000, 0, 0, 0],
              ![0, 0, 0, 1000, 1000, 0, 0], ![0, 0, 0, 0, 2000, 0, 0], ![0, 0, 0, 0, 200, 0, 0],
              ![0, 0, 0, 0, 100, 100, 0], ![0, 0, 0, 0, 0, 200, 0], ![0, 0, 0, 0, 0, 20, 0], ![0, 0, 0, 0, 0, 10, 10],
              ![0, 0, 0, 0, 0, 0, 20], ![0, 0, 0, 0, 0, 0, 2], ![1, 0, 0, 0, 0, 0, 1], ![2, 0, 0, 0, 0, 0, 0]])
          rule ∧
      ¬AppliedModelingLib.Alignment.Axioms.MajorityConsistent
          (AppliedModelingLib.Alignment.Axioms.LinearFeasibleRanking
            ![![2000000, 0, 0, 0, 0, 0, 0], ![0, 2000000, 0, 0, 0, 0, 0], ![0, 200000, 0, 0, 0, 0, 0],
              ![0, 100000, 100000, 0, 0, 0, 0], ![0, 0, 200000, 0, 0, 0, 0], ![0, 0, 20000, 0, 0, 0, 0],
              ![0, 0, 10000, 10000, 0, 0, 0], ![0, 0, 0, 20000, 0, 0, 0], ![0, 0, 0, 2000, 0, 0, 0],
              ![0, 0, 0, 1000, 1000, 0, 0], ![0, 0, 0, 0, 2000, 0, 0], ![0, 0, 0, 0, 200, 0, 0],
              ![0, 0, 0, 0, 100, 100, 0], ![0, 0, 0, 0, 0, 200, 0], ![0, 0, 0, 0, 0, 20, 0], ![0, 0, 0, 0, 0, 10, 10],
              ![0, 0, 0, 0, 0, 0, 20], ![0, 0, 0, 0, 0, 0, 2], ![1, 0, 0, 0, 0, 0, 1], ![2, 0, 0, 0, 0, 0, 0]])
          rule
\end{ReviewVerbatim}



\paragraph{Lean proof endpoint (built separately)}\begin{ReviewVerbatim}
GeEtAl2024AlignmentAxioms.theoremC_4_linearKemenyFailsParetoAndMajorityConsistency
\end{ReviewVerbatim}

\paragraph{Recorded source-to-Spec screening}
\textbf{Verdict:} \texttt{matches}
\\
{\raggedright
\textbf{Reason:} Compared the literal twenty-\allowbreak{}candidate, seven-\allowbreak{}coordinate feature table and six-\allowbreak{}voter counterexample.\allowbreak{} Candidate 18 expands to twenty candidates and the displayed feature matrix agrees entry by entry with the source.\allowbreak{} For every rule selecting a feasible minimizer of the exact Kemeny disagreement objective on every profile, the target concludes both failure of unanimity-\allowbreak{}based Pareto optimality and failure of strict-\allowbreak{}majority top-\allowbreak{}choice consistency.\allowbreak{}
\par}
\reviewmatch{Matches source input}
\noindent\textbf{Reviewer annotation}\par
\reviewerbox
\clearpage
\hypertarget{source-claim-e6c3a6a88028de60}{}
\section*{Review row 14}
\textbf{Source locator:} {\footnotesize\raggedright cited publication, lines 1509\par}
\paragraph{Verbatim source input}
\begin{ReviewVerbatim}
Theorem C.5. Linear Kemeny subject to PO violates separability and majority consistency.

[Next verbatim source excerpt]

A possibly easy fix to the problem that linear Kemeny does not satisfy PO would be to enforce the
PO criterion, as we did for the LCPO, i.e., to restrict to parameter vectors that respect PO. However,
we show that if we do that, then linear Kemeny subject to PO does not satisfy separability anymore.
Theorem C.5. Linear Kemeny subject to PO violates separability and majority consistency.

Proof. First, consider a set of candidates and a profile similar to the one that is given in the proof
of Theorem C.4. When we restrict to reward functions that 1 above 2, then we can check that Linear
Kemeny subject to PO outputs one of the rankings that are in the input profile. Without loss of
generality, assume that it outputs the ranking of v1 .
Second, consider the same set of candidates and three voters, with rankings induced by the following
parameter vectors:

                                      Voter    Parameter Vector
                                        v1′    (2, 1, 7, 6, 5, 4, 3)
                                        v2′    (2, 1, 7, 6, 5, 4, 3)
                                        v3′    (1, 7, 6, 5, 4, 3, 2)

In this case, Linear Kemeny subject to PO outputs the ranking of v1′ which is the same with this of
voter v1 .
When the two profiles are combined, then we do not anymore restrict on rankings in which 1 is above
2, since 1 does not Pareto dominates 2 anymore. Then, we can check that linear Kemeny subject to
PO outputs a different ranking than before in which 2 is ranked above 1. From this example, we see
that linear Kemenery subject to PO still violates majority consistency.

[Next verbatim source excerpt]

Proof. Consider the scenario with 20 candidates whose feature vectors are represented in the table
below:

                             Candidate        Feature Vector
                                  1           (2000000, 0, 0, 0, 0, 0, 0)
                                  2           (0, 2000000, 0, 0, 0, 0, 0)
                                  3           (0, 200000, 0, 0, 0, 0, 0)
                                  4           (0, 100000, 100000, 0, 0, 0, 0)
                                  5           (0, 0, 200000, 0, 0, 0, 0)
                                  6           (0, 0, 20000, 0, 0, 0, 0)
                                  7           (0, 0, 10000, 10000, 0, 0, 0)
                                  8           (0, 0, 0, 20000, 0, 0, 0)
                                  9           (0, 0, 0, 2000, 0, 0, 0)
                                 10           (0, 0, 0, 1000, 1000, 0, 0)
                                 11           (0, 0, 0, 0, 2000, 0, 0)
                                 12           (0, 0, 0, 0, 200, 0, 0)
                                 13           (0, 0, 0, 0, 100, 100, 0)
                                 14           (0, 0, 0, 0, 0, 200, 0)
                                 15           (0, 0, 0, 0, 0, 20, 0)
                                 16           (0, 0, 0, 0, 0, 10, 10)
                                 17           (0, 0, 0, 0, 0, 0, 20)
                                 18           (0, 0, 0, 0, 0, 0, 2)
                                 19           (1, 0, 0, 0, 0, 0, 1)
                                 20           (2, 0, 0, 0, 0, 0, 0)


                                                     25

[form-feed in source extraction]
Each vector is constructed such that candidates are prioritized based on the magnitude of their
first non-zero entry, leading to a natural ordering within grouped subsets: {1, 2} ≻ {3, 4, 5} ≻
{6, 7, 8} ≻ {9, 10, 11} ≻ {12, 13, 14} ≻ {15, 16, 17} ≻ {18, 19, 20}. We will have six voters,
with rankings induced by the parameter vectors described below.

                                      Voter    Parameter Vector
                                        v1     (2, 1, 7, 6, 5, 4, 3)
                                        v2     (3, 2, 1, 7, 6, 5, 4)
                                        v3     (4, 3, 2, 1, 7, 6, 5)
                                        v4     (5, 4, 3, 2, 1, 7, 6)
                                        v5     (6, 5, 4, 3, 2, 1, 7)
                                        v6     (7, 6, 5, 4, 3, 2, 1)

Each voter ranks candidates within the group in increasing order except for a single reversed group.
For instance, v1 ranks candidates as 1 ≻ 2 ≻ 5 ≻ 4 ≻ 3 and so forth. Under this setup, no param-
eter vector θ can create a ranking that satisfies all voters’ preferences due to the cyclic nature and
individual group preferences.
We can check that linear Kemeny rule outputs a ranking in which candidate 2 is ranked above 1.
From this analysis, we also conclude that the rules does not satisfy majority consistency either.
\end{ReviewVerbatim}

% report-context-presentation-sha256: 422827fc2f870981d9a9e288ba1e0082da0be346a37f5720f173bff1111f3733
\paragraph{Settled source reading}
The checked C.\allowbreak{}5 counterexample retains the source proof’s local choice of v1 on its first profile as an explicit premise.\allowbreak{} Recording that premise does not itself prove a reduction from every Pareto-\allowbreak{}Kemeny selector.\allowbreak{}

\paragraph{Expanded PaperInterface specification}
\begin{ReviewVerbatim}
∀
  (rule :
    (voterCount : ℕ) →
      AppliedModelingLib.Alignment.Axioms.LinearRankAggregationRule (Fin voterCount) 18
        (AppliedModelingLib.Alignment.Axioms.LinearFeasibleRanking
          ![![2000000, 0, 0, 0, 0, 0, 0], ![0, 2000000, 0, 0, 0, 0, 0], ![0, 200000, 0, 0, 0, 0, 0],
            ![0, 100000, 100000, 0, 0, 0, 0], ![0, 0, 200000, 0, 0, 0, 0], ![0, 0, 20000, 0, 0, 0, 0],
            ![0, 0, 10000, 10000, 0, 0, 0], ![0, 0, 0, 20000, 0, 0, 0], ![0, 0, 0, 2000, 0, 0, 0],
            ![0, 0, 0, 1000, 1000, 0, 0], ![0, 0, 0, 0, 2000, 0, 0], ![0, 0, 0, 0, 200, 0, 0],
            ![0, 0, 0, 0, 100, 100, 0], ![0, 0, 0, 0, 0, 200, 0], ![0, 0, 0, 0, 0, 20, 0], ![0, 0, 0, 0, 0, 10, 10],
            ![0, 0, 0, 0, 0, 0, 20], ![0, 0, 0, 0, 0, 0, 2], ![1, 0, 0, 0, 0, 0, 1], ![2, 0, 0, 0, 0, 0, 0]])),
  GeEtAl2024AlignmentAxioms.IsParetoKemenySelector rule →
    (rule 6).1
          ![Equiv.swap 2 4, Equiv.swap 5 7, Equiv.swap 8 10, Equiv.swap 11 13, Equiv.swap 14 16, Equiv.swap 17 19] =
        Equiv.swap 2 4 →
      ¬AppliedModelingLib.Alignment.Axioms.RankingSeparability
            (AppliedModelingLib.Alignment.Axioms.LinearFeasibleRanking
              ![![2000000, 0, 0, 0, 0, 0, 0], ![0, 2000000, 0, 0, 0, 0, 0], ![0, 200000, 0, 0, 0, 0, 0],
                ![0, 100000, 100000, 0, 0, 0, 0], ![0, 0, 200000, 0, 0, 0, 0], ![0, 0, 20000, 0, 0, 0, 0],
                ![0, 0, 10000, 10000, 0, 0, 0], ![0, 0, 0, 20000, 0, 0, 0], ![0, 0, 0, 2000, 0, 0, 0],
                ![0, 0, 0, 1000, 1000, 0, 0], ![0, 0, 0, 0, 2000, 0, 0], ![0, 0, 0, 0, 200, 0, 0],
                ![0, 0, 0, 0, 100, 100, 0], ![0, 0, 0, 0, 0, 200, 0], ![0, 0, 0, 0, 0, 20, 0], ![0, 0, 0, 0, 0, 10, 10],
                ![0, 0, 0, 0, 0, 0, 20], ![0, 0, 0, 0, 0, 0, 2], ![1, 0, 0, 0, 0, 0, 1], ![2, 0, 0, 0, 0, 0, 0]])
            rule ∧
        ¬AppliedModelingLib.Alignment.Axioms.MajorityConsistent
            (AppliedModelingLib.Alignment.Axioms.LinearFeasibleRanking
              ![![2000000, 0, 0, 0, 0, 0, 0], ![0, 2000000, 0, 0, 0, 0, 0], ![0, 200000, 0, 0, 0, 0, 0],
                ![0, 100000, 100000, 0, 0, 0, 0], ![0, 0, 200000, 0, 0, 0, 0], ![0, 0, 20000, 0, 0, 0, 0],
                ![0, 0, 10000, 10000, 0, 0, 0], ![0, 0, 0, 20000, 0, 0, 0], ![0, 0, 0, 2000, 0, 0, 0],
                ![0, 0, 0, 1000, 1000, 0, 0], ![0, 0, 0, 0, 2000, 0, 0], ![0, 0, 0, 0, 200, 0, 0],
                ![0, 0, 0, 0, 100, 100, 0], ![0, 0, 0, 0, 0, 200, 0], ![0, 0, 0, 0, 0, 20, 0], ![0, 0, 0, 0, 0, 10, 10],
                ![0, 0, 0, 0, 0, 0, 20], ![0, 0, 0, 0, 0, 0, 2], ![1, 0, 0, 0, 0, 0, 1], ![2, 0, 0, 0, 0, 0, 0]])
            (rule 9)
\end{ReviewVerbatim}



\paragraph{Lean proof endpoint (built separately)}\begin{ReviewVerbatim}
GeEtAl2024AlignmentAxioms.theoremC_5_paretoKemenyFailsSeparabilityAndMajorityConsistency
\end{ReviewVerbatim}

\paragraph{Recorded source-to-Spec screening}
\textbf{Verdict:} \texttt{matches}
\\
{\raggedright
\textbf{Reason:} For the literal C.\allowbreak{}5 feature instance, the target requires Pareto-\allowbreak{}constrained Kemeny minimization and only the local premise that the first six-\allowbreak{}voter profile outputs v1;\allowbreak{} it then concludes failure of separability and majority consistency.\allowbreak{} The approved source context identifies that premise as the proof's explicit WLOG selection and confirms that no global tie convention is assumed.\allowbreak{}
\par}
\reviewmatch{Matches source input}
\noindent\textbf{Reviewer annotation}\par
\reviewerbox
\clearpage
\hypertarget{source-claim-c7e3d91ee78bfd29}{}
\section*{Review row 15}
\textbf{Source locator:} {\footnotesize\raggedright cited publication, lines 1543-\allowbreak{}1544\par}
\paragraph{Verbatim source input}
\begin{ReviewVerbatim}
Theorem C.6. Leximax Plurality satisfies majority consistency, winner monotonicity and separa-
bility.

[Next verbatim source excerpt]

C.3   Leximax Plurality

Plurality is probably the most ubiquitous voting rule in the world. Its ranking variant ranks the can-
didates in decreasing order with respect to their plurality scores. The plurality score of a candidate
is equal to the number of her appearances in the first position. This rule is known to satisfy several
axioms but in linear social choice cannot be directly applied, as not all rankings are feasibly.
Similarly to leximax Copeland, we define Leximax Plurality as follows. It ranks first the candidate
with the highest plurality score that can be ranked first under some parameter vector. Subject to
this first position, it ranks second the candidate with the highest plurality score that can feasibly be
ranked second, and so on, until all the positions are filled.


                                                  26

[form-feed in source extraction]
Theorem C.6. Leximax Plurality satisfies majority consistency, winner monotonicity and separa-
bility.

Proof. Note that leximax Plurality always returns a ranking in which the candidate with the highest
plurality score is ranked first, since there exists at least one feasible ranking in which this candidate
is ranked first. From this observation, we immediately see that leximax Plurality satisfies majority
consistency.
Now, suppose that on input profile π, the rule outputs a ranking such that candidate a is ranked first.
This means that a has the highest plurality score. Now, consider a profile π ′ which is similar to π
with the only exception being a ranking in which a is ranked in a higher position. It is clear that a
continues to have the highest plurality score and therefore leximax Plurality will output a ranking in
which a is ranked first. Therefore, winner monotonicity is satisfied. .
It remains to prove that the rule satisfies separability. Suppose that under two different profiles π1
and π2 , the rule outputs σ, and under the aggregated profile π3 , it outputs σ ′ . We will show that
σ = σ ′ . One main observation is that if a candidate a has a higher plurality score than a candidate b
under both π1 and π2 , then a has a higher plurality score than b under π3 as well. We will show the
desired property by induction on the positions of the ranking σ. Start from the first position in which
say candidate a is ranked first. From above, we know that a has the highest plurality score under
both π1 and π2 , which remains true in π3 , and therefore a is ranked first in σ ′ . Now, assume that up
to position t − 1, σ and σ ′ are similar and denote with a′ the candidate that is ranked at position t of
σ. We denote by S1 , S2 and S3 the set of candidates that are not ranked among the first t positions
in σ and have higher plurality score than a′ under π1 , π2 , and π3 respectively. Since, a′ is ranked at
the t-th postion in σ, we get that, subject to the fixed first t − 1 position, no candidate in S1 ∪ S2 can
be ranked at the t-th position. The theorem follows by noticing that S3 ⊆ S1 ∪ S2 , which follows
from the main observation above.




                                                   27
\end{ReviewVerbatim}

% report-context-presentation-sha256: 4a862756d48b1da06bd1f7fc56799f77487941a409c0bc42f79b1e2a17e23482
\paragraph{Settled source reading}
The checked C.\allowbreak{}6 rule uses a fixed profile-\allowbreak{}independent tie key for the sequential highest-\allowbreak{}plurality prescription.\allowbreak{} Strict plurality comparisons are unchanged;\allowbreak{} the theorem has this specified rule-\allowbreak{}family scope.\allowbreak{}

\paragraph{Expanded PaperInterface specification}
\begin{ReviewVerbatim}
∀ {n : ℕ} (feasible : AppliedModelingLib.SocialChoice.Ranking.Ranking n → Prop)
  (rule : (voterCount : ℕ) → AppliedModelingLib.Alignment.Axioms.LinearRankAggregationRule (Fin voterCount) n feasible),
  (∀ (voterCount : ℕ), GeEtAl2024AlignmentAxioms.IsFixedTieLeximaxPluralitySelector feasible (rule voterCount)) →
    (∀ (voterCount : ℕ), AppliedModelingLib.Alignment.Axioms.MajorityConsistent feasible (rule voterCount)) ∧
      (∀ (voterCount : ℕ), AppliedModelingLib.Alignment.Axioms.WinnerMonotonic feasible (rule voterCount)) ∧
        AppliedModelingLib.Alignment.Axioms.RankingSeparability feasible rule
\end{ReviewVerbatim}



\paragraph{Lean proof endpoint (built separately)}\begin{ReviewVerbatim}
GeEtAl2024AlignmentAxioms.theoremC_6_fixedTieLeximaxPlurality
\end{ReviewVerbatim}

\paragraph{Recorded source-to-Spec screening}
\textbf{Verdict:} \texttt{matches}
\\
{\raggedright
\textbf{Reason:} The target states that every member of the rule family follows the source's sequential feasible highest-\allowbreak{}plurality choice and concludes majority consistency, winner monotonicity, and separability.\allowbreak{} The approved source context limits the fixed candidate order to equality ties, so it does not alter any strict plurality comparison in the source construction.\allowbreak{}
\par}
\reviewmatch{Matches source input}
\noindent\textbf{Reviewer annotation}\par
\reviewerbox

\end{document}
